\documentclass[11pt,a4paper]{article}
\usepackage{fontspec}
\setmainfont[BoldFont={Noto Sans CJK JP Bold}]{Noto Serif CJK JP}
\setsansfont[BoldFont={Noto Sans CJK JP Bold}]{Noto Sans CJK JP}
\usepackage{amsmath,amssymb,amsthm,mathtools}
\usepackage[margin=24mm]{geometry}
\usepackage{xcolor,booktabs,longtable,array,enumitem,fancyhdr,titlesec}
\usepackage{tocloft}\setlength{\cftsecnumwidth}{3.2em}\setlength{\cftpartnumwidth}{3.2em}
\usepackage[colorlinks=true,linkcolor=blue!50!black,citecolor=blue!50!black,urlcolor=blue!50!black]{hyperref}
\setlength{\parindent}{1em}
\setlength{\parskip}{2pt}
\linespread{1.17}
\emergencystretch=4em
\setlist{itemsep=1pt,topsep=3pt}
\setlength{\tabcolsep}{4pt}

%%% --- Japanese line breaking for XeTeX (no xeCJK available) ---
\makeatletter
\XeTeXinterchartokenstate=1
\newXeTeXintercharclass\cjkc
\newXeTeXintercharclass\cjkclose
\newXeTeXintercharclass\cjkopen
\def\setcjkrange#1#2{\@tempcnta=#1 \@whilenum\@tempcnta<#2 \do{\XeTeXcharclass\@tempcnta=\cjkc \advance\@tempcnta by 1 }}
\setcjkrange{"3000}{"30FF}
\setcjkrange{"4E00}{"9FFF}
\setcjkrange{"FF00}{"FFEF}
\setcjkrange{"2010}{"2027}
\setcjkrange{"25A0}{"25FF}
\setcjkrange{"2190}{"21FF}
\@tempcnta=0
\def\setclose#1{\XeTeXcharclass#1=\cjkclose}
\def\setopen#1{\XeTeXcharclass#1=\cjkopen}
\setclose{"3001}\setclose{"3002}\setclose{"FF0C}\setclose{"FF0E}\setclose{"FF09}\setclose{"300D}\setclose{"300F}\setclose{"3011}\setclose{"FF1A}\setclose{"FF1B}\setclose{"FF01}\setclose{"FF1F}\setclose{"30FC}
\setopen{"FF08}\setopen{"300C}\setopen{"300E}\setopen{"3010}
\def\cjkglue{\hskip0pt plus 0.25pt minus 0.1pt\relax}
\XeTeXinterchartoks\cjkc\cjkc={\cjkglue}
\XeTeXinterchartoks\cjkclose\cjkc={\cjkglue}
\XeTeXinterchartoks\cjkclose\cjkopen={\cjkglue}
\XeTeXinterchartoks\cjkc\cjkopen={\cjkglue}
\XeTeXinterchartoks\cjkclose\cjkclose={}
\XeTeXinterchartoks\cjkc\cjkclose={}
\XeTeXinterchartoks 0 \cjkc={\cjkglue}
\XeTeXinterchartoks \cjkc 0={\cjkglue}
\XeTeXinterchartoks 0 \cjkopen={\cjkglue}
\XeTeXinterchartoks \cjkclose 0={\cjkglue}
\makeatother

%%% --- status tags ---
\newcommand{\tP}{\textcolor{green!45!black}{\textbf{【証明済】}}}
\newcommand{\tC}{\textcolor{blue!55!black}{\textbf{【文献】}}}
\newcommand{\tV}{\textcolor{violet!80!black}{\textbf{【計算で確認】}}}
\newcommand{\tK}{\textcolor{red!70!black}{\textbf{【予想】}}}
\newcommand{\tQ}{\textcolor{orange!85!black}{\textbf{【要確認】}}}

%%% --- theorem-like environments ---
\theoremstyle{definition}
\newtheorem{defn}{定義}
\newtheorem{thm}[defn]{定理}
\newtheorem{prop}[defn]{命題}
\newtheorem{lem}[defn]{補題}
\newtheorem{cor}[defn]{系}
\newtheorem{ex}[defn]{例}
\newtheorem{rem}[defn]{注意}
\newtheorem{obs}[defn]{観察}
\newtheorem{conj}[defn]{予想}
\newcommand{\src}[1]{\par\noindent{\small\textcolor{blue!55!black}{\textbf{[出典]}~#1}}\par}
\renewcommand{\proofname}{証明}

\newenvironment{hyptab}{\par\medskip\noindent{\small\textbf{仮定と、その役割}}\par\nopagebreak\small\noindent\begin{tabular}{>{\raggedright\arraybackslash}p{3.3cm}>{\raggedright\arraybackslash}p{5.5cm}>{\raggedright\arraybackslash}p{5.9cm}}\toprule 仮定 & なぜ必要か & 証明のどこで効くか\\\midrule}{\\\bottomrule\end{tabular}\par\medskip}

%%% --- notation ---
\newcommand{\R}{\mathbb R}\newcommand{\C}{\mathbb C}\newcommand{\Z}{\mathbb Z}
\newcommand{\bT}{\mathbb T}
\newcommand{\Lb}{\overline L}
\newcommand{\Jc}{\mathcal J}
\newcommand{\dbar}{\bar\partial}
\newcommand{\ch}{\mathrm{ch}}\newcommand{\Td}{\mathrm{Td}}
\newcommand{\rk}{\mathrm{rk}}
\newcommand{\HH}{\mathrm{HH}}
\newcommand{\Res}{\mathrm{Res}}
\newcommand{\dlog}{\mathrm{dlog}\,}
\newcommand{\sh}{\sharp}

\pagestyle{fancy}\fancyhf{}
\fancyhead[L]{\small 一般化複素幾何の三種の神器：統合論文}
\fancyhead[R]{\small 研究ノート 2026-11-22（統合版）}
\fancyfoot[C]{\small\thepage}
\titleformat{\section}{\large\bfseries\sffamily}{\thesection}{0.8em}{}
\titleformat{\subsection}{\normalsize\bfseries\sffamily}{\thesubsection}{0.8em}{}


\newcommand{\Q}{\mathbb Q}\newcommand{\N}{\mathbb N}
\newcommand{\Hc}{H_c}
\newcommand{\we}{\wedge}
\setcounter{tocdepth}{1}
\newcommand{\oldname}[1]{\textsf{\small（総合版：#1）}}

\renewcommand{\contentsname}{目次}\renewcommand{\abstractname}{概要}

\begin{document}
\sloppy
\begin{center}
{\LARGE\bfseries\sffamily 一般化複素幾何の「三種の神器」へのアプローチ}\\[6pt]
{\Large\bfseries\sffamily 正則 Poisson 構造の例から（統合論文）}\\[12pt]
{\large 調査・執筆：Claude\ ／\ 問題設定：森}\\[3pt]
2026年11月22日（統合版 v1。原稿は日本語、英語版は別ファイル）\\[3pt]
{\small （査読を受けていない研究ノート。Part I--V の総合版を主題別に書き直した。\\未確認の箇所は本文中に\tQ\ で明示してある。）}
\end{center}
\begin{abstract}\noindent
複素幾何の三種の神器（Riemann--Roch、Serre 消滅、Kodaira 埋め込み）の一般化複素幾何版を、最も簡単な非自明な例（正則 Poisson 構造 $\sigma$ の $\beta$-変換で得る、型が跳ぶ一般化複素構造と、その上の一般化正則直線束 $L$）から調べる。
(1) 一般化正則構造は Poisson 加群であり、$\sigma$ が対数シンプレクティックなら平坦な対数接続で、留数 $r_i$ が定数になる。
(2) Serre 型の消滅は横断次数での消滅として成り立ち、全次数では成り立たない。双対性 $\HH^k(L)^*\cong\HH^{2n-k}(L^\vee\otimes K^2)$ が任意次元で成り立つ。
(3) 全 Euler 標数は $(-1)^n\rk(E)\chi(X)$ で $c_1$ に依らない。曲面・トーリックでは Poincar\'e 多項式を留数で書き下し、準斉次な特異点の局所公式を証明した。
(4) 一般次元は簡約定理により商複体の条件に帰着し、留数が非整数なら $H^k(U;L_U)=0\ (k\ne n)$ を示した。
(5) 正則型では次数ごとに $m$ の多項式、対数シンプレクティック型では周期的となる。
(6) Kodaira 型では、一般化正則切断による埋め込みは不可能であり、葉空間版とトーリックの Poisson--Kodaira を示した。
指数定理（重み付き b-計算）の一般化は予想にとどまり、一般の Poisson--Kodaira は未解決である。
\end{abstract}
\tableofcontents
\clearpage
%==============================================================
\section{序論と主結果}\label{sec:intro}
%==============================================================
\subsection{動機と問題設定}
複素幾何の「三種の神器」とは、Hirzebruch--Riemann--Roch の定理（HRR）、Serre の消滅定理、Kodaira の埋め込み定理である。これらを一般化複素幾何（Hitchin, Gualtieri）に拡張できるかが、本研究の最終目標である。調べた範囲では、Hu \cite{Hu1,Hu2} による Lie 亜群上の Kodaira--Nakano 消滅と Riemann--Roch--Hirzebruch、Pal--Poddar \cite{PP} による強一般化正則束（正則型）の消滅定理はあるが、一般の一般化複素構造に対する Kodaira 埋め込みや Riemann--Roch は見つからなかった\tC。

そこで「最も簡単な非自明な例」から攻める。選んだ例は、複素多様体 $X$ 上の正則 Poisson 構造 $\sigma$ の $\beta$-変換で得る一般化複素構造 $\Jc_\sigma$ である。この例には次の利点がある。(i) $\sigma$ が消える因子 $D$ の上で型が跳ぶ（複素型と非退化なシンプレクティック型が共存する）。(ii) 一般化正則直線束が Poisson 加群（平坦 $A$-接続をもつ正則束）として具体的に記述でき、一般化 Dolbeault コホモロジーが正則 Poisson 超コホモロジーとして計算できる。本論文は、この例について得られた結果を主題別に 1 本にまとめたものである。

\paragraph{本論文の方針。}
(a) 本論文は自己完結している：証明を載せた主張はすべて本文で証明し、紙幅のために証明を省いた主張は、省いた旨と根拠（スクリプトや計算の名前）を明記して\tV に留める。(b) すべての主張に状態タグ（\S\ref{subsec:tags}）を付ける。状態は、本文中に書いた証明の範囲を超えて格上げしていない（予想を証明済みにしていない）。(c) 定義、仮定、「その仮定がなぜ必要で証明のどこで効くか」（「仮定と、その役割」の表）、記号表を維持する。(d) 定理・補題等の番号は論文全体の通し番号である。

\subsection{主結果の概観}
以下、$X$ はコンパクト複素多様体（$\dim_\C X=n$）、$\sigma$ は正則 Poisson 構造、$E$ は一般化正則（Poisson 加群）束、$\HH^k(E)$ は一般化 Dolbeault コホモロジー、$P_E(t)=\sum_k\dim\HH^k(E)\,t^k$ は Poincar\'e 多項式である。$\sigma$ が単純正規交差因子 $D$ に沿って対数シンプレクティックなら、$E$ は $D$ に沿った平坦な対数接続をもち、その留数 $r_i$ は定数である。

\begin{enumerate}
\item[(1)] \textbf{準備（\S\ref{sec:prelim}）。}一般化正則束は Poisson 加群であり、$\HH^*(E)=H^*(X;(\wedge^\bullet T\otimes E,\ [\sigma,\cdot]+V\wedge))$ と計算できる。$\sigma$ が対数シンプレクティックなら Poisson 加群は対数的平坦接続で、留数は定数、$c_1(E)=-\sum r_i[D_i]$（定理 \ref{thm:main1}）。一般化正則切断は平坦（Casimir 型）切断である。
\item[(2)] \textbf{Serre 型消滅と双対性（\S\ref{sec:serre}）。}$X$ が射影的で $L_0$ が豊富な Poisson 加群なら、$m\gg0$ で横断次数 $a>0$ のコホモロジーがすべて消える（定理 \ref{thm:serre}）。全次数での消滅は $\sigma=0$ でも成り立たない。双対性 $\HH^k(L)^*\cong\HH^{2n-k}(L^\vee\otimes K_X^2)$ が任意次元・任意の $\sigma$ で成り立つ（定理 \ref{thm:dual}）。
\item[(3)] \textbf{Euler 標数（\S\ref{sec:euler}）。}全 Euler 標数は $(-1)^n\rk(E)\chi(X)$ で $c_1(E)$ に依らない（定理 \ref{thm:E}）。対数 de Rham 複体の Euler 標数は $\chi(U)$、$U=X\setminus D$（定理 \ref{thm:K1}）。
\item[(4)] \textbf{対数シンプレクティック曲面（\S\ref{sec:surface}）。}$r_i\neq1$ なら $\HH^*(E)=H^*(X;\mathcal F)\oplus\C^\mu[-2]$（$\mathcal F$ は $L_U$ の $Rj_*$/$j_!$ 型の拡張、$\mu$ は二重点の数）、Poincar\'e 多項式の閉公式（定理 \ref{thm:main2}, \ref{thm:B}）、$m$ について周期的に安定化（定理 \ref{thm:A2}）。
\item[(5)] \textbf{トーリックと局所項（\S\ref{sec:toric}）。}トーリック多様体の weight 公式、一般の $\sigma$ での閉公式（定理 \ref{thm:T}）、留数が不等な曲面（定理 \ref{thm:C2}）。準斉次な孤立特異点をもつ平面曲線の茎の公式を $s\neq1$ について証明した（定理 \ref{thm:local}）。
\item[(6)] \textbf{一般次元の簡約定理と定理V（\S\ref{sec:reduction}）。}一般次元の問題は、商複体 $Q$ の条件 (GEN$_Q$) に帰着する（定理 \ref{thm:R}）。留数がすべて非整数で $U$ がアフィンなら $H^k(U;L_U)=0\ (k\ne n)$（定理 \ref{thm:V}）。
\item[(7)] \textbf{正則型との比較（\S\ref{sec:regular}）。}正則型（平坦なシンプレクティックファイバー束）では次数ごとの次元が $m$ の多項式（系 \ref{cor:L}）、型が跳ぶ場合は周期的である。
\item[(8)] \textbf{Kodaira 型（\S\ref{sec:kodaira}）。}一般化正則切断による埋め込みは不可能、葉空間版の Kodaira、トーリックでの Poisson--Kodaira を証明した（命題 \ref{prop:one}, 定理 \ref{thm:K1leaf}, \ref{thm:PKtoric}）。
\item[(9)] \textbf{RR 型公式と指数予想（\S\ref{sec:RR}）。}Poincar\'e 多項式の公式 $P(t)=\chi(X)t^n+\sum_{k<n}h^k(U;L_U)(t^k-(-1)^kt^n)$（定理 \ref{thm:K}）。指数定理（重み付き $b$-計算）による統一は予想にとどまる（予想 \ref{conj:index}）。
\end{enumerate}

\subsection{主張の状態表}\label{subsec:ledger}
\begin{longtable}{>{\raggedright\arraybackslash}p{5.3cm}>{\raggedright\arraybackslash}p{2.6cm}>{\raggedright\arraybackslash}p{7.2cm}}
\caption{主な主張の状態（タグの意味は \S\ref{subsec:tags}）}\label{tab:ledger}\\
\toprule 主張 & 状態 & 主な仮定・備考\\\midrule\endfirsthead
\toprule 主張 & 状態 & 主な仮定・備考\\\midrule\endhead
\bottomrule\endfoot
\footnotesize 命題 \ref{prop:ident}（GH 束 $=$ Poisson 加群、$\HH$ の同定） & \tC\ \tP\ \tQ & LGSX Thm.~4.19 と Mokri の定理に依存し、再導出していない\\
\footnotesize 定理 \ref{thm:main1}（Poisson 加群 $=$ 対数的平坦接続） & \tP & 任意の偶数次元、SNC。留数公式 $c_1=-\sum r_i[D_i]$ は \tC\\
\footnotesize 定理 \ref{thm:serre}（Serre 型消滅、横断次数） & \tP & $X$ 射影的、$\sigma$ 任意。全次数では偽（注意 \ref{rem:total}）\\
\footnotesize 定理 \ref{thm:dual}（双対性、任意次元） & \tP\ \tV\ \tQ & 補題 \ref{lem:D}, \ref{lem:E} は $n\le4$ で機械確認。$\dbar$ 部分と $d$ 部分の Koszul 符号の整合は手計算のみ\\
\footnotesize 定理 \ref{thm:E}（全 Euler 標数） & \tP & HRR のみ。$c_1(E)$ に依らない\\
\footnotesize 定理 \ref{thm:K1}（対数版 $\chi(U)$） & \tP\ \tV & SNC、平坦対数接続\\
\footnotesize 定理 \ref{thm:main2}（曲面の主定理） & \tP\ \tV & $r_i\ne1$。命題 \ref{prop:ident} に依存\\
\footnotesize 定理 \ref{thm:B}（曲面の閉公式） & \tP\ \tV & (H1)--(H4)。(H4) は 8 例で確認、一般には $V^1(U)$ の条件\\
\footnotesize 定理 \ref{thm:A2}（周期的安定化） & \tP & 定理 \ref{thm:main2} と $X$ 単連結、$U$ アフィン\\
\footnotesize 定理 \ref{thm:weight}, \ref{thm:T}, \ref{thm:C2}（トーリック） & \tP\ \tV & (V)（$L_0^m$, $m\gg0$ なら定理 \ref{thm:serre}）、(A), (GEN), (R)\\
\footnotesize 観察 \ref{obs:res}（共鳴 $r_i=1$ の公式） & \tV\ \tK & $N$ の定義の転記は\tQ。破れる場合あり\\
\footnotesize 定理 \ref{thm:local}（準斉次な局所公式、$s\ne1$） & \tP\ \tQ & $s$ の符号規約と収束の扱いが\tQ（総合版では観察 \tV\ \tK。更新、巻末の変更点）\\
\footnotesize 局所公式（非準斉次、$s=1$） & \tK & 未証明\\
\footnotesize 予想 \ref{conj:refined}、観察 \ref{obs:triple}, \ref{obs:nr}、予想 \ref{conj:NR} & \tV\ \tK & データへの当てはめ。三重点・非被約\\
\footnotesize 定理 \ref{thm:R}（簡約定理） & \tP & 仮定 (GEN$_Q$)。本論文で濾過に依らない形に書き直した\\
\footnotesize 命題 \ref{prop:grQ}（(GEN$'$) $\Rightarrow$ (GEN$_Q$)、一般 $n$ の $\mathrm{gr}\,Q$） & \tQ & 一般 $n$ の局所計算は見取り図。$n=2$ は補題 \ref{lem:Q}\\
\footnotesize 定理 \ref{thm:V}（定理V：$H^k(U;L_U)=0$） & \tP & SNC、$U$ アフィン、すべて $r_i\notin\Z$。混合（一部の $r_i\in\Z$）は\tK\\
\footnotesize 定理 \ref{thm:K}（RR 型 Poincar\'e 多項式） & \tP\ \tV & 定理 \ref{thm:R} の仮定、$U$ アフィン、$n$ 偶数\\
\footnotesize 定理 \ref{thm:L}、系 \ref{cor:L}（正則型） & \tP\ \tQ & hard Lefschetz。$\partial$ と Gauss--Manin の同定は概略\\
\footnotesize 命題 \ref{prop:one}, \ref{prop:leaf}、定理 \ref{thm:K1leaf}（Kodaira 型） & \tP & 不可能性と葉空間版\\
\footnotesize 定理 \ref{thm:PKtoric}（トーリック Poisson--Kodaira） & \tP & 観察 \ref{obs:P2}（$\mathbb P^2$, $m=2$）は数値的存在 \tV\\
\footnotesize 予想 \ref{conj:PK}（一般の Poisson--Kodaira） & \tK & $\mathbb P^2$, $m\ge3$ は未計算\\
\footnotesize 予想 \ref{conj:index}（指数予想） & \tK & 未証明\\
\end{longtable}

\subsection{三種の神器との対応}
\begin{center}\small
\begin{tabular}{>{\raggedright\arraybackslash}p{2.4cm}>{\raggedright\arraybackslash}p{6.5cm}>{\raggedright\arraybackslash}p{6.0cm}}
\toprule 神器 & 成り立った形 & 成り立たない／未解決\\\midrule
Riemann--Roch & 全 Euler 標数 $(-1)^n\rk(E)\chi(X)$（定理 \ref{thm:E}）、対数版 $\chi(U)$（定理 \ref{thm:K1}）、Poincar\'e 多項式の公式（曲面：定理 \ref{thm:B}、一般次元：定理 \ref{thm:K}）、正則型の次数別多項式（系 \ref{cor:L}） & $c_1(E)$ に依存する素朴な HRR は不成立。指数定理による統一は未証明（予想 \ref{conj:index}）\\\midrule
Serre 消滅 & 横断次数 $a>0$ での消滅（$m\gg0$、$X$ 射影、$\sigma$ 任意、定理 \ref{thm:serre}） & 全次数での消滅は $\sigma=0$ でも不成立\\\midrule
Kodaira 埋め込み & 葉空間版（定理 \ref{thm:K1leaf}）、トーリックでの Poisson--Kodaira（定理 \ref{thm:PKtoric}） & 一般化正則切断による埋め込みは不可能（命題 \ref{prop:one}）。一般の Poisson--Kodaira は予想（予想 \ref{conj:PK}）。一般の一般化複素構造への一般化は未着手\\\bottomrule
\end{tabular}
\end{center}

\subsection{全体を貫く 3 つの考え}
(i) \emph{平坦接続への翻訳。}$\beta$-変換構造上の一般化 Dolbeault 複体は $(\wedge^\bullet T\otimes E,\ \dbar+[\sigma,\cdot])$ と同一視でき\tC（Laurent-Gengoux--Sti\'enon--Xu \cite{LGSX1}）、一般化正則束は Poisson 加群である。$\sigma$ が $D$ に沿って対数シンプレクティックなら、これは $D$ に沿った平坦な対数接続になり、留数が定数になる。
(ii) \emph{局所と大域の分離。}対数 de Rham 複体と $Rj_*L_U$ の比較（Deligne \cite{Del}）により $U=X\setminus D$ の局所系のコホモロジーが現れ、留数が整数のとき（共鳴）に商複体の寄与が出る。大域は $H^*(U;L_U)$ と $\chi(X)$、局所は $D$ の特異点の茎の公式で書ける。
(iii) \emph{3 つの振る舞い。}正則型では次数ごとの次元が $m$ の多項式、対数シンプレクティック型では周期的、全 Euler 標数はどちらでも $m$ に依らない（\S\ref{sec:regular}, \S\ref{sec:RR}）。

\subsection{タグの意味、記号表、読み方}\label{subsec:tags}
\paragraph{状態タグ。}\tC\ 文献からの引用（原典での確認の範囲は個々に記す）。\tP\ 本論文に証明を書いたもの（引用した結果への依存は明記）。\tV\ 厳密な線形代数（有理数または有限体 $\mathbb F_{32003}$, $\mathbb F_{2^{61}-1}$）などによる計算機確認で、証明ではない。\tK\ 予想・見通し。\tQ\ 未確認の点（記憶による引用、導出が見取り図、符号規約の未確認など）。複数付く場合は、主張の各部分に対応する。

\paragraph{読み方。}\S\ref{sec:prelim}--\S\ref{sec:euler} が基礎である。\S\ref{sec:surface} は曲面の完全な結果で、\S\ref{sec:toric} のトーリック・局所項、\S\ref{sec:reduction} の一般次元がそれに続く。\S\ref{sec:regular}（正則型）、\S\ref{sec:kodaira}（Kodaira 型）は互いにほぼ独立で、\S\ref{sec:RR} は \S\ref{sec:reduction} と \S\ref{sec:euler} を使う。

\paragraph{記号表。}
\begin{longtable}{>{\raggedright\arraybackslash}p{3.5cm}>{\raggedright\arraybackslash}p{11.3cm}}
\toprule 記号 & 意味\\ \midrule\endhead
\bottomrule\endfoot
\footnotesize $X,\ n$ & コンパクト複素多様体、$\dim_\C X=n$（対数シンプレクティックでは $n$ は偶数）\\
\footnotesize $T,\ T^*,\ \Omega^p,\ K_X$ & 正則接束、余接束、正則 $p$-形式の層、$K_X=\Omega^n$\\
\footnotesize $\bT=T\oplus T^*$ & 一般化接束。内積 $\langle X+\xi,Y+\eta\rangle=\tfrac12(\xi(Y)+\eta(X))$、Courant（Dorfman）括弧 $[X+\xi,Y+\eta]=[X,Y]+\mathcal L_X\eta-\iota_Yd\xi$\\
\footnotesize $\Jc,\ L$ & 一般化複素構造と、その $+i$ 固有束（$\bT_\C$ の極大迷零部分束、複素 Lie 亜群）。$L^*\cong\Lb$。直線束には $E$, $L_0$, $\mathcal O(m)$ 等も用いる（文脈で区別）\\
\footnotesize $\sigma,\ \sigma^\sh,\ \Jc_\sigma$ & 正則 Poisson 構造（$\sigma\in H^0(\wedge^2T)$, $[\sigma,\sigma]=0$）、$\sigma^\sh:T^*\to T$, $\xi\mapsto\sigma(\xi,\cdot)$、その $\beta$-変換 $L_\sigma=T^{0,1}\oplus\{\sigma^\sh\xi+\xi:\xi\in T^{*1,0}\}$\\
\footnotesize $A=(T^*)_\sigma$ & 正則 Lie 亜群：アンカー $\sigma^\sh$、括弧 $[a,b]_\sigma=\mathcal L_{\sigma^\sh a}b-\mathcal L_{\sigma^\sh b}a-d\sigma(a,b)$\\
\footnotesize Poisson 加群 & 平坦 $A$-接続 $\nabla$ をもつ正則ベクトル束 $(E,\nabla)$（$=$ 一般化正則束、命題 \ref{prop:ident}）\\
\footnotesize $V$ & 局所枠 $e$ に対し $\nabla_\beta e=\beta(V)e$ で定まる正則ベクトル場（Poisson ベクトル場）\\
\footnotesize $C(A;E)$、$\HH^k(E)$ & $(\wedge^bT\otimes E,\ d_E=[\sigma,\cdot]+V\wedge)$（$A$ の Chevalley--Eilenberg 複体）、$\HH^k(E)=H^k(X;C(A;E))$\\
\footnotesize $E_1^{a,b}$ & $H^a(X,\wedge^bT\otimes E)\Rightarrow\HH^{a+b}$。$a$ を横断次数と呼ぶ\\
\footnotesize $P_E(t),\ [\cdots]$ & $P_E(t)=\sum_k\dim\HH^k(E)t^k$、表では $[\dim\HH^0,\dots,\dim\HH^{2n}]$ と略記\\
\footnotesize $D=\sum D_i,\ U$ & $\sigma$ の退化因子（簡約、SNC、曲面では通常二重点のみも可）、$U=X\setminus D$。$D_T=\bigcap_{i\in T}D_i$、$D_{(n)}$ は $n$ 重点の集合\\
\footnotesize $T_D,\ B,\ \mathrm{DR}_{\log}(E)$ & $T_D=T_X(-\log D)$、$B=\Omega^1(\log D)=T_D^*$、対数 de Rham 複体 $(\Omega^\bullet(\log D)\otimes E,\nabla^{\log})$\\
\footnotesize 対数シンプレクティック & $\sigma\in H^0(\wedge^2T_D)$ かつ $\sigma^\sh:B\to T_D$ が同型\\
\footnotesize $r_i,\ L_U$ & $D_i$ に沿った留数（定数）、$U$ 上の平坦切断の局所系（階数 1、$D_i$ まわりのモノドロミー $e^{\mp2\pi\sqrt{-1}r_i}$）\\
\footnotesize $Q$ & 商複体 $C(A;E)/C(T_D;E)$。$D$ 上に台をもつ\\
\footnotesize $\mu,\ c,\ \rho_D,\ b_1$ & （曲面）二重点の個数 $\mu$、$D$ の既約成分数 $c$、$[D_i]\in H^2(X;\C)$ の張る空間の次元 $\rho_D$、$b_1=c-\rho$\\
\footnotesize $j_1,j_2,\ \mathcal F$ & 開埋め込み $U\hookrightarrow X'\hookrightarrow X$、$\mathcal F=j_{2!}Rj_{1*}L_U$（定理 \ref{thm:main2}）\\
\footnotesize $N,M,\Sigma,u_i$ & トーリック：格子 $N\cong\Z^n$、双対 $M$、ファンと光線生成元\\
\footnotesize $\rho,\ A,\ \chi,\ \Lambda$ & トーリック：$\nabla=d+\rho$（$\rho\in M_\C$）、定数歪対称行列 $A$、指標 $\chi\in M$、weight 公式の指標集合\\
\footnotesize $h,\ (w_z,w_w),\ N_h,\ \delta$ & 準斉次な孤立特異点の定義関数、重み、次数、$\delta=N_h-w_z-w_w$（\S\ref{subsec:local}）\\
\footnotesize $\mu_p,\ \mathrm{Sp},\ \mathrm{lct}$ & Milnor 数、スペクトル、log canonical threshold $(w_z+w_w)/N_h$\\
\footnotesize $s$ & 局所モデル $\mathcal O h^s$ の留数パラメータ\\
\footnotesize $Y\to M\to X$、$V^j$ & 平坦なシンプレクティックファイバー束、ファイバーコホモロジーの Gauss--Manin 束（\S\ref{sec:regular}）\\
\footnotesize $\partial U$ & $D$ の管状近傍の境界（実 $2n-1$ 次元の閉多様体、$\chi(\partial U)=0$）\\
\end{longtable}

%==============================================================
\section{準備：一般化正則直線束と一般化 Dolbeault コホモロジー}\label{sec:prelim}
%==============================================================
\subsection{古典の三種の神器：主張と、仮定が証明のどこで効くか}\label{subsec:classical}
一般化の目標を明確にするため、古典版の主張と証明の骨格を整理する。内容は標準的で新しい結果ではない\tC。

\begin{itemize}
\item \textbf{HRR。}$M$ コンパクト複素多様体、$E$ 正則ベクトル束とすると $\chi(M,E)=\int_M\ch(E)\Td(T^{1,0}M)$。\src{Hirzebruch, \emph{Neue topologische Methoden in der algebraischen Geometrie}, Springer 1956（射影的な場合）；Atiyah--Singer, Ann.\ Math.\ 87 (1968) 546--604（一般のコンパクト複素多様体、Dolbeault 複体の指数）；Huybrechts \cite{Huy} \S5.1。}
\item \textbf{Serre の消滅定理。}$X$ 射影的、$\mathcal O(1)$ 非常に豊富、$\mathcal F$ 連接のとき、ある $m_0$ があり $m\ge m_0$ で $H^q(X,\mathcal F(m))=0\ (q>0)$。\src{Serre \cite{Serre}；Hartshorne \cite{Hart} Thm.\ III.5.2。}
\item \textbf{Kodaira--Akizuki--Nakano 消滅。}$M$ コンパクト Kähler、$E$ 正の正則 Hermite 直線束のとき $p+q>n$ ならば $H^q(M,\Omega^p\otimes E)=0$。\src{Kodaira \cite{Kod53}；Akizuki--Nakano \cite{AN}；Griffiths--Harris \cite{GH} Ch.~1 \S2。}
\item \textbf{Kodaira の埋め込み定理。}$M$ コンパクト、$L$ 正の正則直線束ならば、$M$ は射影的で、$m\gg0$ で $|L^m|$ による写像 $M\to\mathbb P(H^0(M,L^m)^*)$ は埋め込み。\src{Kodaira \cite{Kod54}；Griffiths--Harris Ch.~1 \S4；Huybrechts \S5.3。}
\end{itemize}

\begin{center}\small
\begin{tabular}{>{\raggedright\arraybackslash}p{2.3cm}>{\raggedright\arraybackslash}p{2.9cm}>{\raggedright\arraybackslash}p{4.4cm}>{\raggedright\arraybackslash}p{4.6cm}}
\toprule 古典の定理 & 仮定 & なぜ必要か・証明のどこで効くか & 本論文での運命\\\midrule
HRR & コンパクト、正則構造 & Dolbeault 複体が楕円型で Fredholm（指数が定義できる）、$\dbar^2=0$。シンボルから $\Td$ が現れ、$E$ は $\ch(E)$ として入る & 全 Euler 標数は $c_1$ を見ない（定理 \ref{thm:E}）。細分が必要（\S\ref{sec:RR}）\\\midrule
Serre & 射影的、豊富、連接 & $\bigoplus\mathcal O(-N_i)\twoheadrightarrow\mathcal F$ と $\mathbb P^N$ の \v Cech 計算、下降帰納法 & 連接層 $\wedge^bT\otimes E$ に適用：横断次数で成立（定理 \ref{thm:serre}）\\\midrule
KAN & コンパクト、Kähler、正 & Hodge 理論、Kähler 恒等式 $\Rightarrow$ BKN 恒等式、曲率項 $[i\Theta_E,\Lambda]$ の正値性 & 一般化 Kähler では直線束係数版は未記述。正則型は Pal--Poddar、Lie 亜群は Hu（\S\ref{subsec:compare}）\\\midrule
Kodaira 埋め込み & コンパクト、正 & 爆発 $\tilde M$ 上の KAN 消滅から $H^0(L^m)\to L^m\otimes\mathcal O/\mathfrak m_x^2$ の全射性（点と接ベクトルを分離する切断が多い） & 一般化正則切断は平坦で点を分離できない（\S\ref{sec:kodaira}）\\\bottomrule
\end{tabular}
\end{center}

\begin{rem}\tK\ （一般化への示唆）
古典では (1) Hodge 理論（楕円型）、(2) Kähler 恒等式 $\Rightarrow$ BKN 恒等式、(3) 曲率の正値性、(4) 大域切断の多さ、の 4 点が役割を分担している。一般化複素幾何では (1) は楕円型複体 $(\wedge L^*\otimes E,d_L)$ として成立する（Gualtieri \cite{Gua11}）。(2), (3) の直線束係数版は一般化 Kähler の枠組みではまだ書かれていない（調べた範囲）。(4) は、一般化正則切断が Casimir 型の平坦切断になる（観察 \ref{obs:casimir}）ため、古典と同じ形では成り立たない。
\end{rem}

%--------------------------------------------------------------
\subsection{一般化複素構造と $\beta$-変換}
%--------------------------------------------------------------
\begin{defn}[一般化複素構造・型]\tC\label{defn:GC}
$\Jc:\bT\to\bT$ が $\Jc^2=-1$ かつ内積を保つとき一般化概複素構造という。$+i$ 固有束 $L\subset\bT_\C$ が Courant 括弧で閉じるとき可積分（一般化複素構造）という。$\mathrm{type}(\Jc)_p=\mathrm{codim}_\C\,\pi_T(L_p)$ を型といい、局所定数な点を正則点という。型は一般に跳ぶ。複素構造 $I$ は型 $n$（$L=T^{0,1}\oplus T^{*1,0}$）、シンプレクティック構造は型 $0$ である。
\src{Hitchin \cite{Hit03}；Gualtieri \cite{Gua11} \S3（型 Def.~3.5、標準束 Def.~3.7、一般化 Calabi--Yau Def.~3.17、Picard 群 Def.~3.18。番号は原文で確認）；Abouzaid--Boyarchenko \cite{AB}。}
\end{defn}
一般化 Dolbeault 複体 $(\wedge^\bullet L^*,d_L)$ は楕円型複体で、コンパクト多様体上でコホモロジーは有限次元である（Gualtieri \cite{Gua11} Prop.~3.13, Cor.~3.14\tC）。

\begin{prop}[$\beta$-変換]\tC\label{prop:Jsigma}
$X$ を複素多様体、$\sigma$ を正則 Poisson 構造とする。(a) $L_\sigma=T^{0,1}\oplus\{\sigma^\sh\xi+\xi:\xi\in T^{*1,0}\}$ は $\bT_\C$ の可積分な極大迷零部分束で、対応する一般化複素構造 $\Jc_\sigma$ の型は $p$ で $n-\mathrm{rk}_\C\sigma^\sh_p$。特に $\sigma$ が非退化な点では型 0、$\sigma=0$ の点では型 $n$。(b) $(\Xi,\xi)\mapsto(\Xi+\sigma^\sh\xi,\xi)$ は複素 Lie 亜群の同型 $T^{0,1}\bowtie(T^{1,0})^*_\sigma\to L_\sigma$（matched pair）。
\src{(a) Gualtieri \cite{GuaBr} Prop.~5；Laurent-Gengoux--Sti\'enon--Xu（LGSX）\cite{LGSX1} Thm.~2.7, Rem.~2.8（$\Jc_\pi$, $\pi=\pi_R+i\pi_I$）。型の公式は筆者の整理で、Gualtieri の例と整合する\tQ。(b) LGSX Thm.~4.21, Lemma 4.22（$\Jc_{4\pi}$ の因子 4 は $\pi=\pi_R+i\pi_I$ の規格化による）。固有値の符号の取り方は文献で異なるが複素共役で移り、以下に影響しない。}
\end{prop}

\begin{defn}[一般化正則束・Poisson 加群]\tC\label{defn:GH}
複素ベクトル束 $E$ に平坦 $L$-接続 $\dbar_E:\Gamma(E)\to\Gamma(L^*\otimes E)$（$\dbar_E^2=0$）が与えられたものを一般化正則束という。正則ベクトル束 $E$ と平坦 $A$-接続 $\nabla:\Gamma(E)\to\Gamma(A^*\otimes E)$（$\nabla_\xi(fs)=f\nabla_\xi s+(\sigma^\sh\xi)(f)s$、$[\nabla_a,\nabla_b]=\nabla_{[a,b]_\sigma}$、$a,b$ は正則 1-形式）の組を Poisson 加群という。
\src{Gualtieri \cite{Gua11} \S3（一般化正則束は Thm.~3.15 の後に散文で導入、番号付き定義はない）；Gualtieri \cite{GuaBr} \S4.1；LGSX Def.~4.15, Lemma 4.16。}
\end{defn}

\begin{defn}[一般化 Dolbeault コホモロジー $\HH$]\label{defn:HH}
Poisson 加群 $E$ に対し、Chevalley--Eilenberg 複体 $C(A;E)=(\wedge^\bullet T\otimes E,\ d_E)$（微分は $\C$-線形、局所正則枠 $e$ では $d_E=[\sigma,\cdot]+V\wedge$、$V$ は $\nabla_\beta e=\beta(V)e$ で定まる）の超コホモロジー $\HH^k(E):=H^k(X;C(A;E))$ を一般化 Dolbeault コホモロジーと呼ぶ。$b$ による濾過から、$E_1^{a,b}=H^a(X,\wedge^bT\otimes E)\Rightarrow\HH^{a+b}(E)$ なるスペクトル系列がある（$X$ コンパクトなら有限次元）。$P_E(t)=\sum_k\dim\HH^k(E)\,t^k$。
\end{defn}

\begin{prop}[一般化正則束 $=$ Poisson 加群、Dolbeault モデル]\label{prop:ident}
(a) $\Jc_\sigma$ に関する一般化正則束は、Poisson 加群構造をもつ正則ベクトル束と同じものである。\tC\ Gualtieri \cite{GuaBr} Prop.~8；LGSX Lemma 4.16。直接の証明の概略を下に示す\tP。
(b) $L^*=T^{*0,1}\oplus T^{1,0}$ より $\Gamma(\wedge^kL^*\otimes E)=\bigoplus_{a+b=k}A^{0,a}(\wedge^bT\otimes E)$ で、$d_L=\dbar+(-1)^ad_A$。したがって一般化 Dolbeault コホモロジー $H^k(\Gamma(\wedge^\bullet L^*\otimes E),d_L)$ は $\HH^k(E)$ に等しい。\tC\ LGSX Thm.~4.19（$H^*(A,E)=H^*(T^{0,1}\bowtie A^{1,0},E)$、$H^*(A,E)$ は Def.~4.17 で層コホモロジーとして定義）。\tP\ 証明の構造は下に示すが、matched pair に関する Mokri の定理（\cite{Mok}）と LGSX Lemma 4.22 に依存し、これらは再導出していない\tQ。
\end{prop}
\begin{proof}[証明の概略]
(a) ($\Leftarrow$) 正則枠 $\{e_s\}$ に対し $\nabla_{\bar\Xi}\sum f_se_s=\sum(\bar\Xi f_s)e_s$、$\nabla_\xi\sum f_se_s=\sum((\sigma^\sh\xi)f_s)e_s+f_s\nabla_\xi e_s$ とおく（Leibniz 則から強制される唯一の拡張）。平坦性は、$(\bar\Xi,\bar\Upsilon)$ 型は正則構造の可積分性、$(\bar\Xi,\xi)$ 型は $\nabla_\xi$ が正則切断を正則切断に写すこと、$(\xi,\eta)$ 型は Poisson 加群の平坦性（曲率はテンソルなので $C^\infty$ に拡張できる）による。($\Rightarrow$) $T^{0,1}$ への制限が正則構造を与え（Koszul--Malgrange）、混合型の平坦性から $\nabla_\xi$ は正則切断を正則切断に写し、Leibniz 則と平坦性が加群の条件を与える。
(b) $L\cong T^{0,1}\bowtie A^{1,0}$ の Chevalley--Eilenberg 複体は matched pair の一般論（Mokri）により、$T^{0,1}$ 方向（$\dbar$）と $A^{1,0}$ 方向（$d_A$）の二重複体の全複体になる。2 つの微分は Koszul 符号 $(-1)^a$ のもとで反可換で、全複体は $(\wedge^\bullet T\otimes E,d_A)$ の Dolbeault 分解なので、そのコホモロジーは超コホモロジーである。\qedhere
\end{proof}
以下、結果はすべて $\HH^*(E)=H^*(X;C(A;E))$ についての主張であり、一般化 Dolbeault コホモロジーとの同一視が必要な箇所では命題 \ref{prop:ident}(b) を使う。

\begin{obs}[一般化正則切断は平坦切断]\tP\label{obs:casimir}
$\HH^0(E)=\ker(d_E:H^0(E)\to H^0(T\otimes E))=\{s\in H^0(X,E):\nabla_\xi s=0\ \forall\xi\}$、すなわち平坦（Casimir 型）切断である。$\sigma$ が非退化な点の近くでは $\nabla$ は $\sigma^\sh$ により通常の平坦接続に対応し、平坦切断は局所定数になる。
\end{obs}
一般化正則切断が $X$ から射影空間への写像として埋め込みを与えられない理由はここにある（\S\ref{sec:kodaira}）。

%--------------------------------------------------------------
\subsection{対数シンプレクティックな場合：Poisson 加群 $=$ 対数的平坦接続}\label{subsec:logsymp}
%--------------------------------------------------------------
$X$ をコンパクト複素多様体（$\dim_\C X=n=2k$）、$D\subset X$ を簡約な因子（単純正規交差 SNC、曲面では通常二重点のみも許す）、$T_D=T_X(-\log D)$、$B=\Omega^1(\log D)$ とする。
\begin{defn}[対数シンプレクティック]\tC\label{defn:logsymp}
$\sigma\in H^0(\wedge^2T_D)$ かつ $\sigma^\sh:B\to T_D$ が同型であるとき、$\sigma$ は $D$ に沿って対数シンプレクティックであるという（同値：$\omega:=\sigma^{-1}\in H^0(\Omega^2(\log D))$ が非退化）。このとき $\mathrm{div}(\sigma^k)=D$、すなわち $-K_X=\mathcal O(D)$。曲面では、$\sigma\in H^0(-K_X)$ で零因子が簡約かつ通常二重点のみ、ということである。局所的に $\sigma=h\,\partial_1\wedge\partial_2$、$h=zu$（滑らかな点）または $z_1z_2u$（二重点）、$u$ は単元。
\src{Goto \cite{Goto}；Gualtieri--Li \cite{GL}；Pym \cite{Pym} Lett.\ Math.\ Phys.\ 108 (2018) 573--632（書誌は確認）。Goto と Gualtieri--Li の掲載情報は\tQ。}
\end{defn}

\begin{thm}[Poisson 加群 $=$ 対数的平坦接続]\tP\label{thm:main1}
$(X,D,\sigma)$ を対数シンプレクティック（任意の偶数次元、$D$ は SNC）、$E$ を Poisson 加群をもつ正則直線束とする。
\begin{enumerate}
\item[(a)] 局所枠 $e$ に対し $\nabla_\beta e=\beta(V)e$ で正則ベクトル場 $V$ が定まり $[\sigma,V]=0$（Poisson ベクトル場）。枠の取り替え $e\mapsto ge$ で $V$ は Hamilton 項 $X_{\log g}$ だけ変わる。
\item[(b)] $V$ は自動的に対数的：$V\in T_D$。したがって $A$-加群構造は一意に $B$-加群構造に延び、$\sigma^\sh:B\cong T_D$ を通して、対数極をもつ平坦接続 $\nabla^{\log}:E\to\Omega^1(\log D)\otimes E$ と同じものになる。
\item[(c)] $\nabla^{\log}$ の $D_i$ に沿った留数は定数 $r_i\in\C$ で、$H^2(X;\C)$ において $c_1(E)=-\sum_ir_i[D_i]$。
\item[(d)] $D$ の点の近くで $\nabla^{\log}e=\sum_ir_i\,\dlog z_i\,e$ となる枠が取れる。
\item[(e)] 部分複体 $C(T_D;E)=(\wedge^bT_D\otimes E,d_E)\subset C(A;E)$ は、$\wedge^b\sigma^\sh$ により対数 de Rham 複体 $\mathrm{DR}_{\log}(E)$ と同型。
\end{enumerate}
\end{thm}
\begin{hyptab}
対数シンプレクティック（$\sigma^\sh:B\cong T_D$） & $B$ 上の Lie 亜群構造 $(T_D^*)_\sigma$ と $T_D$ を同一視し、Poisson 加群を（対数極つき）平坦接続として扱うため & 補題 \ref{lem:iso}、(b), (e)\\\midrule
$D$ 簡約・SNC & 局所方程式 $f=z_1\cdots z_s$ の各因子から極の位数を独立に評価するため & 補題 \ref{lem:log}（$V(f)\in(f)\Rightarrow V(z_i)\in(z_i)$）\\\midrule
$X$ コンパクト & $c_1(E)$ と留数の関係は大域的なコホモロジー類の等式 & (c)。留数が定数であることは局所的\\\midrule
$[\sigma,\sigma]=0$ & $A$ が Lie 亜群になり $d_E^2=0$ & 補題 \ref{lem:flat}
\end{hyptab}

\begin{lem}[平坦性 $\Leftrightarrow$ Poisson ベクトル場]\tP\label{lem:flat}
局所枠 $e$ について $\Theta(\beta):=\nabla_\beta e/e$ とおくと $\nabla_a\nabla_be=(\sigma^\sh a)(\Theta b)e+\Theta(a)\Theta(b)e$。よって平坦性 $[\nabla_a,\nabla_b]=\nabla_{[a,b]_\sigma}$ は $(\sigma^\sh a)(\Theta b)-(\sigma^\sh b)(\Theta a)=\Theta([a,b]_\sigma)$、すなわち $\Theta$ が Lichnerowicz 複体の 1-コサイクル（$\Theta(\beta)=\beta(V)$, $[\sigma,V]=0$）と同値。曲面で $\sigma=h\partial_1\wedge\partial_2$、$a=\Theta(dz_1)$、$b=\Theta(dz_2)$ のとき条件は $h(a_1+b_2)=h_1a+h_2b$、すなわち $V(h)=h\,\mathrm{div}V$。微分は $C^0$ で $d_E(fe)=(-\sigma^\sh(df)+fV)e$、$c=X\otimes e$、$X=p\partial_1+q\partial_2$ に対し
\begin{equation}\label{eq:dE}
(d_Ec)(dz_1,dz_2)=h(p_1+q_2)-h_1p-h_2q+qa-pb.
\end{equation}
\end{lem}
\begin{lem}[$V$ は対数的]\tP\label{lem:log}
$D$ の局所方程式を $f=z_1\cdots z_s$ とする。$L_V\sigma=0$ から $L_V\sigma^k=0$。$\sigma^k=f'\partial_1\wedge\cdots\wedge\partial_n$（$f'=f\cdot$単元）なので $V(f')-f'\,\mathrm{div}V=0$、よって $V(f)\in(f)$。$V(f)/f=\sum_iV(z_i)/z_i$ の第 $i$ 項の極は $\{z_i=0\}$ のみで、全体が正則だから各 $V(z_i)/z_i$ が正則、すなわち $V(z_i)\in(z_i)$。
\end{lem}
\begin{lem}[$\sigma^\sh:B\to T_D$ は Lie 亜群の同型]\tP\label{lem:iso}
$\sigma\in\wedge^2T_D$ は Lie 亜群 $T_D$ 上の非退化 Poisson 構造である。$\omega=\sigma^{-1}$ が $d_{T_D}$-閉であることは、$U$ 上で $\omega$ が閉（Poisson 構造の逆）で、$\Omega^3(\log D)$ の切断が稠密開集合上で 0 なら 0 であることから従う。閉 2-形式による非退化 Poisson 構造では $\sigma^\sh$ が括弧を保つので $(T_D^*)_\sigma=B\cong T_D$。
\end{lem}
\begin{lem}[留数・局所標準形・$c_1$]\tP\ \tC\label{lem:res}
$\nabla^{\log}e=\theta e$、$\theta\in\Omega^1(\log D)$ は平坦性から閉。$D_i$ に沿った留数は、$d\theta$ を $D_i$ に制限して定数と分かる。$\theta-\sum r_i\dlog z_i$ は閉正則 1-形式で $=dg$、枠 $e'=e^{-g}e$ で $\nabla^{\log}e'=\sum r_i\dlog z_i\,e'$。$c_1(E)=-\sum r_i[D_i]$ は Esnault--Viehweg の留数公式による（\tC\ \cite{EV} \S2；\cite{Del} Ch.~II。符号の確認：標準切断が平坦な $\mathcal O(D)$ の留数は $-1$、$c_1=[D]$）。同じ $c_1$ をもつ加群構造の全体は $H^0(\Omega^1(\log D))=\ker(\C^c\to H^2(X;\C))$ 上のトーサー（Deligne）で、次元は $c-\rho_D$。$D$ が既約なら一意、可約なら留数ベクトルのアフィン族になる。
\end{lem}
\begin{proof}[定理 \ref{thm:main1} の証明]
(a) 補題 \ref{lem:flat}。(b) 補題 \ref{lem:log} で $\Theta$ が $B$ に延び、平坦性は稠密開集合上の等式なので延長後も成り立つ。補題 \ref{lem:iso} と合わせる。(c), (d) 補題 \ref{lem:res}。(e) 補題 \ref{lem:iso}。
\end{proof}

%==============================================================
\section{Serre 型消滅と双対性}\label{sec:serre}
%==============================================================
\subsection{Serre 型消滅（横断次数）}
古典の Serre 消滅（連接層 $\mathcal F$ と豊富な $\mathcal O(1)$）との対応は次の通りである\tK（対応づけは筆者の整理）。連接層 $\mathcal F$ に当たるのは Lie 亜群複体の各項 $K^b=\wedge^bT\otimes E$、豊富な $\mathcal O(1)$ に当たるのは豊富かつ Poisson 加群である正則直線束 $L_0$（Poisson 加群構造は $C(A;\cdot)$ が複体になるためだけに必要で、正値性には関与しない）、結論 $H^q(\mathcal F(m))=0\ (q>0)$ に当たるのは $E_1^{a,b}=H^a(X,\wedge^bT\otimes E\otimes L_0^m)=0$（$a>0$：\emph{横断（Dolbeault）次数}）である。

\begin{thm}[Serre 型定理（$\beta$-変換）]\tP\label{thm:serre}
$X$ をコンパクト複素射影多様体（$\dim_\C X=n$）、$\sigma$ を\emph{任意の}正則 Poisson 構造（非退化性や対数条件は不要）、$E$ を Poisson 加群（正則ベクトル束）、$L_0$ を Poisson 加群構造をもつ豊富な正則直線束とする。このとき $X,\sigma,E,L_0$ のみに依存する $m_0$ があって、$m\ge m_0$ ならば
\begin{enumerate}
\item[(a)] すべての $a>0$、すべての $b$ で $H^a(X,\wedge^bT\otimes E\otimes L_0^m)=0$、
\item[(b)] $\HH^k(E\otimes L_0^m)\cong H^k\bigl(\Gamma(X,\wedge^\bullet T\otimes E\otimes L_0^m),\,d\bigr)$。特に $k>n$ なら $\HH^k=0$、
\item[(c)] $\HH^0(E\otimes L_0^m)=$ 平坦（Casimir）正則切断全体。
\end{enumerate}
\end{thm}
\begin{hyptab}
$X$ 射影的、$L_0$ 豊富 & 古典 Serre の仮定そのもの & (a)：Hartshorne III.5.2 を有限個（$b=0,\dots,n$）の連接層 $\wedge^bT\otimes E$ に適用\\\midrule
$L_0$、$E$ が Poisson 加群 & $L_0^m$ にテンソル積の加群構造が入り $C(A;E\otimes L_0^m)$ が複体になる & 複体 $K^\bullet$ の定義。正値性とは無関係\\\midrule
$[\sigma,\sigma]=0$ & $d^2=0$ & 同上\\\midrule
（不要）$\sigma$ の非退化性・正規交差 & --- & この定理では使わない
\end{hyptab}
\begin{proof}
$L_0^m$ はテンソル積の加群構造（Leibniz 則）をもつ。Serre の定理（Hartshorne III.5.2；$L_0$ が豊富なだけのときは、$L_0^{m_1}$ が非常に豊富となる $m_1$ をとり、有限個の捻り $\mathcal F\otimes L_0^j\ (0\le j<m_1)$ にも適用する）を $n+1$ 個の連接層 $\wedge^bT\otimes E$ に適用し、$m_0$ をそれらの最大にとれば (a)。各 $K^b$ は $\Gamma$-非輪状なので、有界な複体 $K^\bullet=C(A;E\otimes L_0^m)$（微分は $\C$-線形）に対し $E_1^{a,b}=H^a(K^b)$ は $a=0$ に集中し、$R\Gamma(K^\bullet)=\Gamma(K^\bullet)$（非輪状分解の補題）。これが (b)。$b\le n$ なので $k>n$ で消える。(c) は $\Gamma$-複体の $H^0$ で、観察 \ref{obs:casimir} による。\qedhere
\end{proof}
\begin{rem}[全次数での消滅は偽]\tP\label{rem:total}
$\sigma=0$（複素構造）のとき $\HH^k(L)=\bigoplus_{a+b=k}H^a(X,\wedge^bT\otimes L)$ で、$m\gg0$ でも $H^0(\wedge^bT\otimes L^m)\ne0$ により $\HH^b\ne0\ (b=0,\dots,n)$。$\mathbb{CP}^2$ の $m=0$ では $[1,8,10,0,0]$。消えるのは横断次数 $a>0$ と全次数 $>n$ のみで、「全次数での Serre 消滅」は成り立たない。
\end{rem}
\begin{rem}[有効な閾値\ \tV]
Bott の公式から、$E_1^{a>0}=0$ となる条件は $\mathbb{CP}^2$（$E=\mathcal O(m)$）で $m\ge-2$（$m=-3$ では $(2,0),(1,1),(0,2)$ の 3 項がすべて非零）、$\mathbb{CP}^1\times\mathbb{CP}^1$（$L_0=\mathcal O(1,1)$）で $m\ge-1$（$m=-2$ は $K_X$ で例外）。どちらも、計算した表（\S\ref{sec:surface}）で次数 $\ge3$ が消える範囲と一致する。
\end{rem}
\begin{rem}[Pal--Poddar との関係]
Pal--Poddar \cite{PP} の Thm.~11.12(2) は正則型での類似の消滅だが、群は同じではない（葉方向の de Rham 因子の有無など）。比較は \S\ref{subsec:PPcompare}（命題 \ref{prop:PPcompare}）で行う。本定理は正則型や葉空間の軌道体を仮定せず、型が跳ぶ場合を含む代わりに、基礎の複素多様体が射影的であること（$\beta$-変換の場合に限る）を必要とする。一般の一般化複素多様体では「連接な濾過 $K^b$ の $E_1$ 項」を定義する方法がなく、これが一般化の主な障害である\tK。
\end{rem}

%--------------------------------------------------------------
\subsection{Serre 双対性}
%--------------------------------------------------------------
\begin{thm}[一般化正則直線束の Serre 双対性]\tP\label{thm:dual}
$X$ を任意のコンパクト複素多様体（$\dim_\C X=n$）、$\sigma$ を任意の正則 Poisson 構造（Kähler、射影、対数シンプレクティック、非退化のいずれも不要。$\sigma=0$ を含む）、$L$ を Poisson 加群である直線束とすると、すべての $k$ で
\[\HH^k(L)^*\;\cong\;\HH^{2n-k}\bigl(L^\vee\otimes K_X^{2}\bigr).\]
$L^\vee\otimes K_X^2$ はテンソル積の加群構造をもつ：局所枠 $e$ の双対枠 $e^\vee$ と $\omega=dz_1\wedge\cdots\wedge dz_n$ に対する枠 $\omega^2$ で $V'=-V+2V_K=-V-2M$。ここで $K_X$ の標準加群構造は $V_K^k=\sum_i\partial_i\sigma^{ki}$、$M=-V_K$（$M^j=\sum_i\partial_i\sigma^{ij}$）は $\omega$ に関する modular ベクトル場。
\end{thm}
\begin{hyptab}
$X$ コンパクト & 楕円型複体の閉値域性・有限次元性、Stokes で境界項が出ない & 手順 (4)（双対の同定）、手順 (3)\\\midrule
$[\sigma,\sigma]=0$、$L$ が Poisson 加群 & $D^2=0$、$D'^2=0$（$L'$ も加群） & 複体として扱うため\\\midrule
（不要）Kähler／非退化／対数 & 楕円性を主シンボルから直接確かめ、転置は局所的な作用素の等式なので大域的な原始関数が不要 & 手順 (1), (3)
\end{hyptab}
\begin{proof}
(0) Dolbeault モデル：$\mathrm{Tot}^k(L)=\bigoplus_{a+b=k}A^{0,a}(\wedge^bT\otimes L)$、$D=\dbar+d_K$（$d_K(\phi\otimes Pe)=(-1)^a\phi\otimes d(Pe)$、$d=[\sigma,\cdot]+V\wedge$）。$D^2=0$ は $\dbar^2=0$、$d^2=0$、$\dbar d_K+d_K\dbar=0$ による（命題 \ref{prop:ident}(b)）。

(1) 楕円性：実余接ベクトル $\xi\ne0$ での $D$ の主シンボルは $v(\xi)=(\xi^{0,1},\sigma^\sh\xi^{1,0})\in L^*=T^{*0,1}\oplus T^{1,0}$ との外積で、$(\wedge L^*,v\wedge)$ は $v\ne0$ のときに限り完全。$v=0$ なら $\xi^{0,1}=0$、$\xi$ 実なので $\xi=0$。よって $\sigma$ が退化する点も含め、すべての点で楕円型（Gualtieri \cite{Gua11} Prop.~3.13 をこの場合に直接確かめたもの）。$L'=L^\vee\otimes K^2$ についても同様。

(2) 完全対：$E^b(L)\times E^{n-b}(L')\to\wedge^nT\otimes K^2=K$ と $A^{0,a}\times A^{0,n-a}\to A^{0,n}(K)=A^{n,n}$ の積分により $\mathrm{Tot}^k(L)\times\mathrm{Tot}^{2n-k}(L')\to\C$ が各点で完全対をなす。

(3) 転置：補題 \ref{lem:E}（下）から $d$ 部分の形式的転置が $d'$ に（符号を除いて）一致する。補題 \ref{lem:E} の右辺は $\partial_iF^i\,\omega\wedge d\bar z^{1\cdots n}=\pm d(F^i\iota_{\partial_i}\omega\wedge d\bar z^{1\cdots n})$ の形の全微分（$d\bar z^{1\cdots n}$ との積で $\dbar$ 成分が消える）で、Stokes により積分は 0。$\dbar$ 部分は古典的 Serre 双対性の計算で $\langle\dbar\alpha,\beta\rangle=-(-1)^{a+b}\langle\alpha,\dbar\beta\rangle$、$d$ 部分は $\langle d_K\alpha,\beta\rangle=-(-1)^{a+b}\langle\alpha,d'_K\beta\rangle$ で、2 つの符号が一致することを手計算で確認した（機械では未確認\tQ）。$D^t$ は局所作用素なので、座標近傍での等式から大域的に $\langle D\alpha,\beta\rangle=-(-1)^{\deg\alpha}\langle\alpha,D'\beta\rangle$。

(4) 結論：$D,D'$ は閉多様体上の楕円型複体で、完全対に関して互いの転置（符号を除く）。楕円型複体の Hodge 理論により $\HH^k$ は有限次元で値域は閉、$\HH^k(L)^*$ は転置複体のコホモロジー（超関数係数）で、楕円正則性により滑らかな係数のもの $\HH^{2n-k}(L')$ に等しい。\qedhere
\end{proof}
\begin{lem}[補題 D：Schouten 括弧の最高次]\tP\ \tV\label{lem:D}
$T=\sum_j(-1)^ju_j\,\partial_1\wedge\cdots\widehat{\partial_j}\cdots\wedge\partial_n\in\wedge^{n-1}T$ に対し $[\sigma,T]_{\rm top}=-\sum_{i,j}\partial_i(\sigma^{ij}u_j)+2\sum_jM^ju_j$、$(M\wedge T)_{\rm top}=\sum_jM^ju_j$。
\end{lem}
\begin{lem}[補題 E：転置恒等式]\tP\ \tV\label{lem:E}
$P\in\wedge^bT$、$Q\in\wedge^{n-b-1}T$（滑らかな係数）、$d=[\sigma,\cdot]+V\wedge$、$d'=[\sigma,\cdot]+V'\wedge$、$V'=-V-2M$、$\langle P,Q\rangle:=(P\wedge Q)_{\rm top}$ のとき
\[\langle dP,Q\rangle+(-1)^b\langle P,d'Q\rangle=-\sum_i\partial_i\Bigl(\sum_j\sigma^{ij}u_j\Bigr),\qquad T=P\wedge Q=\sum_j(-1)^ju_j\partial_{\hat j}.\]
証明：Leibniz 則 $[\sigma,P\wedge Q]=[\sigma,P]\wedge Q+(-1)^bP\wedge[\sigma,Q]$（$\sigma$ は偶数次数）、$V$-項は $V\wedge P\wedge Q+(-1)^bP\wedge(-V)\wedge Q=0$、$M$-項は $-2M\wedge T$。補題 D から従う。
\tV\ $n=2,3,4$ で、ランダムな多項式・Laurent 単項式データ、任意の二重ベクトル場 $\sigma$（Poisson でなくても恒等式自体は成立）に対し、恒等式を厳密算術で確認した（スクリプト \texttt{duality\_ndim\_exact.py}）。定数 $c=-2$ と符号 $(-1)^b$ は、導出とは独立に連立方程式として解いて求めた（\texttt{duality\_ndim.py}）。$[\sigma,M]=0$ と $d^2=0$（$V=M$）も確認した。
\end{lem}
\begin{rem}[整合性の確認\ \tV]
(i) $\sigma=0$：古典的 Serre 双対性（$\wedge^bT^*=\wedge^{n-b}T\otimes K$）。(ii) $\mathbb{CP}^2$、$K^2=\mathcal O(-6)$：$m\leftrightarrow-m-6$；$\mathbb{CP}^1\times\mathbb{CP}^1$、$K^2=\mathcal O(-4,-4)$：$m\leftrightarrow-m-4$（\S\ref{sec:surface} の表と一致）。(iii) 対数シンプレクティックな場合、$K_X$ の標準加群構造の留数は 1 である（$\wedge^nT=K_X^{-1}$ の切断 $\sigma^k$ が平坦で $D$ に沿って 1 位に消えるため）ので、$L'$ の留数は $2-r_i$。これが曲面の主定理の公式の $r\leftrightarrow2-r$ の対称性の起源である（任意の偶数次元）。
\end{rem}

%==============================================================
\section{Euler 標数}\label{sec:euler}
%==============================================================
\subsection{全 Euler 標数は位相不変量}
\begin{thm}[Euler 標数の定理]\tP\label{thm:E}
$X$ をコンパクト複素多様体（$\dim_\C X=n$）、$\sigma$ を任意の正則 Poisson 構造、$E$ を平坦 $A$-接続をもつ階数 $\rk E$ の正則ベクトル束とする。複体 $K^\bullet=(\wedge^\bullet T\otimes E,\ d_E)$（微分は $\C$-線形）の超コホモロジー $\HH^k(E)$ は有限次元で
\[\sum_k(-1)^k\dim\HH^k(E)=\sum_b(-1)^b\chi(X,\wedge^bT\otimes E)=(-1)^n\,\rk(E)\,\chi(X).\]
\end{thm}
\begin{hyptab}
$X$ コンパクト & 各 $H^a(X,K^b)$ が有限次元になり、HRR が使える & (i), (ii)\\\midrule
$K^b$ が連接層 & Euler 標数を HRR で計算でき、下降帰納法が回る & (i)：切り捨て三角形の帰納法\\\midrule
平坦 $A$-接続 & $d_E^2=0$（複体になる） & 超コホモロジーの定義。Euler 標数の計算自体には使わない\\\midrule
（使わない）$\sigma$ の非退化性、$E$ の正値性 & --- & 証明に現れない。したがって $c_1(E)$ にも依存しない
\end{hyptab}
\begin{proof}
(i) 微分は $\C$-線形で $\mathcal O$-線形ではないが、切り捨て $K^{\ge b}$ は部分複体で、三角形 $K^{\ge b+1}\to K^{\ge b}\to K^b[-b]\to$ が成り立つ。各 $K^b$ は連接で $X$ はコンパクトなので $H^*(X;K^b)$ は有限次元。$b$ についての下降帰納法で $H^*(X;K^{\ge b})$ はすべて有限次元で、$\chi(K^{\ge b})=\chi(K^{\ge b+1})+(-1)^b\chi(X,K^b)$。したがって $\sum_k(-1)^k\dim\HH^k=\sum_b(-1)^b\chi(X,\wedge^bT\otimes E)$。
(ii) HRR より $\chi(X,\wedge^bT\otimes E)=\int\ch(E)\ch(\wedge^bT)\Td(T)$。$T$ の Chern 根を $x_i$ とすると $\sum_b(-1)^b\ch(\wedge^bT)=\prod(1-e^{x_i})$、$\Td(T)=\prod x_i/(1-e^{-x_i})$、$(1-e^x)/(1-e^{-x})=-e^x$ なので被積分形式は $(-1)^n\ch(E)e^{c_1(T)}c_n(T)$。$c_n(T)$ は既に最高次なので、$\ch(E)e^{c_1(T)}$ の次数 0 の部分 $\rk E$ だけが残り、Gauss--Bonnet で $(-1)^n\rk E\int c_n(T)=(-1)^n\rk E\,\chi(X)$。\qedhere
\end{proof}
\tV\ $\mathbb P^n$（$n\le6$）、$E=\mathcal O(m)$（$m=-12,\dots,12$）で $\sum_b(-1)^b\chi(\mathbb P^n,\wedge^bT(m))=(-1)^n(n+1)$ を Bott の公式から再計算して確認した（スクリプト \texttt{roundD\_check.py}）。

\begin{rem}[意味\ \tP]\label{rem:Emeaning}
左辺は $c_1(E)$ に依存しない。例：$\sigma=0$、$\mathbb P^1$、$E=\mathcal O(m)$ で $\chi(\mathcal O(m))-\chi(\mathcal O(m+2))=(m+1)-(m+3)=-2$。古典の $\chi(X,E)=\int\ch(E)\Td(X)$ は $b=0$ の「列」$H^\bullet(X,E)$ の Euler 標数にあたり、全複体のものではない。したがって、一般化 Dolbeault 複体の「全」Euler 標数に対する Riemann--Roch 公式は位相的な公式であり、古典の HRR の一般化にはなっていない。「$m$ についての多項式」としての Riemann--Roch は定数になり、$m$ への依存は Poincar\'e 多項式（次数ごとの次元）にのみ現れる（\S\ref{sec:surface}, \S\ref{sec:regular}, \S\ref{sec:RR}）。
\end{rem}
\begin{rem}[一般の一般化複素構造\ \tK]
一般の $\Jc$ と平坦 $L$-接続 $E$ について、Atiyah--Singer により $\mathrm{ind}(\wedge L^*\otimes E)=\int\ch(E)\cdot S$（$S$ は $(\wedge L^*,d_L)$ のシンボル類）。$L$-接続の平坦性から Chern--Weil で $\ch_j(E)$ は $H^{2j}_{dR}(M)\to H^{2j}(L)$ の核に入るはずなので、階数だけが効くと予想される。証明していない。
\end{rem}

%--------------------------------------------------------------
\subsection{対数 de Rham 複体の Euler 標数と壁越え}
%--------------------------------------------------------------
\begin{thm}[対数版の Euler 標数]\tP\label{thm:K1}
$X$ をコンパクト複素多様体、$D\subset X$ を SNC 因子、$L$ を平坦な対数接続 $\nabla:L\to L\otimes\Omega^1_X(\log D)$ をもつ直線束とする。このとき超コホモロジー $\mathbb H^k(X;\mathrm{DR}_{\log}(L))$ は有限次元で
\[\sum_k(-1)^k\dim\mathbb H^k(X;\mathrm{DR}_{\log}(L))\;=\;\int_Xc_n\bigl(T_X(-\log D)\bigr)\;=\;\chi(U)\;=\;\sum_{T}(-1)^{|T|}\chi(D_T)\]
が成り立ち、$L$ にも留数 $r_i$ にも依らない。
\end{thm}
\begin{hyptab}
$X$ コンパクト & $H^q(X;\Omega^p(\log D)\otimes L)$ が有限次元、HRR が使える & (1), (2)\\\midrule
$D$ が SNC & $\Omega^p(\log D)$ が局所自由、$\Omega^\bullet(\log D)\simeq Rj_*\C_U$（log Poincar\'e 補題）が使える & (3)（$L=\mathcal O$ の場合に位相的な値を得る）\\\midrule
$\nabla$ が平坦な対数接続 & $\mathrm{DR}_{\log}(L)$ が複体になる（階数 1 なので留数は定数） & (1)\\\midrule
（不要）留数 $r_i\notin\Z_{\ge1}$ & --- & 同定 $\mathrm{DR}_{\log}(L)\simeq Rj_*L_U$ にだけ必要で、Euler 標数には使わない\\\midrule
（不要）$\sigma$ の対数シンプレクティック性 & --- & 一般化 Dolbeault 複体との関係（定理 \ref{thm:main1}(e)）を使うときだけ
\end{hyptab}
\begin{proof}
(1) $\Omega^p(\log D)\otimes L$ に沿う濾過の $E_1$ 項は $H^q(X;\Omega^p(\log D)\otimes L)$ で有限次元。よって $\chi(\mathrm{DR}_{\log}(L))=\sum_p(-1)^p\chi(X;\Omega^p(\log D)\otimes L)$。
(2) HRR により、これは $\int_X\bigl[\sum_p(-1)^p\ch(\Omega^p(\log D))\bigr]\ch(L)\Td(T_X)$。階数 $n$ の束 $F=\Omega^1(\log D)$ について $\sum_p(-1)^p\ch(\wedge^pF)=c_n(F^\vee)\,\Td(F^\vee)^{-1}$ は次数 $n$ から始まる。したがって $\ch(L)\Td(T_X)\Td(F^\vee)^{-1}$ の定数項 $1$ だけが効き、結果は $\int_Xc_n(T_X(-\log D))$ で $L$ に依らない。
(3) $L=\mathcal O$（$\nabla=d$）のとき、log Poincar\'e 補題（Deligne \cite{Del}）$\Omega^\bullet(\log D)\simeq Rj_*\C_U$ により超コホモロジーは $H^*(U;\C)$ で、(1) の Euler 標数は $\chi(U)$。(2) よりこれは $\int c_n(T_X(-\log D))$ に等しい。最後に $\chi(U)=\sum_T(-1)^{|T|}\chi(D_T)$ は、開多様体 $U$ で $\chi_c=\chi$（Poincar\'e 双対性、実偶数次元）と $\chi_c$ の加法性から従う。\qedhere
\end{proof}
\tV\ $X=\mathbb P^2$、$D$ を次数 $d=1,\dots,7$ の滑らかな曲線、$L=\mathcal O(a)$（$a=-12,\dots,12$）で $\sum_p(-1)^p\chi(\Omega^p(\log D)\otimes L)$ を直接計算し、全 175 例で $3-3d+d^2=\chi(U)$ と一致（スクリプト \texttt{roundK\_check.py}。再実行して確認した）。

\begin{cor}[壁越えの恒等式]\tP\label{cor:wall}
$L_U$ が階数 1 の局所系のとき $\chi_c(U;L_U)=\chi(U;L_U)=\chi(U)$。これは、三角形 $j_!L_U\to Rj_*L_U\to i_*i^*Rj_*L_U$ の Euler 標数において、$D_T$（$|T|\ge1$）上の $Rj_*L_U$ の茎が $H^*((S^1)^{|T|};\cdot)$ で Euler 標数 0 である（$\chi(\partial U)=0$ と同じ理由）ことから従う。したがって、留数が整数を横切って $j_!$ と $Rj_*$ が入れ替わっても Euler 標数は変わらない。
\end{cor}
\begin{rem}[定理 \ref{thm:E} との関係]\tP
定理 \ref{thm:E} と本定理は同じ現象「Euler 標数は $\ch(L)$ の高次項を見ない」の 2 つの顔である。一般化 Dolbeault 複体の Euler 標数が $(-1)^n\chi(X)$、その部分複体 $B\simeq\mathrm{DR}_{\log}(L)$ の Euler 標数が $\chi(U)$ で、差 $(-1)^n\chi(X)-\chi(U)$ は商複体 $Q$ の Euler 標数が受け持つ（\S\ref{sec:reduction}）。
\end{rem}

%==============================================================
\section{対数シンプレクティック曲面の閉公式}\label{sec:surface}
%==============================================================
\subsection{設定と主定理}
$X$ をコンパクト複素曲面、$\sigma\in H^0(-K_X)$ は簡約な零因子 $D=\bigcup D_i$（既約成分 $D_i$、特異点は通常二重点のみ、その個数 $\mu$）をもち、$E$ は Poisson 加群をもつ直線束とする。留数 $r_i$ は定理 \ref{thm:main1}(c) による（各分枝に沿って定数。二重点の 2 つの枝は別の枝と見なし、既約な結節曲線では両枝の留数は等しい）。
\[I_*=\{i:r_i\in\Z_{\le0}\},\quad I_!=\{i:r_i\notin\Z_{\le0}\},\quad X'=X\setminus\bigcup_{i\in I_!}D_i,\quad j_1:U\hookrightarrow X',\ j_2:X'\hookrightarrow X,\]
\[\mathcal F:=j_{2!}\,Rj_{1*}L_U\]
（$I_*$ の成分では $Rj_*$、$I_!$ の成分では零拡張）とおく。

\begin{thm}[曲面の主定理]\tP\label{thm:main2}
上の設定で、すべての分枝について $r_i\ne1$ とする。このとき
\[\HH^k(E)\;\cong\;H^k(X;\mathcal F)\ \oplus\ \C^{\mu}\quad(\text{第 2 項は }k=2\text{ のみ}).\]
特に (i) すべての $r_i\notin\Z_{>0}$ なら $\HH^*(E)=H^*(U;L_U)\oplus\C^\mu[-2]$、(ii) すべての $r_i\notin\Z_{\le0}$ なら $\HH^*(E)=H^*_c(U;L_U)\oplus\C^\mu[-2]$、(iii) 全 Euler 標数は $\chi(U)+\mu=\chi(X)$（$E$ によらない、定理 \ref{thm:E} と整合）。
\end{thm}
\begin{hyptab}
曲面、$\sigma\in H^0(-K_X)$、$D$ 簡約・通常二重点のみ & 局所標準形 $h=zu$ または $z_1z_2u$ を使い、商複体 $Q$ を明示的に計算するため & 補題 \ref{lem:Q}、命題 \ref{prop:moser}、補題 \ref{lem:conn}\\\midrule
$r_i\ne1$（全分枝） & $r_i=1$ だと $Q$ の微分 $-u(1-r)p$ が 0 になり、$Q$ のコホモロジー層が無限次元になって簡約が破れる & 補題 \ref{lem:Q}（$r_i=1$ は $E=K_X$ のときの特別扱い、命題 \ref{prop:KX}）\\\midrule
留数が正整数でない／非正整数でない & 対数複体と $Rj_*L_U$, $j_!L_U$ の比較（Deligne）が擬同型になる条件 & 補題 \ref{lem:stalk}\\\midrule
$X$ コンパクト & $R\Gamma(X;j_!L_U)=R\Gamma_c(U;L_U)$ と Euler 標数の計算 & (ii), (iii)\\\midrule
$-K_X=\mathcal O(D)$ & $\chi(D)=\mu$ を導く（$\omega_D=\mathcal O_D$、$p_a(D)=1$） & (iii)
\end{hyptab}

\subsection{証明}
\begin{lem}[対数複体の茎と Deligne の比較]\tP\ \tC\label{lem:stalk}
定理 \ref{thm:main1}(d) の局所枠で、$D_I$（$|I|=s$）の点での $\mathrm{DR}_{\log}$ の茎は、可換作用素 $z_i\partial_i+r_i\ (i\in I)$ の Koszul 複体と、残りの変数の通常の de Rham 複体のテンソル積である。単項式 $z^a$ 上では $z_i\partial_i+r_i=a_i+r_i$（スカラー）なので、茎のコホモロジーは、すべての $i\in I$ で $r_i\in\Z_{\le0}$ のときに限り $\wedge^*\C^s$（次元 $2^s$、Euler 標数 0）で、そうでなければ 0。一方 $Rj_*L_U$ の同じ点での茎は、$e^{2\pi\sqrt{-1}r_i}=1$ がすべての $i\in I$ で成り立つとき（$r_i\in\Z$）に限り $\wedge^*\C^s$。ゆえに
\begin{itemize}
\item 自然な写像 $\mathrm{DR}_{\log}\to Rj_*L_U$ は、正整数の $r_i$ が一つもないとき擬同型（Deligne \cite{Del} Ch.~II \S3、番号は\tQ。ここでは直接計算で確認）。逆に、$r_i\in\Z_{>0}$ なる $i$ があれば、$D_i$ の滑らかな点で茎が 0 対 $\wedge^*\C$ となり擬同型でない。
\item 零拡張 $j_!L_U\to\mathrm{DR}_{\log}$（平坦切断の零拡張は正則）は、$r_i\in\Z_{\le0}$ となる $i$ が一つもないとき擬同型。
\item 混合の場合、$X'$ 上で $\mathrm{DR}_{\log}\simeq Rj_{1*}L_U$、$\bigcup_{I_!}D_i$ 上で茎が 0 なので $\mathrm{DR}_{\log}\simeq j_{2!}j_2^*\mathrm{DR}_{\log}\simeq\mathcal F$。
\end{itemize}
係数 $1/(a_i+r_i)$ による割り算は、有限個の $a$ を除き $|a_i+r_i|\ge c>0$ なので収束を壊さない。
\end{lem}
\begin{lem}[商複体 $Q$]\tP\ \tV\label{lem:Q}
$Q:=C(A;E)/C(T_D;E)$ は $D$ 上に台をもち、次数 1, 2 のみ。$Q^1=(T/T_D)\otimes E$、$Q^2=(\wedge^2T/\wedge^2T_D)\otimes E$。
\begin{itemize}
\item 滑らかな点（$h=zu$）：$Q^1=\partial_z\otimes\mathcal O_D$、$Q^2=\partial_z\wedge\partial_w\otimes\mathcal O_D$、$d_Q(p)=-u(1-r)\,p|_D$。よって $r\ne1$ のとき非輪状。
\item 二重点（$h=z_1z_2u$）：$Q^1=\mathcal O_{D_1}\partial_1\oplus\mathcal O_{D_2}\partial_2$、$Q^2=\mathcal O_D$、$d_Q(p,q)=-u[(1-r_1)z_2p+(1-r_2)z_1q]\bmod z_1z_2$。$r_1,r_2\ne1$ なら単射で像は極大イデアル $\mathfrak m_D$ なので $H^1(Q)_p=0$、$H^2(Q)_p=\mathcal O_D/\mathfrak m_D=\C$。
\end{itemize}
ゆえに $r_i\ne1$ のとき、導来圏で $Q\simeq\bigoplus_{\text{二重点 }p}\C_p[-2]$。
\end{lem}
\begin{proof}
式 \eqref{eq:dE} を $z=0$（二重点では $z_1=0$ または $z_2=0$）に制限する。留数の符号規約は $V=-\sigma^\sh\theta$（$\nabla^{\log}e=\theta e$）で、滑らかな点では $a|_D=\theta_2h$、$b|_D=-ru$（$\theta=r\,\dlog z+\theta_2\,dw$）、二重点では $a=r_2z_1u$、$b=-r_1z_2u$。$p\partial_z$ に対し $h(p_1+q_2)-h_1p-h_2q+qa-pb$ の $q=0$、$z=0$ での値は $-up+rup=-u(1-r)p$。二重点でも同様に計算でき、$d_Q(p,q)$ の単射性は、$z_1=0$ に制限して $(1-r_1)z_2p|_{z_1=0}=0\Rightarrow p=0$ とすることで分かる。像が $(z_1,z_2)/(z_1z_2)=\mathfrak m_D$ であることから余核は $\C$。\qedhere
\end{proof}
\tV\ ランダムな多項式 $u,g,p,q$ と留数で、滑らかな点・二重点とも各 20 回、$d_E^2=0$ と $d_Q$ の公式を厳密算術で確認した（\texttt{ncd\_cochain\_check.py}、再実行して確認）。

\begin{prop}[二重点での対数 Moser の標準形]\tP\label{prop:moser}
二重点 $p$ の近傍で $\omega=\sigma^{-1}=u^{-1}\dlog z_1\wedge\dlog z_2$、$c:=u(0)$ とする。$D$ の各枝を保ち $p$ を固定する双正則写像 $\phi$ で $\phi^*\omega=c^{-1}\dlog z_1\wedge\dlog z_2$ となるものが存在する。$c$ は局所不変量で、一般には 1 にできない。
\end{prop}
\begin{proof}
$\omega_0=\dlog z_1\wedge\dlog z_2$、$k=1/u-1/c$（$k(0)=0$）、$\omega_t=(c^{-1}+tk)\omega_0$ とおく。$\omega_t$ は閉で、$p$ の近くで各 $t\in[0,1]$ で非退化。$k=\sum k_{ab}z_1^az_2^b$ に対し
\[\alpha=\sum_{a\ge1,b\ge0}\tfrac{k_{ab}}{a}z_1^az_2^b\,\dlog z_2-\sum_{b\ge1}\tfrac{k_{0b}}{b}z_2^b\,\dlog z_1\]
は収束する対数 1-形式で $d\alpha=k\omega_0$。$X\mapsto\iota_X\omega_t$ は $T_D\to\Omega^1(\log D)$ の同型なので、$\iota_{X_t}\omega_t=-\alpha$ を満たす $X_t\in T_D$ が一意に定まり、各枝に接し $p$ で 0 になる。その流れ $\phi_t$ は $D$ を保ち、$\tfrac d{dt}\phi_t^*\omega_t=\phi_t^*(d\iota_{X_t}\omega_t+k\omega_0)=\phi_t^*(-d\alpha+k\omega_0)=0$。ゆえに $\phi_1^*\omega=\phi_1^*\omega_1=\omega_{t=0}=c^{-1}\omega_0$。\qedhere
\end{proof}
\begin{lem}[連結射の消滅]\tP\ \tV\label{lem:conn}
完全三角 $C(T_D;E)\to C(A;E)\to Q\to C(T_D;E)[1]$ の連結射の二重点 $p$ での成分 $e_p\in H^3_p(X;\mathrm{DR}_{\log})$ は、$(r_1,r_2)\ne(1,1)$ のとき 0 である。
\end{lem}
\begin{proof}[概略]
命題 \ref{prop:moser} と補題 \ref{lem:res} により $\sigma=c\,z_1z_2\partial_1\wedge\partial_2$、$\nabla^{\log}e=(r_1\dlog z_1+r_2\dlog z_2)e$ としてよい。すべての複体はトーラスウェイト $\gamma\in\Z^2$ で次数付けられ、微分はウェイトを保つ（$\lambda_i=\gamma_i+r_i$）。局所コホモロジーは Laurent 尾 $H^2_p(\mathcal O)=\mathrm{span}\{z^{-a}:a\ge1\}$ で計算できる。(a) $H^3_p(\mathrm{DR}_{\log})$ は $\lambda_1=\lambda_2=0$（$\gamma=-r$）のウェイトにのみ現れ、$r_1,r_2$ が正整数のときに限り非零。(b) $H^2_p(Q)=\C$ は単一ウェイト $\gamma=(-1,-1)$ にあり（$\mathcal O_D$ の局所コホモロジーの余分な類 $H^1_p(\mathcal O_D)$ による）、$(-1,-1)=-r$ は $r_1=r_2=1$ のときに限る。連結射はウェイトを保つので $(r_1,r_2)\ne(1,1)$ のとき $e_p=0$。すべての写像が単項式上で対角なので収束の問題はない。
\tV\ \texttt{ncd\_local\_weights.py}：ウェイト $|\gamma_i|\le8$、$r_i\in\{-3,\dots,3,\pm\tfrac12,\tfrac13,-\tfrac23,\tfrac52\}$ の 78 組（共鳴 $r_i=1$ を含む）で、茎のコホモロジー・局所コホモロジーの錐・$H(\mathrm{LOG})\to H(\mathrm{FULL})$ の単射性を厳密算術で確認した（再実行して確認）。$r_i\neq1$ では違反 0、$r=(1,1)$ では連結射は非零。上の (a), (b) のウェイトごとの手計算の詳細は紙幅のため省略し、根拠はこの計算である。\qedhere
\end{proof}
\begin{proof}[定理 \ref{thm:main2} の証明]
定理 \ref{thm:main1}(e)、補題 \ref{lem:stalk} により $C(T_D;E)\simeq\mathrm{DR}_{\log}\simeq\mathcal F$。補題 \ref{lem:Q}, \ref{lem:conn} により $C(A;E)\simeq\mathrm{DR}_{\log}\oplus\bigoplus_p\C_p[-2]$（$D^b(\C_X)$ 内）。$R\Gamma(X;\cdot)$ を取る。(i), (ii)：$\mathcal F=Rj_*L_U$ または $j_!L_U$、$X$ コンパクトより $R\Gamma(X;j_!L_U)=R\Gamma_c(U;L_U)$。(iii)：$\chi(X;\mathcal F)=\sum_S\chi_c(S)\chi(\mathcal F_x)$ で、$D$ 上の茎の Euler 標数は 0 なので $=\chi_c(U)=\chi(U)=\chi(X)-\chi(D)$。$D$ は反標準で $\omega_D=\mathcal O_D$、$\chi(\mathcal O_D)=0$ から $\sum_j(2-2\tilde g_j)=2\mu$（$\tilde D_j$ は正規化）、よって $\chi_{\rm top}(D)=\mu$、$\chi(U)+\mu=\chi(X)$。（SNC の場合の別証：$\chi(U)=\int c_2(T_X(-\log D))$、定理 \ref{thm:K1}。）\qedhere
\end{proof}

\begin{prop}[$r_i=1$：$E=K_X$]\tP\label{prop:KX}
$E=K_X$（標準加群構造、全分枝で $r_i=1$）なら $\wedge^bT\otimes K_X=\Omega^{2-b}$ で $E_1^{a,b}=H^{2-b,a}(X)$（全次数 $2+a-p$、$p=2-b$）。$h^{p,q}(X)=0\ (p\ne q)$ なら $E_1$ は全次数 2 のみに集中し、微分は全次数を 1 上げるので退化し、$\HH^k=0\ (k\ne2)$、$\dim\HH^2=\sum_ph^{p,p}=\chi(X)$（$\mathbb{CP}^2$ で 3、$\mathbb{CP}^1\times\mathbb{CP}^1$ で 4）。
\end{prop}

%--------------------------------------------------------------
\subsection{Poincar\'e 多項式の閉公式}
%--------------------------------------------------------------
定理 \ref{thm:main2} の設定で、次を追加で仮定する。(H1) $H^1(X;\C)=0$、(H2) $U$ がアフィン、(H3) 全成分で留数が等しい $r_i=r\ne1$、(H4) $r\notin\Z$ のとき $H^0(U;L_U)=H^1(U;L_U)=0$。$b_1=c-\rho_D$（$c$ は成分数、$\rho_D$ は $[D_i]\in H^2(X;\C)$ の張る空間の次元）、$\chi_U=\chi(X)-\mu$、$b_2=\chi_U-1+b_1$ とおく。

\begin{thm}[Riemann--Roch 型の閉公式（曲面）]\tP\label{thm:B}
上の仮定のもとで
\begin{align*}
P_E(t)&=1+b_1t+(b_2+\mu)t^2 &&(r\in\Z_{\le0}),\\
P_E(t)&=(b_2+\mu)t^2+b_1t^3+t^4 &&(r\in\Z_{\ge2}),\\
P_E(t)&=\chi(X)\,t^2 &&(r\notin\Z).
\end{align*}
同値な書き方：$P_E(t)=\chi(X)t^2+[r\le0]\,\Delta(t)+[r\ge2]\,t^4\Delta(1/t)$、$\Delta(t)=(1+t)(1+(b_1-1)t)$。
\end{thm}
\begin{hyptab}
(H1) $H^1(X;\C)=0$ & Gysin 完全列から $H^1(U)=\ker(\C^c\to H^2(X))$ & $b_1(U)=c-\rho_D$\\\midrule
(H2) $U$ アフィン & Artin 消滅（\cite{SGA4}） & $b_3=b_4=0$、したがって $b_2=\chi_U-1+b_1$\\\midrule
(H3) 留数が等しい & $\mathcal F$ が $Rj_*$ か $j_!$ のどちらか一方になる & 定理 \ref{thm:main2}(i), (ii)\\\midrule
(H4) $r\notin\Z$ で $H^0(U;L)=H^1(U;L)=0$ & $H^*(U;L)=[0,0,\chi_U]$ になるため & $r\notin\Z$ の行。$\mathbb{CP}^2$ の結節曲線では $\pi_1(U)$ がアーベル群（Fulton \cite{Ful}、Deligne \cite{DelFG}）\tC\\\midrule
$r_i\ne1$ & 定理 \ref{thm:main2} の仮定 & 補題 \ref{lem:Q}
\end{hyptab}
\begin{proof}
定理 \ref{thm:main2}(i), (ii) に $H^*(U)=[1,b_1,b_2,0,0]$（$r\in\Z_{\le0}$、自明な局所系）、$H_c^*(U)=[0,0,b_2,b_1,1]$（Poincar\'e 双対、$r\in\Z_{\ge2}$）、$H^*(U;L)=[0,0,\chi_U,0,0]$（$r\notin\Z$：(H2), (H4), $\chi(U;L)=\chi(U)$）を代入し $\mu t^2$ を加える。同値な形は $\chi(X)=\chi_U+\mu=1-b_1+b_2+\mu$ による。\qedhere
\end{proof}
\begin{rem}[性質\ \tP]
(i) $P_E(-1)=\chi(X)$（定理 \ref{thm:E}, $n=2$ と整合）。(ii) 双対性 $P_{2-r}(t)=t^4P_r(1/t)$ は式から明らか（定理 \ref{thm:dual} の公式レベルの帰結）。(iii) $\mu$ と $b_2$ は $\chi(X)t^2$ に吸収され、跳びの項 $\Delta$ は $b_1$ のみに依存する。(iv) $E=L_0^m$ のとき、古典の RR が与える $m$ の擬多項式は定数 $\chi(X)$ で、$m$ への依存は留数 $r=-ma$ の整数性（周期 $N$）と符号のみに現れる。
\end{rem}
\begin{rem}[(H4) について\ \tQ]\label{rem:G2}
(H4) を、検証した 8 例（$\mathbb{CP}^2$ の結節曲線 4 種、$\mathbb{CP}^1\times\mathbb{CP}^1$ の 4 種）で確かめた。$\mathbb{CP}^1\times\mathbb{CP}^1$ の滑らかな $(2,2)$ 曲線と 2 本の $(1,1)$ 曲線については Timmerscheidt \cite{Tim} と Arapura \cite{Ara}（トーション平行移動の定理）を引用した（適用範囲の確認は未了\tQ）。一般には $H^1(U;L_\chi)=0\Leftrightarrow\chi\notin V^1(U)$（第 1 特性多様体）で、非ノーダルな $D$（例：3 本共点線）では成り立たないことがある（Libgober, Arapura）。
\end{rem}

%--------------------------------------------------------------
\subsection{例と周期的安定化}
%--------------------------------------------------------------
\begin{ex}[$\mathbb{CP}^2$ と滑らかな 3 次曲線]\tP\ \tV\label{ex:cp2}
$F$ を $\mathbb C^3$ 上の滑らかな 3 次斉次多項式、$\Pi^{ij}=\epsilon_{ijk}\partial_kF$ とする。$\Pi$ は Euler ベクトル場との縮約が 0 の斉次 Poisson 構造で、$\mathbb{CP}^2$ 上の正則 Poisson 構造 $\sigma\in H^0(\wedge^2T)=H^0(\mathcal O(3))$ を定め、$D=\{F=0\}$ で 1 位に消える（標準的な構成。出典の同定は\tQ、$D$ での消え方は局所計算で確認）。
\emph{加群：}すべての $m\in\Z$ に対し $\mathcal O(m)$ はただ一つの Poisson 加群構造をもち、留数は $r=-m/3$ である。実際、定理 \ref{thm:main1}(c) より $r=-m/3$（$[D]=3H$）。存在：$\mathcal O(m)$ の Atiyah 類が $H^1(\Omega^1(\log D))$ で消えることと $c_1\in\C[D]$ が同値（$H^0(\mathcal O_D)\to H^1(\Omega^1)\to H^1(\Omega^1(\log D))$）で、$m[H]\in\C\cdot3[H]$。曲率は $H^0(\Omega^2(\log D))=\C\omega$ の元で留数が $dr=0$ だから平坦。一意性：差は $H^0(\Omega^1(\log D))=0$。
\emph{結果：}$\HH^*(\mathcal O(m))$ の次元 $[\dim\HH^0,\dots,\dim\HH^4]$ は
\begin{center}\small
\begin{tabular}{lll}\toprule
$m$ & 次元 & 解釈\\\midrule
$m\in3\Z_{\ge0}$ & $[1,0,2,0,0]$ & $H^*(U;\C)$。$\HH^0=\C F^{m/3}$ は Casimir 切断\\
$m\in3\Z_{\le-6}$ & $[0,0,2,0,1]$ & $H^*_c(U;\C)$\\
上記以外（$m=-3$ を含む） & $[0,0,3,0,0]$ & $m\notin3\Z$：$H^*(U;L_m)$。$m=-3$ は $E=K_X$（命題 \ref{prop:KX}）\\\bottomrule
\end{tabular}
\end{center}
全 Euler 標数は常に $3=\chi(\mathbb{CP}^2)$。これは定理 \ref{thm:B} から従う：$c=\rho_D=1$, $b_1=0$, $\mu=0$, $\chi_U=3$, $b_2=2$。(H4) は、$U=\{F=1\}/\mu_3$、$\{F=1\}$（Milnor ファイバー、$8$ 個の $S^2$ の束）が単連結で $\pi_1(U)=\Z/3$ だから、非自明な指標に対し $H^0=H^1=0$ で成り立つ（\tC\ Milnor \cite{Milnor}、細部は\tQ）。$r=1$ は $m=-3$ のみ。
\end{ex}
\tV\ 上の表は $m=-15,\dots,15$、Fermat 型と Hesse 型の 3 次曲線について、スペクトル系列 $E_1^{a,b}=H^a(\wedge^bT(m))$ の微分 $d_1$ を直接（$m\ge-2$ は大域切断の複体、$m\le-4$ は Laurent 単項式による \v Cech 計算）素数 $32003$ を法とする厳密な線形代数で計算した結果と一致する（スクリプト \texttt{poisson.py, table.py, neg\_check.py, predict.py}）。$m\gg0$ でも $\HH^2$ は 2 または 3 次元のまま残り、全次数での Serre 消滅は成り立たない。

\begin{ex}[$\mathbb{CP}^1\times\mathbb{CP}^1$]\tP\ \tV\label{ex:p1p1}
$\sigma=\Pi_F$、$F\in H^0(\mathcal O(2,2))$（$-K=\mathcal O(2,2)$）。$D$ が滑らかなとき、$\mathcal O(a,b)$ が Poisson 加群をもつのは $a=b=m$ のときに限り、留数は $-m/2$、構造は一意（上と同じ議論）。$m=-9,\dots,8$ の厳密計算（\texttt{p1p1.py}）：
\begin{center}\small
\begin{tabular}{lll}\toprule
$D$ & $m$ 偶数 $\ge0$ / 偶数 $\le-4$ & $m$ 奇数、$m=-2$\\\midrule
滑らかな $(2,2)$ 曲線 & $[1,0,3,0,0]$ / $[0,0,3,0,1]$ & $[0,0,4,0,0]$\\
横断する 2 つの $(1,1)$ 曲線（$\mu=2$） & $[1,1,4,0,0]$ / $[0,0,4,1,1]$ & $[0,0,4,0,0]$\\
4 本の線 $x_0x_1y_0y_1$（$\mu=4$） & $[1,2,5,0,0]$ / $[0,0,5,2,1]$ & $[0,0,4,0,0]$\\
$(1,1)+(1,0)+(0,1)$（$\mu=3$） & $[1,1,4,0,0]$ / $[0,0,4,1,1]$ & $[0,0,4,0,0]$\\\bottomrule
\end{tabular}
\end{center}
周期 2、Serre 双対性 $m\leftrightarrow-m-4$、$\chi=4$ が常に成立。可約な $D$ では Poisson 加群構造は一意でなく（$H^0(\Omega^1(\log D))\ne0$）、表は成分ごとに留数が等しい対称な構造についての計算である。
\end{ex}
\tV\ 定理 \ref{thm:main2}、\ref{thm:B} を、$\mathbb{CP}^2$ 上の滑らかな 3 次曲線、結節 3 次曲線（$\mu=1$）、円錐＋割線（$\mu=2$）、三角形（$\mu=3$）（$m=-15..15$）と上の $\mathbb{CP}^1\times\mathbb{CP}^1$ の 4 例（$m=-9..8$）、計 196 個の $m$ について厳密計算の結果と比較して食い違い 0（\texttt{ncd\_global\_check.py}）。閉公式の双対性 $r\leftrightarrow2-r$ は 73 個の有理数 $r$ の格子で確認（\texttt{serre\_rr\_check.py}、再実行して確認）。尖点 3 次、円錐＋接線、3 線共点、非簡約 $D$ は本節の対象外で、\S\ref{subsec:local} で扱う。

\begin{thm}[周期的安定化（対数シンプレクティックな曲面）]\tP\label{thm:A2}
$X$ を単連結な射影曲面、$\sigma$ を $D$（反標準、通常二重点のみ）に沿って対数シンプレクティック、$L_0$ を Poisson 加群をもつ直線束で留数が $\rho_i=-a_i$（$a_i\in\mathbb Q_{>0}$、各成分に沿って定数）、$U=X\setminus D$ をアフィン（$D$ 豊富）とする。$N$ を $a_i$ の分母の最小公倍数とする。$m\ge1$ ならば
\[\HH^k(L_0^m)=0\ (k\ge3),\qquad\HH^*(L_0^m)=H^*(U;L_U^{(m)})\oplus\C^\mu[-2].\]
ここで $L_U^{(m)}$ は $D_i$ まわりのモノドロミー $e^{-2\pi\sqrt{-1}ma_i}$ の局所系。特に $\HH^*(L_0^m)$ は $m\bmod N$ のみに依存する。
\end{thm}
\begin{hyptab}
留数 $-ma_i<0$ & $r_i\ne1$ かつ正整数でない & 定理 \ref{thm:main2}(i)：$\mathcal F=Rj_*L_U$\\\midrule
$U$ アフィン（$D$ 豊富） & Artin 消滅 $H^k(U;L)=0\ (k>2)$ & 次数 $\ge3$ の消滅（定理 \ref{thm:serre}(b) は射影性と豊富性で同じ結論を与え、アフィン性を不要にする）\\\midrule
$X$ 単連結 & $\pi_1(U)$ のアーベル化が子午線で生成される & 局所系が $m\bmod N$ のみで決まる\\\midrule
$a_i\in\mathbb Q_{>0}$ & $c_1(L_0)$ が境界成分の正の有理結合 & 周期性
\end{hyptab}
\begin{proof}
$r_i=-ma_i<0$ なので定理 \ref{thm:main2}(i) が使え、$\HH^*(L_0^m)=H^*(U;L_U^{(m)})\oplus\C^\mu[-2]$。Artin の消滅定理で $k\ge3$ の項は消える。$X$ が単連結なら $\pi_1(U)$ は子午線で生成されるので、階数 1 の局所系はモノドロミー $e^{-2\pi\sqrt{-1}ma_i}$ で決まり、これは $m\bmod N$ のみに依存する。\qedhere
\end{proof}
\tV\ $\mathbb{CP}^2$（$L_0=\mathcal O(1)$、$a=1/3$）で周期 3、$\mathbb{CP}^1\times\mathbb{CP}^1$（$L_0=\mathcal O(1,1)$、$a=1/2$）で周期 2、$m\ge1$。計算した 8 例すべてで確認（\texttt{serre\_rr\_check.py}、再実行して確認）。

%==============================================================
\section{トーリック多様体と局所項}\label{sec:toric}
%==============================================================
\subsection{トーリック対数シンプレクティック構造と weight 公式}
$X$ を滑らかな射影的トーリック多様体（$N=\Z^n$、$n=2k$、$M=\mathrm{Hom}(N,\Z)$）、$u_i\in N$ を光線の生成元、$D=\sum D_i$ をトーリック境界（SNC）とする。対数枠 $E_j=z_j\partial_j$ に関し、定数の非退化歪対称行列 $A$ により $\sigma=\tfrac12\sum A_{ij}E_i\wedge E_j$（トーラス不変で $D$ に沿って対数シンプレクティック）。$L=\mathcal O(\sum a_iD_i)$ に平坦対数接続 $d+\rho$（$\rho\in M_\C$）を与える。留数は $r_i=\langle\rho,u_i\rangle-a_i$、$c_1(L)=-\sum r_i[D_i]$。すべての加群はこの形である（$H^0(\Omega^1(\log D))=M_\C$、補題 \ref{lem:res}）\tP。$\chi\in M$ に対し
\begin{gather*}c_i(\chi)=\langle\chi,u_i\rangle+a_i,\quad C(\chi)=\{i:c_i(\chi)=-1\},\\ W(\chi)=\mathrm{span}(u_i:i\in C(\chi)),\quad s(\chi)=\dim W(\chi).\end{gather*}

\begin{thm}[weight 公式]\tP\ \tV\label{thm:weight}
$E_1^{a,b}=H^a(X,\wedge^bT\otimes L)=0\ (a>0)$ を仮定する（例：$L=L_0^m$、$m\gg0$、$L_0$ は豊富な Poisson 加群。定理 \ref{thm:serre}）。このとき
\[P_{\HH}(t)=\sum_{\chi\in\Lambda}t^{s(\chi)}(1+t)^{n-s(\chi)},\qquad\Lambda=\{\chi\in M:\ c_i(\chi)\ge-1\ \forall i,\ A(\chi+\rho)\in W(\chi)\}.\]
\end{thm}
\begin{hyptab}
トーリック、$\sigma$ トーラス不変、定数 $A$ & 大域切断の複体がトーラスの指標 $\chi\in M$ で次数付けられ、各 weight で有限次元の線形代数になる & weight $\chi$ の部分複体 $G^b_\chi$ の決定\\\midrule
$E_1^{a>0}=0$ & $\HH^*$ が大域切断の複体のコホモロジーに等しい（定理 \ref{thm:serre}(b) の議論） & 最初の等式\\\midrule
$A$ 非退化 & $\sigma^\sh:B\cong T_D$ が成立し、対数 de Rham 微分 $(\chi+\rho)\wedge$ が $v_\chi\wedge$（$v_\chi=A(\chi+\rho)$）に写る & 補題 \ref{lem:iso}、定理 \ref{thm:main1}(e)
\end{hyptab}
\begin{proof}
$z^\chi w$（$w\in\wedge^bN_\C$）が $\wedge^bT\otimes L$ の正則切断であるための条件は、各 $D_i$ に沿って $c_i(\chi)\ge-1$、かつ $c_i(\chi)=-1$ なら $u_i\wedge w=0$（$w$ が $E_i=z_i\partial_i$ で割れる）。よって weight $\chi$ の部分は $G^b_\chi=\{w\in\wedge^bN_\C:u_i\wedge w=0\ (i\in C(\chi))\}$（ある $c_i<-1$ なら空）で、$\wedge^sW$ の生成元 $\omega_W$ を用いて $G^\bullet_\chi=\omega_W\wedge\wedge^{\bullet-s}(N_\C/W)$。微分は $\wedge\sigma^\sh$ により $v_\chi\wedge$ に対応する。$v_\chi\in W$ のとき微分は 0 で、Poincar\'e 多項式 $t^s(1+t)^{n-s}$。そうでなければ $v_\chi\bmod W\neq0$ の Koszul 複体として非輪状。\qedhere
\end{proof}
\tV\ 閉公式と、線形代数による直接計算（\texttt{toric\_hh.py}）が 60 個のランダムな場合（$\mathbb P^2$, $\mathbb P^1\times\mathbb P^1$, $\mathbb P^4$, $\mathbb P^2\times\mathbb P^2$, $(\mathbb P^1)^4$）で一致した。

\begin{cor}\tP\label{cor:weight}
(a) $E_1^{a>0}=0$ なら全 Euler 標数は $s(\chi)=n$ となる weight の個数（頂点の個数）$=\chi(X)$（定理 \ref{thm:E} と整合）。
(b) 寄与の由来：$x=\chi+\rho$ とおく。$s=0$ は $\chi=-\rho$、$\rho\in M$ かつ全 $r_i\le0$ のとき $(1+t)^n=H^*(U)$（$Rj_*$ の場合）。$s=n$ は常に条件を満たし頂点 weight、$t^n$ を各 1 個。$0<s<n$ は $x\in A^{-1}W$ かつ $\langle x,u_i\rangle=r_i-1\ (i\in C)$ が必要で、積分格子上の「共鳴」でのみ解がある。
(c) $s=n-1$ で全 $r_i=1\ (i\in C)$ のとき条件は自動的に成り立ち、トーリック曲線 $C$ 上の格子点の個数 $\deg(L|_C)+1$ 個のクラス $t^{n-1}(1+t)$ が現れる（次元 2 の $r_i=1$ の共鳴に対応）。
\end{cor}

%--------------------------------------------------------------
\subsection{トーリック曲面：留数が不等な場合}
%--------------------------------------------------------------
$\Delta_\Sigma$ をファン $\Sigma$ の光線を頂点とする単体複体（$\Sigma$ の錐の光線の集合を単体とする）、$V_S$ を頂点集合 $S$ で張られる full な部分複体とし、$\tilde h^q(V_S)$ を被約コホモロジーの次元とする（$V_\emptyset=\emptyset$ では $\tilde h^{-1}=1$、他は 0 と約束）。

\begin{thm}[$\mathrm{DR}_{\log}$ のトーリックなコホモロジー]\tP\ \tC\label{thm:C1}
$X=X_\Sigma$ を完備な滑らかなトーリック多様体、$L=\mathcal O(\sum a_iD_i)$、$\nabla=d+\rho$ とし、$S=\{i:\langle\rho,u_i\rangle-a_i>0\}$ とおく。このとき
\[\sum_k\dim H^k(X;\mathrm{DR}_{\log}(L))\,t^k=\begin{cases}(1+t)^n\sum_q\tilde h^{q-1}(V_S)\,t^q&(\rho\in M),\\ 0&(\rho\notin M).\end{cases}\]
\end{thm}
\begin{hyptab}
トーラス不変な構造 & 複体が weight $\chi\in M$ で次数付けられ、各 weight で独立に計算できる & 証明の最初の分解\\\midrule
完備性 & 超コホモロジーが有限次元、\v Cech 複体が有限個のアファイン開集合で計算できる & \v Cech 複体\\\midrule
$\Omega^1(\log D)\cong\mathcal O\otimes M$ & Koszul 複体との直積分解（Künneth） & weight $\chi$ 部分が $\check C^\bullet_\chi\otimes(\wedge M_\C,(\chi+\rho)\wedge)$ の形\\\midrule
Cox--Little--Schenck の weight 公式 & $H^a(X;L)_\chi$ が単体複体の被約コホモロジーで書ける & 最後の同定
\end{hyptab}
\begin{proof}
$T$-同変なアファイン開被覆の \v Cech 複体は weight $\chi$ ごとに分解し、weight $\chi$ 部分は $\check C^\bullet(L)_\chi\otimes\wedge^\bullet M_\C$（微分 $d_{\check C}\otimes1\pm1\otimes(\chi+\rho)\wedge$）のテンソル積複体。Künneth より $H^k(X;\mathrm{DR}_{\log})_\chi=\bigoplus_{a+b=k}H^a(X;L)_\chi\otimes H^b(\wedge M_\C,(\chi+\rho)\wedge)$。Koszul 複体は $\chi+\rho\neq0$ で非輪状、$\chi+\rho=0$（$\rho\in M$ が必要）で次元 $\binom nb$ の $\wedge^bM_\C$。ゆえに $\rho\notin M$ なら全体が 0、$\rho\in M$ なら $\chi=-\rho$ のみが寄与し $(1+t)^n\sum_ah^a(X;L)_{-\rho}t^a$。最後に Cox--Little--Schenck \cite{CLS}（Thm.~9.1.3\tC\ \tQ）：$H^a(X,\mathcal O(\sum a_iD_i))_m\cong\tilde H^{a-1}(V_{D,m};\C)$、$V_{D,m}$ は $\langle m,u_i\rangle<-a_i$ の光線で張られる full な部分複体。$m=-\rho$ で条件は $\langle\rho,u_i\rangle>a_i$、すなわち $i\in S$。\qedhere
\end{proof}
\begin{thm}[トーリック曲面、留数が不等]\tP\ \tV\label{thm:C2}
$X$ を完備な滑らかなトーリック曲面、$A$ を $2\times2$ の非退化歪対称行列、すべての $r_i\ne1$ とする。このとき
\[P_{\HH}(t)=\chi(X)\,t^2+[\rho\in M]\,(1+t)^2\sum_q\tilde h^{q-1}(V_S)\,t^q,\qquad S=\{i:r_i\ge2\}.\]
特に $S=\emptyset$（すべて $r_i\le0$）で $(1+t)^2+\chi(X)t^2$、$S$ が全光線で $t^2(1+t)^2+\chi(X)t^2$。
\end{thm}
\begin{hyptab}
トーリック曲面、$\sigma$ が対数シンプレクティック & 定理 \ref{thm:main2} の適用条件：$D$ は被約で通常二重点のみ（二重点はトーラス不変点）、$\mu=\chi(X)$ & $\HH^*=H^*(X;\mathcal F)\oplus\C^\mu[-2]$\\\midrule
$r_i\ne1$（全成分） & 補題 \ref{lem:Q}（商複体 $Q$ の非輪状性）の条件 & 同上\\\midrule
$\mathrm{DR}_{\log}\simeq\mathcal F$ & 補題 \ref{lem:stalk} の混合の場合 & $H^*(X;\mathcal F)=H^*(X;\mathrm{DR}_{\log})$ に帰着し定理 \ref{thm:C1} を適用
\end{hyptab}
\begin{proof}
定理 \ref{thm:main2} により $\HH^*=H^*(X;\mathcal F)\oplus\C^\mu[-2]$、トーリック曲面では $\mu=$ 極大錐の個数 $=\chi(X)$。補題 \ref{lem:stalk} により $H^*(X;\mathcal F)=H^*(X;\mathrm{DR}_{\log})$、定理 \ref{thm:C1} で $\rho\in M$（すなわちすべての $r_i\in\Z$）のとき $(1+t)^2\sum\tilde h^{q-1}(V_S)t^q$、$S=\{r_i\ge1\}=\{r_i\ge2\}$（$r_i\ne1$）。\qedhere
\end{proof}
\tV\ $\mathbb{CP}^2$（三角形、$m=-9..6$、留数 $\{-2..3,\pm\tfrac12\}$）で 755 行、$\mathbb{CP}^1\times\mathbb{CP}^1$（4 本の直線）で 729 行について、weight 分解した \v Cech 超コホモロジーの厳密な線形代数（\texttt{roundC\_toric\_cech.py, roundC\_p2\_sweep*.py, roundC\_p1p1\_sweepA.py}）と食い違い 0。$\mathbb{CP}^1\times\mathbb{CP}^1$ で $S=$ 対辺の 2 本（$V_S$ は 2 点）のとき公式が予言する $t(1+t)^2$ の項が実際に現れる。等しい留数 $r_i=r$ では $S=\emptyset$ または全光線で、定理 \ref{thm:B} の三つの場合に一致する（例：三角形で $r\le0$：$(1+t)^2+3t^2=[1,2,4]$）。\tC\ これらの大規模な表は総合版の記録からの転記で、本論文の執筆中には再実行していない\tQ。

%--------------------------------------------------------------
\subsection{トーリック多様体（任意の偶数次元）：一般の $\sigma$}
%--------------------------------------------------------------
\begin{thm}[トーリック、一般の $\sigma$]\tP\ \tV\label{thm:T}
$X$ を滑らかな射影トーリック多様体（$\dim X=n$ は偶数）、$\sigma$ をトーラス不変な対数シンプレクティック構造（定数行列 $A$）、$L=\mathcal O(\sum a_iD_i)$、$\nabla=d+\rho$ とし、次を仮定する。
\begin{itemize}
\item[(V)] $H^a(X,\wedge^bT\otimes L)=0$（$a>0$、すべての $b$）。
\item[(A)] $L(D):=L\otimes\mathcal O(\sum D_i)$ が ample。
\item[(GEN)] $\mathrm{Pf}(A)\ne0$ かつ数 $A_{ij}\ (i<j)$ が $K=\mathbb Q(\rho_1,\dots,\rho_n)$ 上 1 次独立。
\item[(R)] すべての $i$ で $r_i\ne1$。
\end{itemize}
このとき $P_{\HH}(t)=[\ \rho\in M\ \text{かつ}\ r_i\le0\ (\forall i)\ ]\,(1+t)^n+\chi(X)\,t^n$。同値な言い方：$L$ に正の整数留数がないとき $\HH^*(L)\cong H^*(U;L_U)\oplus\C^{\chi(X)}[-n]$（$\rho\notin M$ なら $L_U$ は $(\C^*)^n$ 上の自明でない局所系で $H^*(U;L_U)=0$）。
\end{thm}
\begin{hyptab}
(V) & 超コホモロジーが大域切断の複体のコホモロジーに等しくなる & weight 公式（定理 \ref{thm:weight}）の前提。$L=L_0^m$（$m\gg0$）では定理 \ref{thm:serre} による\\\midrule
(A) $L(D)$ ample & 多面体 $\{\chi:c_i(\chi)\ge-1\}$ が単純で法ファンが $\Sigma$、頂点がちょうど $\chi(X)$ 個 & $s=n$：頂点 weight $\chi(X)$ 個、各 $t^n$\\\midrule
(GEN) & $A(\chi+\rho)\in W(\chi)$ が $K$ 上で $x\parallel m$ に退化し、算術的な偶然を排除する & $0<s<n$ の中間層の寄与を消す\\\midrule
(R) $r_i\ne1$ & $r_i=1$ だと $x=0$ または $x\parallel m_0$ の解が現れる（共鳴） & $s=n-1$ と $\chi=-\rho$ の場合に共鳴の寄与を消す\\\midrule
$\mathrm{Pf}(A)\ne0$ & $\sigma$ が対数シンプレクティック & 対数 de Rham 微分 $(\chi+\rho)\wedge$ が $v_\chi\wedge$ に対応
\end{hyptab}
\begin{proof}
定理 \ref{thm:weight} を使う。$x=\chi+\rho$、$C=C(\chi)$、$W=W(\chi)$、$s=\dim W$ とおく。$i\in C$ では $\langle x,u_i\rangle=r_i-1$。条件 $Ax\in W$ は $m\in\mathrm{Ann}_M(W)$ のすべてについて $\langle m,Ax\rangle=\sum_{i<j}A_{ij}(m_ix_j-m_jx_i)=0$ と同値。
\emph{$s=n$：}条件は空。(A) により $c_i\ge-1$ かつ $C(\chi)$ が全空間を張る $\chi$ は多面体の頂点（単純、$X$ が滑らかなので格子点）と一対一で、各極大錐に 1 個、計 $\chi(X)$ 個。各々 $t^n$。
\emph{$1\le s\le n-2$：}$\mathrm{Ann}_M(W)$ は階数 $\ge2$ なので、1 次独立な整数ベクトル $m,m'$ が取れる。係数 $m_ix_j-m_jx_i\in K$ で、(GEN) により $\sum A_{ij}(\cdots)=0$ なら全係数 0、すなわち $x\parallel m$。同様に $x\parallel m'$ で、$x=0$。$C\ne\emptyset$ のとき $x=0$ は $r_i-1=0$ を意味し (R) に反する。解なし。
\emph{$s=n-1$：}$\mathrm{Ann}_M(W)=\Z m_0$。同様に $x=\tau m_0$ で、$i\in C$ について $0=\tau\langle m_0,u_i\rangle=r_i-1$。(R) に反する。解なし。
\emph{$s=0$：}$W=0$ で条件は $x=0$、すなわち $\chi=-\rho$（$\rho\in M$ が必要）。$c_i(\chi)=-r_i\ge-1$ かつ $C=\emptyset$ なので $r_i\le0$。寄与は $(1+t)^n$。
以上を足す。後半は、$\rho\in M$ のとき $L_U$ は自明、すべての $r_i\le0$ と「正の整数留数なし」が (R) のもとで同値であることによる。\qedhere
\end{proof}
\tV\ $\mathbb{P}^2$（$m=1,2,3$）、$\mathbb{P}^1\times\mathbb{P}^1$（$\mathcal O(1,1),\mathcal O(1,2),\mathcal O(2,3)$）、$\mathbb{P}^4$（$m=1$）、$\mathbb{P}^2\times\mathbb{P}^2$（$(1,1)$）、$(\mathbb{P}^1)^4$（$(1,1,1,1)$）、$\rho$ の格子点・有理点、$A$ はランダムな整数（素数 $2^{61}-1$ を法とする）、すべての $r_i\ne1$ の 246 例で定理の式と一致（食い違い 0、\texttt{roundD\_check.py}、再実行して確認）。共鳴を含むより広い 3626 例でも weight 公式の予言と一致している（\texttt{toric\_generic\_check2.py}。再実行していない\tQ）。

\begin{obs}[(GEN) の必要性と $\sigma$ の算術への依存\ \tV]\label{obs:gen}
次元 4 では $\HH^*$ は $(X,D,L,\text{留数})$ だけでは決まらない。$\mathbb P^2\times\mathbb P^2$（$\mathcal O(1,1)$、$\rho=0$）で、$\sigma$ がブロック対角のとき $[1,4,12,16,16]$（2 つの曲面の結果の積、Künneth）、一般の $A$ では $[1,4,6,4,10]$。差 $(0,0,6,12,6)=6t^2(1+t)^2$ は、余次元 2 の 6 個の層 $\{\text{一方の因子の二重点}\}\times(\C^*)^2$ からの寄与である。仕組み（weight 公式による）：$|C|=2$ の weight で $u_i,u_j$ が $\sigma$ に関してシンプレクティックなら $x\in A^{-1}W$ が $\langle x,u_i\rangle=r_i-1$ から一意に決まり、$x\in\rho+M$ となる算術的条件のときだけ $t^2(1+t)^2$ が寄与する。ブロック構造ではこの条件が自動的に成り立つ。\tK\ $\HH$ は対数シンプレクティック構造のパラメータ空間の可算個の真解析部分集合上で上に跳び（上半連続性）、一般の値は最小である。
\end{obs}

%--------------------------------------------------------------
\subsection{共鳴 $r_i=1$}
%--------------------------------------------------------------
$r_i=1$ の成分があると定理 \ref{thm:main2} は適用できない（$d_Q=-u(1-r)p$ が 0）。次は計算に基づく。
\begin{obs}[単一・二重共鳴（$\mathbb{CP}^2$, $\mathbb{CP}^1\times\mathbb{CP}^1$）\ \tV]\label{obs:res}
$R$ を $r_i=1$ となる枝の集合、$N_R$ を共鳴枝の法束と $L$ のテンソル、$K_R=(K_X^{-1}\otimes L)|_{D_R}$、$\nu$ を共鳴枝を含まない二重点の数とするとき
\[P_{\HH}(t)=t\,h^0(N_R)+t^2\bigl(h^1(N_R)+h^0(K_R)+\nu\bigr)+t^3\,h^1(K_R).\]
$\mathbb{CP}^2$ の単一・二重共鳴と $\mathbb{CP}^1\times\mathbb{CP}^1$（非整数留数をもつ 108 例）で食い違い 0。共鳴枝が複数のときの $N_R$ の正確な定義は総合版の記録から転記したもので、確認していない\tQ。全留数が整数（$L_U$ 自明）で $m\ll0$ の場合と全分枝が共鳴する場合（$L=K_X$）は $H_c$ 型の連結写像が働き、公式は破れる（例：$\mathbb{CP}^2$ の二重共鳴 $m=-5$：観測値 $[0,0,5,2]$、公式の予言 $[0,0,6,3]$）。$L=K_X$ は命題 \ref{prop:KX} で別に扱う。\tK\ 共鳴公式は $L_U$ が自明でなく全分枝が共鳴しないときの一般の公式であろうが、証明には連結写像の計算（補題 \ref{lem:conn} の拡張）が必要で未完。トーリックでは系 \ref{cor:weight}(c) が対応する（証明済\tP）。
\end{obs}

%--------------------------------------------------------------
\subsection{非ノーダルな特異点の局所項}\label{subsec:local}
%--------------------------------------------------------------
$D$ に尖点・tacnode・3 重点などがあるとき、定理 \ref{thm:main2} は使えない。局所モデルは $\sigma=h\,\partial_z\wedge\partial_w$、加群 $\mathcal O\,h^s$（$V=s\chi_h$、$\chi_h=h_w\partial_z-h_z\partial_w$）、茎の複体 $\mathcal C_s=(\mathcal O\xrightarrow{d_0}\mathcal O^2\xrightarrow{d_1}\mathcal O)$：
\[d_0f=(hf_w+sfh_w)\partial_z-(hf_z+sfh_z)\partial_w,\qquad d_1(p,q)=h(p_z+q_w)-(1-s)(h_zp+h_wq)\]
（式 \eqref{eq:dE} で $a=sh_w$, $b=-sh_z$ としたもの）。大域では $s=r$ で、$\mathbb{CP}^2$ なら $s=-m/3$、$\mathbb{CP}^1\times\mathbb{CP}^1$ なら $-m/2$。

$h\in\C\{z,w\}$ を、正の整数の重み $(w_z,w_w)$、次数 $N_h$ に関する準斉次多項式で、原点が孤立特異点（したがって $h$ は被約）または滑らかな点とする。$\delta=N_h-w_z-w_w$、$e=w_zz\partial_z+w_ww\partial_w$（$eh=N_hh$）、$J=(h_z,h_w)$、$T=\mathcal O/J$、$\mu=\dim T$、$\mathrm{Sp}$ は多重集合 $\{\alpha_m=(\mathrm{wt}(m)+w_z+w_w)/N_h\}$（$m$ は $T$ の単項式基底）とする。

\begin{thm}[準斉次な局所公式]\tP\ \tQ\label{thm:local}
$s\ne1$ のとき、$\mathcal C_s$ のコホモロジーの Poincar\'e 多項式は
\[\mu\,t^2\;+\;[\,s\in\Z,\ s\le0\,](1+t)\;+\;(t+t^2)\cdot\#\{\alpha\in\mathrm{Sp}:\ s+\alpha\in\Z,\ s+\alpha\le1\}.\]
\end{thm}
\begin{hyptab}
準斉次 & $e\,h=N_hh$ により複体が重みで分解し、各重みが有限次元 & 手順 (2), (3)\\\midrule
孤立特異点 & $h_z,h_w$ が正則列（Hilbert 級数の公式）、$h$ が被約（Saito の基底 $e,\chi_h$） & 手順 (1), (5), (4)\\\midrule
$s\ne1$ & $s=1$ では商 $Q$ の微分が 0 で茎が無限次元 & 手順 (1)\\\midrule
正の重み & 各重み成分 $\mathcal O_k$ が有限次元多項式空間、重みごとの割り算が有界 & 手順 (2)
\end{hyptab}
\begin{proof}
\emph{(1) 分解。}$D=\{h=0\}$ は平面曲線で、$e,\chi_h\in\Theta(-\log D)$（$eh=N_hh$、$\chi_hh=0$）、$e\wedge\chi_h=-N_hh\,\partial_z\wedge\partial_w$ なので、$h$ が被約であることから Saito の判定法（\cite{Saito} Thm.~1.8\tC\ \tQ）により $\Theta(-\log D)=\mathcal Oe\oplus\mathcal O\chi_h$。基底 $(e,\chi_h)$ で書くと（$\chi_h$ の weight は $\delta$、$[e,\chi_h]=\delta\chi_h$）
\begin{align}
d_0f&=-\tfrac1{N_h}\chi_h(f)\,e+\tfrac1{N_h}\bigl(ef+sN_hf\bigr)\chi_h,\label{eq:d0}\\
d_1(a\,e+b\,\chi_h)&=\bigl[(e+sN_h-\delta)a+\chi_h(b)\bigr]\,h\,\partial_z\wedge\partial_w.\label{eq:d1}
\end{align}
（\tV\ ランダムな多項式 $f,a,b$ と有理数 $s$、7 種の準斉次 $h$ で、これらを定義式と照合した（210 検査、不一致 0、\texttt{local\_basis\_check.py}、本論文のために作成）。）したがって $B_s:=(\mathcal O\to\Theta(-\log D)\to h\mathcal O\,\partial_z\wedge\partial_w)$ は $\mathcal C_s$ の部分複体である。商 $Q=\mathcal C_s/B_s$ は、$X\mapsto X(h)\bmod h$ により $Q^1\cong J\mathcal O_D\subset\mathcal O_D$、$Q^2=\mathcal O_D$、$Q^0=0$、微分は $-(1-s)\cdot$（包含）で、$s\ne1$ なら単射。$H^2(Q)=\mathcal O_D/J\mathcal O_D=\mathcal O/J$（$h\in J$：$N_hh=w_zzh_z+w_wwh_w$）。三角 $B_s\to\mathcal C_s\to Q$ の長完全列で $H^3(B_s)=0$、$H^0(Q)=H^1(Q)=0$ だから、$H^k(B_s)=H^k(\mathcal C_s)\ (k\le1)$、$0\to H^2(B_s)\to H^2(\mathcal C_s)\to\mathcal O/J\to0$。すなわち茎 $=\mu t^2+P(B_s)$。

\emph{(2) 重みによる分解。}$\mathcal O=\prod_k\mathcal O_k$（$\mathcal O_k$ は重み $k$ の準斉次多項式）。$B_s$ の重み $k$ の部分は、$a\in\mathcal O_{k+\delta}$、$b\in\mathcal O_k$ として $B^0_k=\mathcal O_k$、$B^1_k=\{a\,e+b\,\chi_h\}$、$B^2_k=\mathcal O_{k+\delta}$、$\kappa:=k+sN_h$ とおくと式 \eqref{eq:d0}, \eqref{eq:d1} より
\[d_0f=\bigl(-\tfrac1{N_h}\chi_hf,\ \tfrac\kappa{N_h}f\bigr),\qquad d_1(a,b)=\kappa a+\chi_hb.\]
$\kappa\ne0$ のとき、ホモトピー $H^2(c)=(c/\kappa,0)$、$H^1(a,b)=N_hb/\kappa$ は $d_1H^2=\mathrm{id}$、$H^1d_0=\mathrm{id}$、$H^2d_1+d_0H^1=\mathrm{id}$ を満たす（直接計算）。$\kappa\ne0$ の重み全体では $|\kappa|\ge c>0$ なので $H$ は収束級数を保ち、$B_s$ は重み $k_0:=-sN_h$ の有限次元部分複体（$k_0\notin\Z$ なら 0）と、可縮な複体の直和に分解する。形式級数と収束級数の比較の問題は、この有界なホモトピーにより生じない。

\emph{(3) 重み $k_0$ の部分。}$\kappa=0$ なので $d_0f=(-\tfrac1{N_h}\chi_hf,0)$、$d_1(a,b)=\chi_hb$。$K=\ker(\chi_h|\mathcal O_{k_0})$、$\mathrm{Cok}=\mathrm{coker}(\chi_h:\mathcal O_{k_0}\to\mathcal O_{k_0+\delta})$ とおくと $H^0=K$、$H^1=K\oplus\mathrm{Cok}$、$H^2=\mathrm{Cok}$（$\mathcal O_j=0\ (j<0)$ とする）。

\emph{(4) $K$ の決定。}$f\in\mathcal O_{k_0}$、$\chi_hf=0$ とする。$\chi_hf=(df\wedge dh)/(dz\wedge dw)$（符号を除く）なので $df\wedge dh=0$。$e$ で縮約すると $k_0f\,dh-N_hh\,df=0$、すなわち $d(f^{N_h}/h^{k_0})=0$（$k_0>0$）。よって $f^{N_h}=c\,h^{k_0}$。$h$ は被約なので、既約因子 $p\mid h$ で $N_hv_p(f)=k_0$。$f\ne0$ なら $N_h\mid k_0$ で、$f/h^{k_0/N_h}$ は $N_h$ 乗が定数の有理関数、すなわち定数。したがって $K=\C\,h^{k_0/N_h}$（$k_0\in N_h\Z_{\ge0}$）、それ以外は $K=0$（$k_0=0$ では $f$ は定数）。

\emph{(5) $\dim\mathrm{Cok}$。}$d=k_0+\delta$、$a=w_z$、$b=w_w$。$\sum_k\dim\mathcal O_kx^k=1/((1-x^a)(1-x^b))$、$h_z,h_w$ は次数 $N_h-a$, $N_h-b$ の正則列なので $\sum\dim T_kx^k=(1-x^{N_h-a})(1-x^{N_h-b})/((1-x^a)(1-x^b))$。恒等式 $x^\delta+x^{N_h}-x^{N_h-a}-x^{N_h-b}=x^\delta(1-x^a)(1-x^b)$ から
\[\frac{(1-x^{N_h-a})(1-x^{N_h-b})}{(1-x^a)(1-x^b)(1-x^{N_h})}-\frac{1-x^\delta}{(1-x^a)(1-x^b)}=\frac{x^\delta}{1-x^{N_h}}.\]
$x^d$ の係数を取ると $\sum_{r\ge0}\dim T_{d-rN_h}-(\dim\mathcal O_d-\dim\mathcal O_{k_0})=[\,k_0\in N_h\Z_{\ge0}\,]=\dim K$。$\dim\mathrm{Cok}=\dim\mathcal O_d-\dim\mathcal O_{k_0}+\dim K$ なので $\dim\mathrm{Cok}=\sum_{r\ge0}\dim T_{d-rN_h}$。

\emph{(6) スペクトルへの翻訳。}$s=-k_0/N_h$。$T_{d-rN_h}$ の基底単項式 $m$（$\mathrm{wt}(m)=d-rN_h$）に対し $\alpha_m=(\mathrm{wt}(m)+w_z+w_w)/N_h=1-r-s$、すなわち $s+\alpha_m=1-r\in\Z_{\le1}$。逆に $s+\alpha\in\Z_{\le1}$ をみたす $\alpha$ はちょうど 1 つの $r\ge0$ に対応する。よって $\dim\mathrm{Cok}=\#\{\alpha\in\mathrm{Sp}:s+\alpha\in\Z,\ s+\alpha\le1\}$、$\dim K=[\,s\in\Z_{\le0}\,]$。(3) と (1) から定理を得る。\qedhere
\end{proof}
\begin{rem}[状態と残る未確認事項]
\tV\ 手順 (3)--(6) の次元（$K$, $\mathrm{Cok}$ と上の公式）を、14 種（$A_1,A_2,A_3,A_6,D_4,D_5,E_6,E_8$, $z^3+w^6$, $z^4+w^5$, $z^5+w^5$, $z^4+w^4$, $z^3+w^7$, $z^5+w^6$）約 700 個の $k_0$ について有理数の厳密計算で照合し、不一致 0（\texttt{roundH\_proof\_check.py}、再実行して確認）。\tQ\ 残る未確認：(a) 局所モデル $\mathcal C_s$（$V=s\chi_h$）が大域的な Poisson 加群の茎として現れること（準斉次でない germ を含め、留数 $s$ の符号規約は式 \eqref{eq:dE} の $V=-\sigma^\sh\theta$ に合わせて固定した）、(b) Saito の判定法の番号。適用範囲は準斉次の孤立特異点（または滑らかな点）と $s\ne1$ のみで、準斉次でない孤立特異点・非孤立特異点・$s=1$ は\tK（未証明）。
\end{rem}
\begin{rem}[解釈\ \tC\ \tK]
Milnor ファイバー束 $h:B\setminus D\to\C^*$ の Wang 完全列（ファイバーは $\mu$ 個の円の束で、$H^1$ のモノドロミー固有値は $e^{-2\pi\sqrt{-1}\alpha}$）から $(Rj_*L_s)_p=[s\in\Z](1+t)+(t+t^2)\#\{\alpha:s+\alpha\in\Z\}$ となる。定理の公式はこれに切り捨て $s+\alpha\le1$（部分複体 $B_s$ が重み $\ge-\delta$ の部分を計算することに由来）を課したものである。総合版（Part II v1）の規則 $[s<\mathrm{lct}]\cdot(\cdots)$ は $\mathrm{lct}<1/2$ のとき誤り（例：$z^4-w^5$、$\mathrm{lct}=9/20$ で $s=11/20=1-\mathrm{lct}$ には $s>\mathrm{lct}$ だが類が存在する）で、$\mathrm{lct}\ge1/2$ では一致する。Calder\'on-Moreno--Narv\'aez-Macarro \cite{CMNM} は自由因子について類似の対数比較定理を与えている（頁は\tQ）。
\end{rem}
\begin{obs}[代表的な型の茎の次元\ \tP]\label{obs:localtable}
定理 \ref{thm:local} から $(\dim H^0,\dim H^1,\dim H^2)$ は次のようになる（$\mu$ は Milnor 数、「$\alpha$」の列は $s+\alpha\in\Z_{\le1}$ をみたす $\alpha\in\mathrm{Sp}$ がちょうど 1 つあり $\alpha\notin\Z$ のとき）：
\begin{center}\small
\begin{tabular}{llcccc}\toprule
特異点 & $\mathrm{Sp}$ & $\mu$ & 一般の $s$ & $s\in\Z_{\le0}$ & $\alpha$\\\midrule
滑らか $(z)$ & $\emptyset$ & 0 & $(0,0,0)$ & $(1,1,0)$ & ---\\
二重点 $(zw)$ & $\{1\}$ & 1 & $(0,0,1)$ & $(1,2,2)$ & ---\\
尖点 $(z^2-w^3)$ & $\{\frac56,\frac76\}$ & 2 & $(0,0,2)$ & $(1,1,2)$ & $(0,1,3)$\\
tacnode $(z^2-w^4)$ & $\{\frac34,1,\frac54\}$ & 3 & $(0,0,3)$ & $(1,2,4)$ & $(0,1,4)$\\
$A_4$ $(z^2-w^5)$ & $\{\frac7{10},\frac9{10},\frac{11}{10},\frac{13}{10}\}$ & 4 & $(0,0,4)$ & $(1,1,4)$ & $(0,1,5)$\\
$D_4$ $(zw(z-w))$ & $\{\frac23,1,1,\frac43\}$ & 4 & $(0,0,4)$ & $(1,3,6)$ & $(0,1,5)$\\
$D_5$ $(w(z^2-w^3))$ & $\{\frac58,\frac78,1,\frac98,\frac{11}8\}$ & 5 & $(0,0,5)$ & $(1,2,6)$ & $(0,1,6)$\\\bottomrule
\end{tabular}
\end{center}
総合版の表（Part I §I.6.6）の値のうち、$A_4$ の最後の列は総合版では未記入であった。
\end{obs}
\noindent \tP\ 訂正：$D_5=z^2w-w^4$ の重みは $(3,2)$、次数 8、$\mathrm{Sp}=\{5/8,7/8,1,9/8,11/8\}$、$\mathrm{lct}=5/8$、$\mu=5$。内部メモにあった重み $(3,1,7)$ は誤りで、そこから出た値は無効（総合版の表と本文は正しい重みで書かれている）。

\begin{obs}[留数が不等な三重点\ \tV]\label{obs:triple}
$h=zw(z-w)$（$D_4$）、加群 $\prod h_i^{s_i}$（各枝に別の $s_i\ne1$）の茎は $\C^4[-2]\oplus R$、$R=H^*(U_{\rm loc};L_s)$（$s_1+s_2+s_3<2$）、$R=0$（それ以外）。$U_{\rm loc}=\C^2\setminus\{3\text{ 本の直線}\}\simeq\C^*\times(\mathbb P^1\setminus3\text{ 点})$。具体的には、すべての $s_i\in\Z_{\le0}$ なら茎 $(1,3,6)$、$\sum s_i\in\Z_{\le1}$ で $s_i$ が全部整数ではないとき $(0,1,5)$、それ以外 $(0,0,4)$。$s_i\in\tfrac16\Z\cap[-2,2]\setminus\{1\}$ のランダムな 533 組で食い違い 0（\texttt{local\_unequal.py, local\_unequal\_rule.py}。再実行していない\tQ）。節点では枝ごとの閾値 $s_i<1$ が $Rj_*$ 部分の有無を決めるのに対し、三重点では留数の\emph{和}が閾値 $2=w_z+w_w$ を超えるかで決まる。等しい留数 $s_1=s_2=s_3=s$ では、定理 \ref{thm:local} の $D_4$ の値と $s\in\tfrac13\Z$, $-10\le3s\le10$ ($s\ne1$) で一致する（照合\tV）。\tK\ 一般原理：閾値は germ の対称トーラスの下での加群の weight と log canonical weight の比較で決まる。
\end{obs}
\begin{conj}[精密化 Riemann--Roch（等留数）]\tK\label{conj:refined}
$X$ をコンパクト曲面、$\sigma\in H^0(-K_X)$ の零曲線 $D$ は簡約で、すべての特異点が孤立・準斉次かつ $\mathrm{lct}_p\ge1/2$（$A_k,D_k,E_{6,7,8}$ 型など）とする。$L$ は $c_1(L)=sK_X$ の Poisson 加群、$\mathcal F_s$ は $U$ 上で $L_s$、$p\in D$ で $s<\mathrm{lct}_p$ なら $Rj_*L_s$ の茎、$s>\mathrm{lct}_p$ なら 0 とする構成可能複体とする。(a) $s\notin\mathrm{Sp}(D)\cup\{1\}$ のとき $\HH^*=H^*(X;\mathcal F_s)\oplus\bigoplus_p\C^{\mu_p}[-2]$。(b) Euler 標数は常に $\chi(X)$（$\chi(D)=\sum\mu_p$ と $\delta$-不変量から\tP）。(c) $s\ll0$ では $\mathcal F_s=Rj_*L_s$ で $\HH^*=H^*(U;L_s)\oplus\bigoplus\C^{\mu_p}[-2]$ は $s\bmod\Z$ のみに依存し次数 $\ge3$ で消える。(d) $s\in\mathrm{Sp}(D)$ では階数 1 のキャンセルが起こりうる。
\end{conj}
\tV\ $\mathbb{CP}^2$（滑らか・結節・尖点・円錐＋割線・円錐＋接線・三角形・3 重点）と $\mathbb{CP}^1\times\mathbb{CP}^1$（4 例）で $m=-15..15$（$-9..8$）に対し (d) 以外で矛盾なし。3 重点 $xy(x-y)$（$\mathbb{CP}^2$）：$m=1,2,4,5,\dots$ と $m=-1$ は $[0,1,4,0,0]$、$m=-5,-7,-8$ は $[0,0,4,1,0]$、$m=-2,-3,-4$ は $[0,0,3,0,0]$。後者の異常（$m=-2,-4$ で $\C^4$ の次数 2 の 1 クラスと $H_c$ の次数 3 の 1 クラスが相殺される）はスペクトル系列の高次微分によると推測するが、証明も定理 \ref{thm:local} による説明もない\tK。証明されていない：大域同定は当てはめで、茎の結果は準斉次 germ のみ。
\begin{obs}[非被約な境界\ \tV]\label{obs:nr}
$\sigma=\Pi_F$ で $F$ が重複因子をもつとき、$\HH^*$ の次元は $m$ について線形に増える（厳密な線形代数、$m=-15,\dots,15$、\texttt{roundF\_nonreduced.py}。再実行していない\tQ）。(a) $F=x^2y$（二重線と線、結節点 1 個、$U\simeq S^1$）：$m\ge-2$：$P(t)=[3\mid m](1+t)+3t^2+(m+2)(t+t^2)$、$m=-3$：$3t^2$、$m\le-4$：Serre 双対 $m\leftrightarrow-m-6$ の像。(b) $F=x^3$（三重線、$U=\C^2$）：$m\ge0$ で $[1,2m+5,2m+7]$、$m=-1$ で $[0,2,5]$、$m=-2,-3,-4$ で $[0,0,3]$。(c) $\mathbb P^1\times\mathbb P^1$、$D=2C$（$C$ は滑らかな $(1,1)$ 曲線）、$L=\mathcal O(m,m)$：$m\ge0$ で $H^0=1$, $H^1=2m+3$, $H^2=2m+6$。
\end{obs}
\begin{conj}[非被約な境界]\tK\label{conj:NR}
$X=\mathbb P^2$ または $\mathbb P^1\times\mathbb P^1$、$D=\sum k_iD_i$（$D_i$ 滑らか、$D_{\mathrm{red}}$ は SNC）、$m\ge0$（または高次の層コホモロジーが消える範囲）で
\[\HH^*(L)=H^*(U;L_U)\oplus\C^{\chi(D_{\mathrm{red}})}[-2]\oplus\bigoplus_i\bigoplus_{j=1}^{k_i-1}\bigl(H^0(D_i,L|_{D_i}\otimes N_{D_i}^{j})[-1]\oplus H^0(D_i,L|_{D_i}\otimes N_{D_i}^{j})[-2]\bigr).\]
\end{conj}
\noindent 確認は観察 \ref{obs:nr} の 3 つの族のみ。全 Euler 標数は $\chi(U)+\chi(D_{\mathrm{red}})=\chi(X)$ で定理 \ref{thm:E} と整合。ヒューリスティクス：$z^k\partial_z\wedge\partial_w$ の Poisson コホモロジーが無限次元の部分をもつ（Monnier, Roger--Vanhaecke 型の結果\tC\ \tQ、未確認）。$m$ が小さいと破れる（$x^3$ の $m=-1$ は合わない）ので、成り立つ範囲の正確な境界は未知。

%==============================================================
\section{一般次元：簡約定理と定理 V}\label{sec:reduction}
%==============================================================
\subsection{設定、商複体、条件 (GEN$_Q$)}
$X$ をコンパクト複素多様体（$\dim X=n$ は偶数）、$D=\sum D_i$ を SNC 因子、$\sigma$ を対数シンプレクティック、$L$ を Poisson 加群（$D$ に沿った平坦な対数接続、留数 $r_i$）とする（定義 \ref{defn:logsymp}、定理 \ref{thm:main1}）。次の複体の層を使う。
\[C=(\wedge^\bullet T\otimes L,\ d_L=[\sigma,\cdot]+V\wedge),\quad B=\wedge^\bullet T_D\otimes L\subset C,\quad Q=C/B.\]
$\HH^k=H^k(X;C)$。$\wedge\sigma^\sh$ により $B\cong\mathrm{DR}_{\log}(L)$（補題 \ref{lem:iso}）\tC。三角 $B\to C\to Q\xrightarrow{\ \partial\ }B[1]$ がある。

\paragraph{$Q$ の構造。}対数座標 $(z_1,\dots,z_k,w)$ では $\wedge T$ の $\mathcal O$-基底は $\partial_S\wedge\partial_w^{Y}$（$S\subset\{1..k\}$）、$\wedge T_D$ の基底は $E_S\wedge\partial_w^Y=z^S\partial_S\wedge\partial_w^Y$。したがって $Q$ は、単項式 $z^a\partial_S\wedge\partial_w^Y$（$a_i=0$ となる $i\in S$ の集合 $T(s)\neq\emptyset$）で張られ、これは「$T(s)$ の各分枝に極をもつ多重ベクトル場」$z^{a-S}E_S\wedge\cdots$ と読める。$\sigma\in\wedge^2T_D$、$V\in T_D$ なので $d_L$ は新しい極を作らない（$[X,z_i^{-1}]\in z_i^{-1}\mathcal O$（$X\in T_D$））。よって $F_kQ=\{|T(s)|\le k\}$ は $Q$ の部分複体の増大列で、$\mathrm{gr}_0Q=0$。最上位 $F_nQ=Q$ では、$|T(s)|=n$ は $n$ 重点 $p$ での単項式 $\partial_1\wedge\cdots\wedge\partial_n$ だけである。次数 $n$ の部分では $Q^n=(\wedge^nT/\wedge^nT_D)\otimes L=K_X^{-1}\otimes L|_D$（$\wedge^nT=K_X^{-1}$、$\wedge^nT_D=K_X^{-1}(-D)$）である。この観察から、$Q$ から点への次の写像が定義される。
\[\varepsilon^n:\ Q^n=K_X^{-1}\otimes L|_D\longrightarrow\bigoplus_{p\in D_{(n)}}(K_X^{-1}\otimes L)_p\otimes\C,\quad(\text{n 重点での値}),\qquad \varepsilon^k=0\ (k\ne n).\]

\begin{prop}[\texttt{gr}\,$Q$ と $\varepsilon$\ \tQ]\label{prop:grQ}
(a) $\varepsilon:Q\to\bigoplus_{p\in D_{(n)}}\C_p[-n]$ は複体の準同型である（すなわち $d_L(Q^{n-1})$ は各 $n$ 重点で値 0 をとる）。
(b) $\mathrm{gr}_kQ=\bigoplus_{|T|=k}Q_T$、$Q_T\cong\bigl(\wedge^\bullet T_{D_T}(-\log D'_T)\otimes L|_{D_T}\bigr)\cdot z^{-T}E_T$（$D'_T$ は $D_T$ に他の成分が切り取る因子）、微分は $d_T=[\sigma_T,\cdot]+\nu_T\wedge$。ここで $\sigma_T$ は $D_T$ 上に誘導される Poisson 構造、$\nu_T=\sum_{i\in T}(r_i-1)H_i$、$H_i=\sigma^\sh(\dlog z_i)$。
\end{prop}
\begin{proof}[状態と根拠]
$n=2$ は補題 \ref{lem:Q}（$d_Q$ の明示的な公式）：滑らかな点では $d_Q=-u(1-r)\cdot$、二重点では像が $\mathfrak m_D$ なので (a) は成り立ち、(b) も定義から従う\tP\ \tV。一般 $n$ では、(b) の「$-1$」は $[\sigma,z_i^{-1}]=-z_i^{-1}H_i$ から、「$r_i$」はモジュール構造 $V=\sigma^\sh(\sum r_i\dlog z_i)+\cdots$ から来て、$[\sigma_T,\nu_T]=0$ は $\sigma$ の Jacobi 恒等式と $\nu_T$ が Poisson ベクトル場の和であることから従う。一行ずつの Schouten 括弧の計算は書き下しておらず、(a), (b) を一般の $n$ で検査した計算もない。本論文はこの命題を、定理 \ref{thm:R} の\emph{仮定} (GEN$_Q$) を局所的な条件に言い換えるためにだけ使い、定理 \ref{thm:R} の証明には使わない。\tQ
\end{proof}

\begin{defn}[条件 (GEN$_Q$)。定理 \ref{thm:T} の (GEN) とは別]\label{defn:GEN}
$\varepsilon:Q\to\bigoplus_{p\in D_{(n)}}\C_p[-n]$ が複体の準同型であり、かつ擬同型であること。
\end{defn}
\begin{defn}[条件 (GEN$'$)]
$1\le|T|\le n-1$ のすべての層 $T$ について、$Q_T$ が茎ごとに非輪状であること。命題 \ref{prop:grQ} が成り立つなら、(GEN$'$) $\Rightarrow$ (GEN$_Q$) である（$F_{n-1}Q$ の次数商がすべて非輪状なので $F_{n-1}Q$ が非輪状、$Q/F_{n-1}Q=\mathrm{gr}_nQ\cong\bigoplus\C_p[-n]$）\tQ。$n=2$ では (GEN$_Q$) $\iff$ すべての $r_i\ne1$（補題 \ref{lem:Q}）\tP。
\end{defn}

\begin{prop}[(GEN$'$) の必要条件と例]\label{prop:genNec}
(a) $\nu_T\equiv0$ なら $Q_T$ は $\sigma_T$ の Lichnerowicz 複体で、局所定数の Casimir が $H^0$ に入る。特に $|T|=1$ で $r_i\ne1$ が必要。\tP\ \tQ（命題 \ref{prop:grQ}(b) に依存）
(b) $|T|=1$、$D$ の滑らかな点では、対数 Darboux 座標で $\sigma=E_z\wedge\partial_{w_1}+\sum_j\partial_{w_{2j}}\wedge\partial_{w_{2j+1}}$ と書けるとき $\nu=(r-1)\partial_{w_1}$、$Q$ はシンプレクティック部分の Lichnerowicz 複体への $(r-1)\cdot\mathrm{id}$ の写像錐で、非輪状 $\iff r\ne1$。\tC\ \tQ\ 対数 Darboux は実の $b$-シンプレクティックでは Guillemin--Miranda--Pires, Adv.\ Math.\ 264 (2014) が知られているが、正則版は同じ Moser の議論（命題 \ref{prop:moser} の対数版）によると思われるものの、確認していない。
\end{prop}

\subsection{簡約定理}
\begin{thm}[簡約定理]\tP\label{thm:R}
上の設定で、(i) どの $r_i$ も $\Z_{\ge1}$ に入らず、(ii) (GEN$_Q$)（定義 \ref{defn:GEN}）が成り立つとする。このとき
\[\HH^*(L)\ \cong\ H^*(U;L_U)\ \oplus\ \bigoplus_{p\in D_{(n)}}\C[-n].\]
特に $\chi(X)=\chi(U)+\#D_{(n)}$（$n$ は偶数）。
\end{thm}
\begin{hyptab}
(i) $r_i\notin\Z_{\ge1}$ & 共鳴を避ける。$\mathrm{DR}_{\log}(L)\simeq Rj_*L_U$（Deligne）が成り立つための条件 & 手順 (1)\\\midrule
(ii) (GEN$_Q$) & $Q$ が $n$ 重点だけの寄与になる。成り立たなければ余分なクラスが出る（観察 \ref{obs:gen}） & 手順 (2)\\\midrule
$D$ が SNC、$\sigma$ が対数シンプレクティック & $B\cong\mathrm{DR}_{\log}(L)$ & 手順 (1)\\\midrule
$X$ コンパクト & $H^*(X;\cdot)$ が有限次元、定理 \ref{thm:E} が使える & 手順 (4)\\\midrule
$n$ 偶数 & 符号 $(-1)^n=1$。対数シンプレクティックなら自動 & 手順 (4)\\\midrule
（不要）射影性、$L$ の豊富性、消滅 (V) & --- & 定理 \ref{thm:T} の weight 計算でのみ使う
\end{hyptab}
\begin{proof}
(1) (i) により Deligne \cite{Del}（II.3.13；Esnault--Viehweg \cite{EV}）\tC\ から $B\simeq\mathrm{DR}_{\log}(L)\simeq Rj_*L_U$。
(2) (ii) により $\varepsilon:Q\xrightarrow{\sim}\bigoplus_p\C_p[-n]$。したがって三角 $B\to C\to Q\xrightarrow\partial B[1]$ の連結射は $\partial\in\mathrm{Hom}_{D(X)}(Q,Rj_*L_U[1])=\mathrm{Hom}(\bigoplus_pi_{p*}\C[-n],Rj_*L_U[1])$。随伴 $j^*\dashv Rj_*$ より $=\mathrm{Hom}(j^*\bigoplus i_{p*}\C[-n],L_U[1])=0$（$p\in D$ なので $j^*i_{p*}=0$）。
(3) よって三角は分裂し、$C\simeq B\oplus Q\simeq Rj_*L_U\oplus\bigoplus_p\C_p[-n]$。$H^*(X;\cdot)$ を取り、$H^*(X;Rj_*L_U)=H^*(U;L_U)$（定義）を使う。
(4) Euler 標数：定理 \ref{thm:E} より $\chi(\HH)=(-1)^n\chi(X)=\chi(X)$、右辺は $\chi(U)+(-1)^n\#D_{(n)}=\chi(U)+\#D_{(n)}$（$L_U$ は階数 1）。\qedhere
\end{proof}
\begin{rem}[旧版との違い\ \tP]
総合版の簡約定理の証明は $Q$ の濾過と次数商 (GEN$'$) を使い、一般の $n$ での次数商の導出が見取り図であった。上の証明は条件を (GEN$_Q$)（$\varepsilon$ が擬同型）の形で直接仮定するので、命題 \ref{prop:grQ} に依存しない。仮定の言い換え (GEN$'$) $\Rightarrow$ (GEN$_Q$) の部分が\tQ\ として残る。
\end{rem}
\tV\ $n=2$ では定理 \ref{thm:main2} と整合し、トーリック $n=2,4$ の 61 例（\texttt{roundK\_check.py}）で簡約定理と閉公式（定理 \ref{thm:K}）の予言が実測と一致（食い違い 0）。トーリックの場合に (GEN$_Q$) が成り立つ条件は、重み $\chi$ の部分の $Q_T$ が $\wedge(N_\C/W_T)$ の Koszul 複体（微分 $v_\chi\wedge$）になることから、$v_\chi\notin W_T$ で非輪状、すなわち定理 \ref{thm:T} の仮定 (GEN)（$A_{ij}$ の $\mathbb Q(\rho)$ 上の独立性）と (R) のもとでの定理 \ref{thm:T} の証明の $x=0$ の排除と同じ議論で成り立つ\tP\ \tQ（命題 \ref{prop:grQ} に依存）。定理 \ref{thm:T} 自身は weight 公式から直接証明されており、この議論には依存しない。

\begin{conj}[(GEN$'$) が成り立たない場合]\tK\label{conj:nonacyc}
$E_1^T=H^*(D_T;Q_T)\Rightarrow H^*(X;Q)$。$\nu_T=0$ の $Q_T$ がすべて「積型」なら、余分な部分は $\sum_Tt^{|T|}H^*(U_T;L_T)$。観察 \ref{obs:gen} の $6t^2(1+t)^2$ は、$T=(\text{一方の因子の二重点})\times Y$ の 6 個の層（$\nu_T=0$）が $Q_T\simeq Rj_*\C_{U_Y}$ を通して $t^2(1+t)^2$ を寄与する、と読める（\tV\ 数値の一致のみ）。
\end{conj}

%--------------------------------------------------------------
\subsection{定理 V：$H^k(U;L_U)$ の消滅}
%--------------------------------------------------------------
簡約定理は $H^*(U;L_U)$ を未知量として残す。留数がすべて非整数ならこれは次のように決まる。

\begin{thm}[定理 V]\tP\label{thm:V}
$X$ をコンパクト複素多様体（$\dim_\C X=n\ge1$）、$D=\sum_{i=1}^sD_i$ を被約な SNC 因子、$U=X\setminus D$ は\emph{アフィン}とする。$L_U$ を $U$ 上の階数 1 の局所系で、$D_i$ のまわりの単純モノドロミーが $e^{-2\pi\sqrt{-1}\,r_i}$、かつ\emph{すべての $r_i\notin\Z$} とする。このとき $H^k(U;L_U)=0\ (k\ne n)$、したがって $\dim H^n(U;L_U)=(-1)^n\chi(U)$。
\end{thm}
\begin{hyptab}
$r_i\notin\Z$（すべての $i$） & 各ストラタムでの茎が消える & (a)。ある層ですべての $r_i\in\Z$ なら茎は $H^*(T^k)\ne0$ で、これが第 5, 6 節の跳びの源\\\midrule
$U$ アフィン & Artin 消滅 & (c)。$D=\emptyset$ など非アフィンでは不成立\\\midrule
SNC & 局所基本群が $\Z^k$ & (a)。非 SNC（三重点など）では局所基本群が $\Z^k$ でなく、茎の規則に $\sum r_i$ が現れる（観察 \ref{obs:triple}）\\\midrule
$X$ コンパクト & $Rj_*$ の大域切断が $U$ のコホモロジーを与え、$j_!=Rj_*$ なら $H^k=H^k_c$ & (b)
\end{hyptab}
\begin{proof}
(a) $p\in D_I$（$|I|=k\ge1$）での $Rj_*L_U$ の茎は、小多重円板 $B$ に対する $H^*(U\cap B;L_U)$。$U\cap B\simeq(\Delta^*)^k\times\Delta^{n-k}\simeq(S^1)^k$ で、$L_U$ は各円周上のモノドロミー $\mu_i=e^{-2\pi\sqrt{-1}r_i}\ne1$ の局所系の外部テンソル積。Künneth により $H^*((S^1)^k;L)=\bigotimes_{i\in I}H^*(S^1;L_{\mu_i})$ で、$\mu\ne1$ なら $H^0=\ker(\mu-1)=0$、$H^1=\mathrm{coker}(\mu-1)=0$。ゆえに $D$ 上で茎が消え、$Rj_*L_U=j_!L_U$。
(b) $X$ がコンパクトなので $H^k(U;L_U)=H^k(X;Rj_*L_U)=H^k(X;j_!L_U)=H^k_c(U;L_U)$。
(c) $U$ がアフィンなので Artin の消滅定理（Andreotti--Frankel）\tC\ \cite{SGA4} により、任意の局所系 $F$ で $H^q(U;F)=0\ (q>n)$。よって $k>n$ では $H^k(U;L_U)=0$。向き付けられた $2n$ 次元多様体 $U$ での Poincar\'e 双対 $H^k_c(U;L_U)\cong H^{2n-k}(U;L_U^\vee)^*$ と (b) より、$2n-k>n$、すなわち $k<n$ で $H^k(U;L_U)=0$。最後に $\chi(U)=(-1)^n\dim H^n$。\qedhere
\end{proof}
\begin{cor}\tP\label{cor:Vgen}
定理 \ref{thm:R} の仮定を満たし、すべての $r_i\notin\Z$、$U$ アフィンなら $P_{\HH}(t)=\chi(X)\,t^n$。
\end{cor}
\begin{proof}
定理 \ref{thm:R} と定理 \ref{thm:V} より $P=\dim H^n(U;L_U)t^n+\#D_{(n)}t^n=(\chi(U)+\#D_{(n)})t^n=\chi(X)t^n$（$n$ 偶数）。\qedhere
\end{proof}
一部の $r_i$ だけが整数の「混合」の場合：$D_I$ 上の茎が残るのは全 $r_i\ (i\in I)$ が整数のときだけだが、$k<n$ での消滅は未証明\tK。
\paragraph{計算による確認\tV。}トーリックで $n=4$ の厳密計算（$2^{61}-1$ を法とする線形代数）：$\mathbb P^4$（$\mathcal O(1)$、$r=(1/2,1/3,1/5,1/7,-457/210)$）が $(0,0,0,0,5)$、$\mathbb P^4$（$\mathcal O(2)$、$r=(1/3,1/3,1/3,1/3,-10/3)$）が $(0,0,0,0,5)$、$\mathbb P^2\times\mathbb P^2$ が $(0,0,0,0,9)$、$(\mathbb P^1)^4$ が $(0,0,0,0,16)$。いずれも $\chi(X)t^4$ で系 \ref{cor:Vgen} の予測に一致する。混合の例（$\mathbb P^1\times\mathbb P^3$ で整数留数 $-2$ を 1 つ含む）も $(0,0,0,0,8)$ となったが、これは観察で定理ではない。検証の範囲は、トーリック、対数枠で定数の $\sigma$、$L$ が十分 ample の場合であり、非トーリックな $U$ での検査ではない。これらの $n=4$ の計算は \texttt{roundJ2\_vanish\_check.py} を本論文の執筆中に再実行して確認した（$\mathbb P^2\times\mathbb P^2$ は $\mathcal O(1,2)$ で全留数が非整数のとき $(0,0,0,0,9)$。整数留数 $-1$ を含む $\mathcal O(1,1)$ の混合例でも $(0,0,0,0,9)$ となったが、これも観察）。

\begin{ex}[積\ \tP\ \tV]\label{ex:prod}
$(X_1,D_1,L_1)$, $(X_2,D_2,L_2)$ が定理 \ref{thm:B}（または \ref{thm:main2}）の対象で、$\sigma=\sigma_1\boxplus\sigma_2$（ブロック対角）のとき、複体は外部テンソル積なので Künneth により $P_{\HH}=P_1\cdot P_2$（$\HH$ の定義が超コホモロジーで有限次元なので）。次元 4 の $\mathbb P^2\times\mathbb P^2$（$\mathcal O(1,1)$、$\rho=0$）では $P_1=P_2=[1,2,4]$ のとき積は $[1,4,12,16,16]$ で、weight 公式による一般の $\sigma$ の値 $[1,4,6,4,10]$ と異なる（観察 \ref{obs:gen}）。積構造をもつ $\sigma$ は (GEN$_Q$) を満たさない特殊な場合である。
\end{ex}

%==============================================================
\section{正則型との比較：平坦なシンプレクティックファイバー束}\label{sec:regular}
%==============================================================
これまでの対象（型が跳ぶ）と対照的な、型が一定の「正則型」を扱い、次元の増え方が質的に違うことを示す。

\subsection{設定と $\omega$ への非依存}
\begin{defn}[設定 (S)]\label{defn:S}
$X$ をコンパクト複素多様体、$\tilde X\to X$ を普遍被覆、$\Gamma=\pi_1(X)$、$(Y,\omega_Y)$ をコンパクト連結なシンプレクティック多様体とする。準同型 $\rho:\Gamma\to\mathrm{Symp}(Y,\omega_Y)$ に対し $M=(Y\times\tilde X)/\Gamma$ とおく。「$Y$ 上のシンプレクティック構造と $X$ 上の複素構造の積」の一般化複素構造は $\Gamma$-不変なので $M$ に降り、型 $k=\dim_\C X$ の正則型の一般化複素構造が得られる。葉はファイバー $Y_x$、葉空間は $X$。$E=\pi^*G$（$G$ は $X$ 上の正則直線束）を SGH（強一般化正則）直線束とする。\tC\ 正則型・SGH の定義：Pal--Poddar \cite{PP} Def.\ 3.1 ほか（番号は\tQ）。
\end{defn}
この設定には、直積 $Y\times X$（$\Gamma$ 自明）、$S^2$-束（ホロノミーが $SO(3)$ に入るもの）、Kodaira--Thurston 型の $T^2$-束などが含まれる。$\HH(E)=H^*(M;C(L;E))$ は第 2 節と同様に定義される。

\begin{prop}\tP\label{prop:omega}
$\tilde\omega$ を $\mathrm{pr}_Y^*\omega_Y$ が降りて得られる $M$ 上の閉 2-形式とする。$e^{i\tilde\omega}$ は Courant 自己同型で、$L=e^{i\tilde\omega}L_0$（$L_0=T_{F,\C}\oplus T^{0,1}_{\mathrm{hor}}\oplus N^*$、$N^*=\pi^*T^{*1,0}_X$）。よって $L\cong L_0$ は Lie 亜群として同型で、一般化 Dolbeault 複体は
\[C(L;E)=\Omega^\bullet_{F_\C}(M;E\otimes\wedge^\bullet T^{1,0}_X),\qquad F_\C=T_{F,\C}\oplus T^{0,1}_{\mathrm{hor}}\ (\text{包合的})\]
と同一視でき、$N$ は SGH 束に自明に作用する。コホモロジーは $\omega$ を見ない（葉付き複素構造だけを見る）。
\end{prop}
\begin{proof}
局所的には積の主張で、すべての対象が $\Gamma$-不変。$d\tilde\omega=0$ なので $e^{i\tilde\omega}$ は Courant 括弧を保つ。$N$ が自明に作用するのは SGH の定義（横断方向に正則、葉に沿って平坦）による。\qedhere
\end{proof}

\subsection{主定理：Leray と Lefschetz}
\begin{thm}[正則型の $\HH^*$]\tP\ \tQ\label{thm:L}
設定 (S) を仮定し、さらに (HL)：$(Y,\omega_Y)$ は hard Lefschetz をみたす（例えば $Y$ が Kähler）とする。$V^j$ を $H^j(Y)$ を係数とする平坦束（Gauss--Manin）、$E=\pi^*G$ とすると
\[\HH^k(M,E)\ \cong\ \bigoplus_{a+b+j=k}H^a\bigl(X,\ \wedge^bT_X\otimes V^j\otimes G\bigr).\]
$X$ に Kähler や射影性は要らない。
\end{thm}
\begin{hyptab}
(S) 平坦なファイバー束 & 命題 \ref{prop:omega} による分解 $L\cong L_0$、Ehresmann 的な局所自明性（ファイバーのコホモロジーが局所系 $V^j$ になる） & 手順 (i)\\\midrule
$Y$ コンパクト連結 & $H^j(Y)$ が有限次元で、葉方向コホモロジー層が $V^j$ の滑らかな切断層になる & 手順 (i)\\\midrule
(HL) hard Lefschetz & $\Lambda=\omega_F\wedge$ が $V^{m_Y-i}\to V^{m_Y+i}$ の同型（平坦束の同型）を与える & 手順 (iii)：Deligne の補題で $d_r=0\ (r\ge2)$\\\midrule
$X$ コンパクト & 各 $H^a(X,\cdot)$ が有限次元 & 次元の和\\\midrule
$E=\pi^*G$（SGH） & $N$ が自明に作用するので $\wedge^bT_X$-次数 $b$ が直和因子（Poisson 的な $d_1$ がない） & 手順 (i)：$E_2$ の形
\end{hyptab}
\begin{proof}
(i)（二重複体）命題 \ref{prop:omega} により $C(L;E)=\Omega^\bullet_{F_\C}(M;E\otimes\wedge^\bullet T_X)$。これは「葉方向の de Rham」と「$X$ 上の Dolbeault（係数 $\wedge^bT_X\otimes G$）」の二重複体を作る。ファイバーがコンパクトなので葉方向コホモロジー層は $V^j$ の $C^\infty$ 切断の層であり、その上の $\bar\partial$ は Gauss--Manin 接続の $(0,1)$-部分で、持ち上げの取り方によらない（持ち上げの違いは葉方向 Lie 微分で、閉形式に対しては完全形式）。よって
\[E_2^{a,j}=\bigoplus_bH^a(X,\wedge^bT_X\otimes V^j\otimes G).\]
(ii) $\omega_F:=\tilde\omega|_F$ の混合成分は 0（$\tilde\omega$ は局所的に $\mathrm{pr}_Y^*\omega_Y$）。したがって $\Lambda=\omega_F\wedge$ は二重複体の鎖準同型（$d_F$ と $\bar\partial$ に可換）で、スペクトル系列の各項に作用し $d_r$ と可換。$E_2$ 上では各点のファイバー上の Lefschetz 作用素 $V^j\to V^{j+2}$ である。
(iii)（Deligne の補題）hard Lefschetz により $V^{m_Y-i}\to V^{m_Y+i}$ が同型。$x$ を重み $s$ の原始元とすると $\Lambda^{s+1}d_2x=0$ だが、$d_2x$ はファイバー次数 $m_Y-s-1$ の側にあり、そこでは $\Lambda^{s+1}$ は単射。よって $d_2=0$。帰納的に $d_r=0\ (r\ge2)$、次元は足し合わさる。\qedhere
\src{Deligne \cite{Del68}（Publ.\ Math.\ IH\'ES 35）。Leray スペクトル系列：Bredon \cite{Bre}（該当箇所は\tQ）。}
\end{proof}
\tQ\ 概略にとどまる点：(i) で、$\bar\partial$ が Gauss--Manin 接続の $(0,1)$-部分であることと、葉方向 Dolbeault 型の二重複体の $E_2$ の同定を詳細には書いていない。(HL) を外すと $d_2\ne0$ になりうる\tK（具体例は計算していない）。

\begin{cor}[次数別 Riemann--Roch]\tP\label{cor:L}
定理 \ref{thm:L} の仮定に加え、$X$ が射影的で $G$ が豊富とする。ある $m_0$ があって、$m\ge m_0$ ならば $H^a(X,\wedge^bT_X\otimes V^j\otimes G^m)=0\ (a>0)$、したがって
\begin{gather*}\HH^k(M,E^m)=\bigoplus_{b+j=k}H^0\bigl(X,\wedge^bT_X\otimes V^j\otimes G^m\bigr),\\ \dim\HH^k(M,E^m)=\sum_{b+j=k}b_j(Y)\int_X\ch(\wedge^bT_X)\,e^{m\,c_1(G)}\,\Td(X).\end{gather*}
右辺は各 $k$ について $m$ の多項式である。
\end{cor}
\begin{hyptab}
$G$ 豊富、$X$ 射影的 & Serre 消滅を連接層 $\wedge^bT_X\otimes V^j$ に使う & 全次数 $a>0$ の項の消滅\\\midrule
$V^j$ が平坦束 & $\ch(V^j)=\mathrm{rk}\,V^j=b_j(Y)$ & $b_j(Y)$ が外に出る\\\midrule
定理 \ref{thm:L} & $\HH^k$ を $\bigoplus H^a$ に分解 & 最初の等式\\\midrule
HRR & $a>0$ が消えるので $H^0$ の次元が $\chi$ に等しい & 最後の等式
\end{hyptab}
\begin{proof}
Serre 消滅 \tC\ \cite{Hart}（III.5.2）を局所自由層 $\wedge^bT_X\otimes V^j$ に適用し、有限個の $(b,j)$ の最大の $m_0$ を取る。残りは HRR と $\ch(V^j)=b_j(Y)$。\qedhere
\end{proof}
全 Euler 標数は $\sum_k(-1)^k\dim\HH^k=\chi(Y)\cdot(-1)^{\dim X}\chi(X)$ で、定理 \ref{thm:E}（のこの設定への拡張）と整合する\tV（例で確認）。各次数ごとの次元は $m$ の多項式として増えるが、全体の交代和は $m$ によらない。

\subsection{Pal--Poddar の定理との関係}\label{subsec:PPcompare}
\begin{thm}[Pal--Poddar, Thm.\ 11.12]\tC
$M$ をコンパクトな正則型一般化複素多様体、$M/\mathcal S$ を Kähler 軌道体、$E$ を正の SGH 直線束とする。(1) $p+q>k$ ならば $H^q(M,(\mathcal G^*M)^p\otimes E)=0$。(2) 任意の SGH 束 $E'$ に対し、$m\ge m_0$、$q>0$ なら $H^q(M,E'\otimes E^m)=0$。
\src{\cite{PP}, Thm.\ 11.12（仮定・主張は \texttt{pp\_recheck.md} での再読による）。}
\end{thm}
\begin{prop}\tP\ \tQ\label{prop:PPcompare}
設定 (S) では、Pal--Poddar の群 $H^q(M,(\mathcal G^*M)^p\otimes E)$ は $H^q(X,\Omega^p_X\otimes G)$（葉に沿って定数な横断形式）として読める。そのとき (1) は $X$ 上の古典的な Kodaira--Nakano 消滅、(2) は $E'=\pi^*(\wedge^bT_X\otimes V^j)$（SGH 束）に対する古典の Serre 消滅であり、系 \ref{cor:L} の消滅 $H^a(X,\wedge^bT_X\otimes V^j\otimes G^m)=0\ (a>0)$ と同じである。したがって横断次数での消滅は Pal--Poddar Thm.\ 11.12(2) を $E'=\wedge^bT_X\otimes V^j$ に適用したものに一致し、系 \ref{cor:L} は「Pal--Poddar の Serre 型消滅 $+$ Leray $+$ Lefschetz」である。
\end{prop}
\begin{rem}[\tQ]
文字通りの $M$ 上の層コホモロジーを取ると、$R^j\pi_*\pi^{-1}\ne0$ のために葉方向因子 $H^j(Y)$ が余分に現れる。Pal--Poddar の Thm.\ 8.9, Cor.\ 8.10 は $H^q(M,(\mathcal G^*M)^p)$ を横断 Dolbeault コホモロジーと同定するので、彼らの $H^q$ は葉空間上で読むべきと思われるが、原文でそのように明記されているかは確認していない。また $\HH^k$ の消滅が言えるのは横断次数 $a>0$ についてだけで、$H^j_{dR}(Y)$ 因子は各次数 $j$ に残る。
\end{rem}

\subsection{例}
\begin{ex}[直積 $Y\times\mathbb P^1$]\tP\ \tV\label{ex:prodY}
$X=\mathbb P^1$、$G=\mathcal O(m)$、$m\ge-1$：定理 \ref{thm:L} と $H^0(\mathbb P^1,\mathcal O(m))=m+1$、$H^0(\mathbb P^1,T\otimes\mathcal O(m))=m+3$ から $P_{\HH}(t)=P_Y(t)\bigl((m+1)+(m+3)t\bigr)$。$Y=S^2$：$(m+1)(1+t^2)+(m+3)(t+t^3)$（$m=0$：$1,3,1,3$）。$Y=T^2$：$m=0$ で $1,5,7,3$、一般に $(m+1,\ 3m+5,\ 3m+7,\ m+3)$。どちらも $m$ について線形に増える。
\end{ex}
\begin{ex}[Kodaira--Thurston：$T^2\to M\to E$]\tP\ \tV\ \tQ
モノドロミーがせん断 $\psi=\begin{pmatrix}1&1\\0&1\end{pmatrix}$ のとき $V^0=V^2=\mathcal O$、$V^1=F_2$（Atiyah \cite{Atiyah} の非自明な冪単な平坦束、$h^0=h^1=1$\tC）。$G=\mathcal O$ では $P_{\HH}(t)=(1+t)^2(1+t+t^2)=1+3t+4t^2+3t^3+t^4$。独立な確認：Heisenberg 代数 $\mathfrak h_3$ の葉方向複体に $\wedge N$ を掛けた Lie 環複体の Betti 数が $(1,3,4,3,1)$ で一致\tV。この確認は不変形式のモデルが $\HH^*$ を計算するという Nomizu 型の定理を仮定しており、その定理は証明していない\tQ（定理 \ref{thm:L} 自体はそれを使わない）。$\deg G=d>0$ では $P_{\HH}(t)=d(1+t)^3$、$\deg G=0$、$G\not\cong\mathcal O$ では $P_{\HH}=0$。
\end{ex}
\begin{ex}[Bieberbach 商、$S^2$-束]\tP\ \tV
$(T^2\times E)/\mathbb Z_2$（$(y,x)\mapsto(-y,x+\tau/2)$）：$V^1=L_{-1}\oplus L_{-1}$（非自明な符号表現）は $G=\mathcal O$ で非輪状、$P_{\HH}(t)=(1+t)^2(1+t^2)$（Künneth の $\mathbb Z_2$-不変部分と一致）。ホロノミーが $SO(3)$ の $S^2$-束では $V^j$ が自明束なので、数値は直積と同じ。
\end{ex}
\begin{rem}\tP
葉方向に非自明な平坦なホロノミー指標 $\chi$ をもつ SGH 束は、コンパクトな底では指標が局所定数になる（最大値原理）。$Y$ がトーラスで $\chi\ne1$ なら $H^*(Y;L_\chi)=0$ で $\HH=0$。
\end{rem}
\tC\ これらの例の数値（Kodaira--Thurston, Bieberbach）は総合版の記録からの転記で、本論文の執筆中には再計算していない\tQ。

%--------------------------------------------------------------
\subsection{型が跳ぶ場合との対比：多項式と周期的}\label{subsec:contrast}
%--------------------------------------------------------------
\begin{center}\small
\begin{tabular}{>{\raggedright\arraybackslash}p{3.4cm}>{\raggedright\arraybackslash}p{4.4cm}>{\raggedright\arraybackslash}p{5.8cm}}\toprule
 & 正則型（$Y=S^2$, $X=\mathbb P^1$） & 型が跳ぶ（$\mathbb{CP}^2$、滑らかな 3 次曲線）\\\midrule
$m$ が大きいとき $\dim\HH^k$ & $m$ の線形関数 $(m+1),(m+3),(m+1),(m+3)$ & 有界：$m\in3\Z_{\ge0}$ で $(1,0,2,0,0)$、それ以外で $(0,0,3,0,0)$\\\midrule
$m=0$ & $(1,3,1,3)$ & $(1,0,2,0,0)$\\\midrule
$m=1$ & $(2,4,2,4)$ & $(0,0,3,0,0)$\\\midrule
$m=3$ & $(4,6,4,6)$ & $(1,0,2,0,0)$\\\midrule
全 Euler 標数 & $\chi(Y)\cdot(-1)^{k}\chi(X)=-4$（$k=\dim X=1$） & $\chi(\mathbb{CP}^2)=3$（$m$ によらない）\\\bottomrule
\end{tabular}
\end{center}
\tV\ 右列は例 \ref{ex:cp2}（$m=-15,\dots,15$ で直接計算と一致）、左列は系 \ref{cor:L}。
\begin{obs}\tP\label{obs:contrast}
両者はともに $E_1^{a,b}=H^a(\wedge^bT\otimes E)$ から出発する。正則型では SGH 束で $N$ が自明に作用するので $d_1=0$（定理 \ref{thm:L}）。型が跳ぶ場合は $d_1=[\sigma,\cdot]+V\wedge$ が $H^0$ の項の間で\emph{ほぼ完全}になり、$H^0(\wedge^bT\otimes L^m)$ が $m$ について増えるにもかかわらず、残るコホモロジーは $U=X\setminus D$ の位相で決まる有界な量になる。「$m$ の多項式」か「周期的」かの違いは、$d_1$ の消滅か（ほぼ）完全性かに帰着する。
\end{obs}
\begin{conj}[統一予想（不正確）]\tK\label{conj:U}
有限個の型の層をもつコンパクトな一般化複素多様体と、正則な層に沿って $N$ が自明に作用する加群 $E$ について、$\HH^*(E)$ は「葉空間上の葉方向 de Rham $\otimes$ 横断 Dolbeault」の層ごとの寄与に濾過され、各層で $(c)$ 型（コンパクト台）か非 $(c)$ 型かは指数根（留数）で決まる。
\end{conj}
\noindent この予想は正確な形では定式化されておらず、上の 2 族（正則型と対数シンプレクティック）以外での検証はない。対応の目安は次の表である。
\begin{center}\small
\begin{tabular}{>{\raggedright\arraybackslash}p{3.0cm}>{\raggedright\arraybackslash}p{5.4cm}>{\raggedright\arraybackslash}p{5.6cm}}\toprule
 & 正則型 & 型が跳ぶ（対数シンプレクティック）\\\midrule
葉方向 de Rham & コンパクトな葉 $Y$ の $H^j(Y)$（局所系 $V^j$） & 開葉 $U$ の $H^j(U;L_U)$\\\midrule
横断方向 & $X$ 上の Dolbeault & $D$ の各層上の複体 $Q_T$\\\midrule
境界の扱い & なし（$Y$ はコンパクト） & 留数による $Rj_*$ / $j_!$ の選択（Deligne）\\\midrule
モノドロミー & Gauss--Manin & $e^{2\pi\sqrt{-1}r}$（$D$ のまわり）\\\midrule
退化 & Lefschetz で $d_r=0$ & $U$ は非コンパクトで類似物なし、$d_1\ne0$\\\bottomrule
\end{tabular}
\end{center}

%==============================================================
\section{Kodaira 型の問題}\label{sec:kodaira}
%==============================================================
\subsection{古典の Kodaira の埋め込み定理：どこで何を使うか}
\begin{thm}[Kodaira の埋め込み定理]\tC\label{thm:Kod}
$X$ をコンパクト複素多様体、$L$ を正の正則直線束とする。このとき $X$ は射影的で、十分大きな $m$ に対し $|L^m|$ が定める写像 $X\to\mathbb P(H^0(X,L^m)^*)$ は埋め込みである。
\src{Kodaira \cite{Kod54}。教科書：Griffiths--Harris \cite{GH} Ch.\ 1 \S4、Huybrechts \cite{Huy} Prop.\ 5.3.1（番号は\tQ）。}
\end{thm}
\begin{hyptab}
$X$ コンパクト & $H^0(X,L^m)$ が有限次元で $\mathbb P^N$ への写像が定義できる。閉埋め込み & 証明の最初と最後\\\midrule
$L$ が正（$i\Theta_L>0$） & 曲率形式が Kähler 形式になり、Kodaira 消滅が使える & 点 $x$ での爆発 $\tilde X$ で $H^1(\tilde X,\pi^*L^m\otimes\mathcal O(-2E_x))=0$ を示す段\\\midrule
正則断面が十分多い & 点と接ベクトルを分離する & 全射 $H^0(L^m)\to L^m\otimes\mathcal O/\mathfrak m_x^2$ が消滅から出る
\end{hyptab}
証明のどの段で、断面が\emph{点を区別する}ことが使われるかが重要である。一般化正則（GH）断面ではこの段が壊れる。

\subsection{一般化正則断面では埋め込めない}
\begin{prop}\tP\label{prop:one}
$X$ を連結なコンパクト複素多様体、$\sigma$ を稠密な開集合 $U_0$ で非退化な正則 Poisson 構造、$L$ を Poisson 加群とする。このとき $\dim_\C\HH^0(L)\le1$。GH 断面は $U_0$ の 1 点での値で決まり、GH 断面による Kodaira 写像は定値（または定義できない）。
\end{prop}
\begin{hyptab}
連結、$U_0$ 稠密 & 一致の定理が使える & (iii)\\\midrule
$\sigma$ が $U_0$ で非退化 & $\sigma^\sh:A|_{U_0}\to T$ が同型になり、$A$-接続が通常の平坦接続に化ける & (i)\\\midrule
$L$ が直線束 & 平坦断面は 1 点の値で決まる & (ii)
\end{hyptab}
\begin{proof}
(i) $U_0$ 上で $\nabla_vs:=D_{(\sigma^\sh)^{-1}v}s$ とおくと、Leibniz 則から $\nabla$ は通常の平坦正則接続になる。GH 断面は $\nabla s=0$ をみたす（観察 \ref{obs:casimir}）。(ii) $U_0$ は連結（連結多様体から真の解析的部分集合を除いたもの）で $L$ は階数 1 なので、平坦断面の空間は $\le1$ 次元で、断面は 1 点の値で決まる。(iii) $s$ は $s|_{U_0}$ で決まる（一致の定理）。\qedhere
\end{proof}
\begin{ex}\tV
$\mathbb{CP}^2$ と $\Pi_F$（$F$ は滑らかな 3 次式）：$\HH^0(\mathcal O(m))=\C F^k$（$m=3k\ge0$）、それ以外は 0（例 \ref{ex:cp2}）。$\mathbb P^1\times\mathbb P^1$ と滑らかな $(2,2)$ 曲線：$\HH^0(\mathcal O(m,m))=\C F^{m/2}$（$m\in2\Z_{\ge0}$、例 \ref{ex:p1p1}）。どの $m$ でも GH 断面による埋め込みは存在しない。
\end{ex}
\begin{prop}\tP\label{prop:leaf}
$M$ をコンパクトな一般化複素多様体、$\Jc$ は稠密な開集合 $W$ 上で正則型 $k<n$（葉の次元 $2(n-k)>0$）、$L$ を一般化正則直線束とする。任意の GH 断面 $s_0,\dots,s_N$ について $\varphi=[s_0:\cdots:s_N]$（基点集合の外で定義）は $W$ の各葉の上で定値で、したがって $W$ 上はしめこみでない。
\end{prop}
\begin{hyptab}
正則型（局所モデルが積） & $\pi_T(L)\supset T_F\otimes\C$ が使える & 局所 Darboux（Gualtieri \cite{Gua11} Thm.\ 4.3\tC）\\\midrule
$k<n$ & 葉の次元が正 & 葉に沿って動ける\\\midrule
$L$ 一般化正則 & $L$-接続 $D$ が $D_\xi(fs)=\pi_T(\xi)(f)s+fD_\xi s$ をみたす & GH 断面 $\Rightarrow$ 葉方向に平行
\end{hyptab}
\begin{proof}
積モデル $L=\mathrm{graph}(\omega_Y)\oplus(T^{0,1}\oplus T^{*1,0})$ では $\pi_T(L)=T_Y\otimes\C\oplus T^{0,1}$。$\xi\in L$ が葉方向の成分をもつとき $D_\xi s=0$ は、$s$ が葉方向の平坦接続 $\nabla^F$ に関して平行であることを意味する。葉内の道に沿った平行移動 $P:L_x\to L_y$ について、すべての GH 断面で $s_i(y)=Ps_i(x)$。ゆえに比 $[s_0:\cdots:s_N]$ は $x,y$ で一致する。\qedhere
\end{proof}
\begin{rem}[古典の証明のどこが壊れるか]\tP
古典の証明が使うのは「断面が点と接ベクトルを分離する」ことである。命題 \ref{prop:one} では分離する断面が（定数倍を除いて）1 つ以下、命題 \ref{prop:leaf} では同じ葉の点を分離できない。仮定（正の曲率）は断面の存在を保証するだけで、断面が GH（平坦）であることは保証しない。Pal--Poddar の「正の SGH 直線束」（Def.\ 11.10, 11.11\tC\ \tQ）は横断方向でだけ正であることを要求しており、この状況に対応する。
\end{rem}

\subsection{葉空間版 Kodaira}
\begin{thm}[葉空間版 Kodaira]\tP\label{thm:K1leaf}
$M$ をコンパクトな正則型一般化複素多様体で、次をみたすとする。
\begin{itemize}
\item[(K1)] シンプレクティック葉への射影 $\pi:M\to X$ は、連結なコンパクトファイバー（$=$ 葉）をもつ固有な沈めこみ（または軌道体の沈めこみ）で、$X$ は $k$ 次元のコンパクト複素多様体（軌道体）。
\item[(K2)] $L$ は SGH 直線束で、$L\cong\pi^*L_X$（SGH 束として）となる正則直線束 $L_X\to X$ がある（葉に沿ったホロノミーが自明）。
\item[(K3)] $L_X$ は豊富。
\end{itemize}
このとき $m_0$ があって、$m\ge m_0$ ならば $\HH^0(M,L^m)=\pi^*H^0(X,L_X^m)$ で、GH 断面による写像 $\varphi_m:M\to\mathbb P^N$ は $\varphi_m=\iota_m\circ\pi$。ここで $\iota_m:X\hookrightarrow\mathbb P^N$ は古典の Kodaira 埋め込み。したがって $\varphi_m$ の各ファイバーはちょうど 1 枚の葉である。軌道体の場合は $\iota_m$ は粗モジュライ空間の埋め込み（$m$ を等方群の位数の倍数にとる、Baily \cite{Bai}。記憶による引用\tQ）。
\end{thm}
\begin{hyptab}
(K1) 連結・コンパクトなファイバー & 葉の上の平坦断面は定数（命題 \ref{prop:leaf} の逆向き）。ファイバーが連結だから $\varphi_m$ の逆像がちょうど 1 枚の葉 & (i), (iii)\\\midrule
(K2) ホロノミー自明 & 葉上で $L^m$ が自明な平坦束になり GH 断面が存在する。非自明なら葉に沿って $\HH^0=0$ & (i)\\\midrule
(K3) $L_X$ 豊富 & 古典 Kodaira（定理 \ref{thm:Kod}）を $X$ に使う & (ii)\\\midrule
軌道体の場合 & 粗空間への Kodaira（Baily） & (ii) の軌道体版
\end{hyptab}
\begin{proof}
(i) $\HH^0(M,L^m)$ は、葉に沿って平坦で横断方向に正則な $L^m$ の断面（SGH 束の定義、Pal--Poddar Cor.\ 8.10 の次数 0 部分\tC\ \tQ）。(K1), (K2) より各ファイバー上で $L^m=\pi^*L_X^m$ は自明な平坦束なので、平坦断面はファイバー上で定数、すなわち $\pi^*t$（$t$ は $L_X^m$ の断面）。横断方向の正則性は $t$ の正則性と同値。ゆえに $\HH^0(M,L^m)=\pi^*H^0(X,L_X^m)$。(ii) (K3) と定理 \ref{thm:Kod} により、$m\ge m_0$ で $\iota_m$ は埋め込み。(iii) $\varphi_m=\iota_m\circ\pi$ で $\iota_m$ は単射なので、$\varphi_m(x)=\varphi_m(y)\iff\pi(x)=\pi(y)$、すなわち同じ葉。\qedhere
\end{proof}
\begin{rem}[位置づけ]
\tP\ 直接の帰着であり深い結果ではない。\tC\ 調べた範囲で、正則型の場合の Kodaira 型の主張を明示した文献は見つからなかった。Pal--Poddar の Thm.\ 11.12 に対応する埋め込み側の主張に当たる。\tQ\ 彼らの「正」の定義との同値は未確認。型が跳ぶ場合や一般の一般化複素構造に対する Kodaira 型の主張は見つかっていない。
\end{rem}

%--------------------------------------------------------------
\subsection{Poisson--Kodaira：Poisson 対の埋め込み}
%--------------------------------------------------------------
GH 断面が（ほぼ）定数なので、問題を変えて、\emph{通常の}正則断面による Kodaira 埋め込み $\nu:X\hookrightarrow\mathbb P^N$ が Poisson 埋め込み、すなわち $\sigma$ が $\mathbb P^N$ 上の Poisson 構造 $\sigma_N$ の制限になっているかを問う。（これは筆者の提案であり、文献にある定式化ではない。）

\begin{lem}[階数 1 の Poisson 加群の $\C^*$ 不変な持ち上げ]\tP\ \tQ\label{lem:lift}
$L$ を Poisson 加群とする。全空間 $C=L^*\setminus0$ には、$\C^*$ 不変で $p_*\hat\pi=\sigma$ をみたす Poisson 構造 $\hat\pi=\sigma+u\,\theta\wedge\partial_u$ が一意にあり（局所枠 $e$、$D_\alpha e=\theta(\alpha)e$、$u$ はファイバー座標）、$L^k$ の断面は次数 $k$ の同次関数に対応する。$\hat\pi$ の Jacobi 恒等式は、$D$ の平坦性と Leibniz 則に同値。
\end{lem}
局所計算の確認は見取り図にとどまる\tQ。例：$\mathbb P^2$ で $\Pi_F$（$\mathbb C^3$ 上の 2 次 Poisson）、$L=\mathcal O(m)$（任意の $m$）。
\begin{prop}\tP\label{prop:graded}
$X$ 射影的、$L$ が豊富な Poisson 加群のとき、$R=\bigoplus_{k\ge0}H^0(X,L^k)$ は次数付き Poisson 環で、括弧の次数は 0（$\{R_a,R_b\}\subset R_{a+b}$）。
\end{prop}
\begin{proof}
補題 \ref{lem:lift} により $R$ は $C$ 上の同次関数の環で、$\hat\pi$ の括弧は次数を保つ。\qedhere
\end{proof}
\begin{prop}[almost-Poisson な 2 次の持ち上げ]\tP\ \tC\label{prop:almost}
$m\ge m_1$ のとき $R^{(m)}=\bigoplus_kR_{mk}=S/I$（$S=\mathrm{Sym}\,R_m$、$I$ は 2 次式で生成）である（Mumford \cite{Mum}。記憶による引用、頁は\tQ）。このとき $S$ 上の次数 0 の歪対称な biderivation $\{\ ,\ \}'$ で、商 $S/I$ に $R$ の括弧を誘導するものがある。
\end{prop}
\begin{proof}
$S_2\to R_{2m}$ の線形な切断 $s$ を選び、$\{y_a,y_b\}:=s(\{\bar y_a,\bar y_b\})$ とおいて Leibniz 則で拡張する。$q\in I_2$、$g\in S_1$ に対し $\{q,g\}'\equiv\{\bar q,\bar g\}=0\bmod I_3$ なので $\{I_2,S_1\}'\subset I_3$。$I$ が $I_2$ で生成されるので Leibniz 則により $\{I,S\}'\subset I$。\qedhere
\end{proof}
持ち上げの Jacobiator $\mathcal J\in\mathrm{Hom}(\wedge^3S_1,I_3)$ は一般に 0 でなく、持ち上げの選び方の差は $\epsilon\in\mathrm{Hom}(\wedge^2S_1,I_2)$。Poisson 構造への拡張は 2 次の方程式 $\mathcal J_0+\mathcal L(\epsilon)+\mathcal J(\epsilon)=0$ が解 $\epsilon$ をもつことと同値。1 次の障害類は $\mathcal J_0\bmod\mathrm{im}\,\mathcal L$。

\begin{thm}[Poisson--Kodaira、トーリック]\tP\label{thm:PKtoric}
$X$ を滑らかな射影トーリック多様体（トーラス $T=N\otimes\C^*$）、$\sigma=\sum\omega^{ij}\xi_i\wedge\xi_j$（$\xi_i=x_i\partial_{x_i}$、$\omega\in\wedge^2N_\C$）とする。$L=\mathcal O(D)$ は豊富で多面体 $P$ をもち、$L^m$ は非常に豊富（断面 $\chi^a$、$a\in mP\cap M$）とする。このとき $\nu:X\hookrightarrow\mathbb P^N$、$\nu^*y_a=\chi^a$ は、$\mathbb P^N$ 上の対角的な 2 次 Poisson 構造
\[\sigma_N=\sum_{a<b}\omega(a,b)\,y_ay_b\,\partial_{y_a}\wedge\partial_{y_b},\qquad\omega(a,b)=\sum\omega^{ij}a_ib_j\]
に関する Poisson 埋め込みである。非常に豊富な任意のべき $L^m$ で成り立つ。
\end{thm}
\begin{hyptab}
トーリック、$\sigma$ がトーラス不変 & 指標の括弧が単項式になる：$\{\chi^a,\chi^b\}=\omega(a,b)\chi^{a+b}$ & (2)。持ち上げを対角的に選べる\\\midrule
$L^m$ が非常に豊富 & $\nu$ が埋め込みになる & 結論の「埋め込み」\\\midrule
$X$ 滑らか・射影的 & 古典 Kodaira の設定 & 同上
\end{hyptab}
\begin{proof}
(1) $\sigma_N$ は Poisson：歪対称な $\lambda_{ij}$ に対し $\{y_i,y_j\}=\lambda_{ij}y_iy_j$ とおくと $\{y_i,\{y_j,y_k\}\}=\lambda_{jk}(\lambda_{ij}+\lambda_{ik})y_iy_jy_k$ の巡回和は歪対称性から 6 項が対で打ち消し合って 0。スケーリングの weight は 0 なので $\mathbb P^N$ に落ちる。(2) トーラス上で $\{\chi^a,\chi^b\}=\sigma(d\chi^a,d\chi^b)=\omega(a,b)\chi^{a+b}=\nu^*(\omega(a,b)y_ay_b)$。$\sigma$ も $\nu^*\sigma_N$ も正則なので、トーラスの閉包 $X$ 上でも $\nu$ は Poisson 写像。\qedhere
\end{proof}

\begin{obs}[$\mathbb P^2$、3 次式、$m=2$：数値的確認]\tV\label{obs:P2}
$X=\mathbb P^2$、$\sigma=\Pi_F$、Veronese $m=2$（$\mathbb C^6$、$\dim I_2=6$、未知数 $\epsilon$ は 90 個、方程式は 940 本）。自然な持ち上げの Jacobiator は Fermat で 18 個、Hesse で 19 個、$xyz$ で 1 個の三つ組、尖点でも、0 でない。完全な 2 次方程式 $\mathcal J_0+\mathcal L(\epsilon)+\mathcal J(\epsilon)=0$ を Levenberg--Marquardt 法で解くと、$F=xyz$、Fermat、Hesse（$x^3+y^3+z^3-21xyz$）、尖点（$x^3-y^2z$）のすべてで残差 $\le10^{-12}$ の解が見つかった。Fermat では有限差分でも独立に確認した（Jacobiator $\sim10^{-9}$、Veronese 錐への制限の誤差 $\sim10^{-10}$）。Fermat では線形部分 $\mathcal J_0+\mathcal L(\epsilon)=0$ だけでは解がなく、2 次の項が本質的である。$m=1$ は自明、$m\ge3$ は未知数が 1215 個で重すぎて計算していない。スクリプト \texttt{roundE\_*.py}（総合版の記録からの転記、再実行していない\tQ）。数値的な存在であって、厳密な代数的証明（certificate）ではない。
\end{obs}
\begin{obs}[$\mathbb P^1\times\mathbb P^1$ の Segre 埋め込み（$m=1$）：Sklyanin 型構造\ \tV]\label{obs:Segre}
$\mathbb P^3$ 上の 2 次 Poisson 構造 $\{x_0,x_i\}=(J_j-J_k)x_jx_k$、$\{x_j,x_k\}=x_0x_i$（$(i,j,k)$ は $1,2,3$ の巡回、$J=(2,3,7)$）について、Jacobi 恒等式（4/4 の三つ組）と、2 次の Casimir がちょうど 2 つあり一般のランクが 2 であることを有理数で厳密に確認した。したがって Casimir の束の各 2 次曲面は Poisson 部分多様体で、滑らかなものは $\mathbb P^1\times\mathbb P^1$、$\sigma$ の零点は滑らかな楕円 4 次曲線 $E=Q_1\cap Q_2$ となる（「単純零点がちょうど $E$」はランクと葉の議論による同定で、計算してはいない）。これは古典的で新しくない（Feigin--Odesskii、Odesskii、Hua--Polishchuk の仕事\tC\ \tQ）。すべての滑らかな $(2,2)$ 曲線がこの形で現れるか、特異な $E$ の場合は未検査\tK。
\end{obs}
\begin{conj}[Poisson--Kodaira]\tK\label{conj:PK}
$(X,\sigma)$ を射影的な正則 Poisson 多様体、$L$ を豊富な Poisson 加群とする。このとき $m_0$ があって、$m\ge m_0$ ならば $|L^m|$ による Kodaira 埋め込みは、$\mathbb P^N$ 上のある 2 次 Poisson 構造 $\sigma_N$ に関する Poisson 埋め込みである。
\end{conj}
\noindent 根拠：定理 \ref{thm:PKtoric}、観察 \ref{obs:P2}、観察 \ref{obs:Segre}。未解決：一般の障害理論（$\mathcal J_0\bmod\mathrm{im}\,\mathcal L$ は Poisson コホモロジー型の類で、$m\gg0$ では Koszul 複体の $H^1$ の Serre 型消滅で消えると期待されるが、証明していない）。$\mathbb P^2$ の $m=3$ の Fermat 型は未計算で、PK に有利でも不利でも証拠がない（$xyz$ の場合はトーリックの定理で証明済みで、独立な検査ではない）。反例は見つかっておらず、$m$ が小さいと破れる可能性は残る。$\sigma=0$ なら古典の Kodaira に一致する。正則型 $M\to X$ で葉が $Y$ のとき GH 断面は $X$ を経由する（命題 \ref{prop:leaf}）ので、$(M,\sigma_Y)$ を Poisson 対として埋め込むには GH でない通常の正則断面を使う必要があり、それが本節の定式化である。

%==============================================================
\section{Riemann--Roch 型の公式と指数予想}\label{sec:RR}
%==============================================================
\subsection{Poincar\'e 多項式の公式}
定理 \ref{thm:E} は全 Euler 標数が位相的であることを述べた。次数ごとの次元は次の形になる。

\begin{thm}[RR 型の Poincar\'e 多項式]\tP\ \tV\label{thm:K}
定理 \ref{thm:R} の仮定（SNC、対数シンプレクティック、$r_i\notin\Z_{\ge1}$、(GEN$_Q$)）に加えて、(i) $U$ がアフィン、(ii) $n$ が偶数（対数シンプレクティックなら自動）とする。このとき
\[P_{\HH}(t)\;=\;\chi(X)\,t^n\;+\;\sum_{k=0}^{n-1}h^k(U;L_U)\,\bigl(t^k-(-1)^kt^n\bigr).\]
\end{thm}
\begin{hyptab}
簡約定理の仮定 & $\HH^*\cong H^*(U;L_U)\oplus\C^{\#D_{(n)}}[-n]$ を与える & (1)\\\midrule
$U$ アフィン & Artin 消滅：局所系で $H^k(U;L_U)=0\ (k>n)$ & (2)。ないと「$k<n$ の和」に収まらない\\\midrule
$n$ 偶数 & $\sum_k(-1)^kh^k+\#D_{(n)}=\chi(X)$ が符号なしで使える（定理 \ref{thm:E}） & (3)\\\midrule
射影性・豊富性 & 不要 & ---
\end{hyptab}
\begin{proof}
(1) 定理 \ref{thm:R} より $P_{\HH}=\sum_{k\le n}h^kt^k+\#D_{(n)}t^n$（$k>n$ は (2) で消える）。(2) Artin の消滅定理 \cite{SGA4} により $k>n$ で $h^k=0$。(3) 定理 \ref{thm:E}（$n$ 偶数）より $\sum_{k\le n}(-1)^kh^k+\#D_{(n)}=\chi(X)$、すなわち $h^n+\#D_{(n)}=\chi(X)-\sum_{k<n}(-1)^kh^k$。これを $t^n$ の係数に代入する。\qedhere
\end{proof}
\begin{cor}\tP
すべての $r_i\notin\Z$ なら（定理 \ref{thm:V}）$P_{\HH}(t)=\chi(X)t^n$。すべての $r_i\in\Z_{\le0}$ で $L_U$ が自明なら $h^k=b_k(U)$。$n=2$ では $h^0=1$、$h^1=b_1$ で $P=\chi(X)t^2+(1-t^2)+b_1(t+t^2)$。
\end{cor}
\paragraph{検証。}\tV\ (a) $n=2$ の表 188 例、(b) トーリック $n=2,4$ の 61 例（$\rho\in M$ のとき $h^k=\binom nk$、それ以外は 0）で、実測の Poincar\'e 多項式と公式の予言が全例一致（食い違い 0、総合版の記録。再実行していない\tQ）。$\mathbb P^2$、$\rho\in M$ のトーリックでは $P=3t^2+(1-t^2)+2(t+t^2)=1+2t+4t^2$。これは定理 \ref{thm:B} の $[r\le0]$ の場合に一致する。$r\ge2$ の側は双対性（定理 \ref{thm:dual}）から得られる。

\subsection{次元の 3 つの振る舞い}
$L=L_0^m$ を取って $m$ を動かすとき、次数ごとの次元は設定によって質的に異なる振る舞いをする。
\begin{center}\small
\begin{tabular}{>{\raggedright\arraybackslash}p{3.0cm}>{\raggedright\arraybackslash}p{4.4cm}>{\raggedright\arraybackslash}p{6.4cm}}
\toprule 振る舞い & 次数 $k$ ごとの次元 & 例\\\midrule
(A) 正則型 & $m$ の多項式（$m\gg0$）\tP & $X=\mathbb P^1$、$Y=T^2$（直積）、$G=\mathcal O(m)$：$(m+1,\ 3m+5,\ 3m+7,\ m+3)$（$m\ge-1$）\\\midrule
(B) 対数シンプレクティック & 留数 $r=-m/3$ の整数性による周期的な跳び\tP\ \tV & $\mathbb P^2$、滑らかな 3 次曲線：$m=3k\ge0$ で $(1,0,2,0,0)$、それ以外 $m\ge-3$ で $(0,0,3,0,0)$（例 \ref{ex:cp2}）\\\midrule
(C) 全 Euler 標数 & $m$ に依らない\tP & (A) では $\sum(-1)^kh^k=0$、(B) では $3=\chi(\mathbb P^2)$\\\bottomrule
\end{tabular}
\end{center}
(A) では $(m+1)-(3m+5)+(3m+7)-(m+3)=0=\chi(Y)\chi(X)$ と整合する。読み方：(C) より、Riemann--Roch の素朴な形「$\chi(E)=\int\ch(E)\Td$ が $c_1(E)$ の多項式」は一般化幾何では退化し（$c_1$ を見ない）、情報は次数つきの次元にある。(A) では次数ごとの多項式、(B) では定理 \ref{thm:K} のように「留数の整数性にだけ依存する位相的な量 $h^k(U;L_U)$」と $\chi(X)$ に分解される。多項式的な振る舞いが全次数で成り立つのは正則型だけである。

\subsection{先行研究との比較}\label{subsec:compare}
\paragraph{Hodge 理論（形式係数）。}\tC\ Gualtieri は一般化 Kähler 多様体上の一般化 Hodge star、Born--Infeld 内積、Laplacian の等式、Hodge 分解、$dd^{\Jc}$-補題を示した（\cite{GuaHodge}, \cite{GuaK}、Cavalcanti \cite{Cav}。定理番号は精読していない\tQ）。これらは自明束または $H$-捩れ係数の理論であり、直線束（曲率つき）係数の Bochner--Kodaira--Nakano 恒等式は、調べた範囲では書かれていない。本論文の結果は、この係数付き版の穴を、Hodge 理論に頼らず、対数複体と構成可能層の比較で $\beta$-変換の場合に埋めたものと位置づけられる。

\paragraph{Pal--Poddar。}\tC\ \cite{PP} は正則型（型 $k$ 一定）のコンパクト一般化複素多様体で葉空間が Kähler 軌道体の場合に、SGH 束の Hodge 理論、横断 Dolbeault 同型（Thm.\ 8.9, Cor.\ 8.10）、Hodge 分解と Serre 双対性（Thm.\ 11.2--11.7）、Kähler の恒等式（Prop.\ 11.8, Lemma 11.9）、正の SGH 直線束（Def.\ 11.10, 11.11）に対する消滅定理（Thm.\ 11.12）を与える（番号は\tQ）。違い：(i) 型が跳ぶ場合は Pal--Poddar の枠外で、本論文の $\mathbb{CP}^2$ などの例はこの場合。(ii) SGH 束は一般化正則束の真部分クラスで、本論文の Poisson 加群は $\beta$-変換構造における一般の GH 束。(iii) Thm.\ 11.12 の正確な主張は \S\ref{subsec:PPcompare}。原典の記述についての疑義（誤植・省略の可能性が高いが未検証\tQ）：主張の $(i/2\pi)$ と証明の $1/2\pi$、$gc_1$ の符号、「正のクラスから正曲率の計量が出る」ことが Def.\ 11.11 から証明されず Thm.\ 11.12 の証明が仮定していること、Def.\ 11.11 にコンパクト性・Kähler 性の明記がないこと。

\paragraph{Hu の仕事。}\tC\ Hu \cite{Hu1,Hu2}。設定は実 Lie 亜群 $\mathfrak g$ にファイバー方向の可積分複素構造を与えたもの（本論文の、$\bT_\C$ の部分束としての複素 Lie 亜群 $L$ とは別物）。Riemann--Roch（2403.13424 Thm.\ 4）は $\dbar_{\mathfrak g}$ の楕円性と単模性の仮定の下で指数公式を主張し、消滅（Thm.\ 7）は Laplacian の核についての主張で、コホモロジーの消滅は導けない。\tP\ 本論文の一般化 Dolbeault 複体 $(T^{0,1}\bowtie A^{1,0})$ は Hu の設定に入らない（$L$ は複素 Lie 亜群で $J$ をもたず、$\wedge L^*$ は一般に $(p,q)$ 二重次数をもたない）。\tK\ Hu の指数は $m$ に多項式的に依存すると見込まれるが、本論文の全 Euler 標数は $m$ に依らず、周期的な跳びを決める留数 $-m/3$ は Hu のデータに含まれない。したがって Hu の定理は「一般化 Dolbeault コホモロジーに対する Riemann--Roch」ではない。原典の不整合（記号の未定義、符号の食い違い）：要約と本文の一部で気づいたが、付録は検証できていない\tQ。

\paragraph{積分可能性。}\tC\ Gualtieri--Li \cite{GL} は対数シンプレクティック多様体の Hausdorff シンプレクティック亜群を分類し、滑らかな $D$ に沿って接する対数ベクトル場の Lie 亜群 $T_X(-\log D)$ が積分可能であることを与える。Bailey--Gualtieri \cite{BG}、Laurent-Gengoux--Stiénon--Xu \cite{LGSX2}、Crainic--Fernandes \cite{CF} は積分可能性の一般論を与える。これらは指数理論の入力としての群を与えるが、複素 Lie 亜群 $L$ 自体の積分は扱わない（定理番号は\tQ）。

\subsection{指数予想}
Hirzebruch--Riemann--Roch は Atiyah--Singer の指数定理 \cite{AS} の特別な場合であり、一般化幾何での Riemann--Roch は「一般化 Dolbeault 複体の指数定理」として定式化されるべきである。以下は\emph{すべて予想であり、証明されていない}（\tK）。

\paragraph{既存の指数定理が使えない理由。}(1) Hu \cite{Hu2} の複体は実 Lie 亜代数に付随し、我々の複体 $(\wedge^\bullet T\otimes L,\dbar+[\sigma,\cdot])$ は複素 Lie 亜代数 $T^{0,1}\ltimes A^{1,0}$（実構造なし）の複体で別物。(2) Pflaum--Posthuma--Tang \cite{PPT} の Lie 亜群の指数はコンパクトな底、ユニモジュラー性等を要する\tQ。(3) 文献に、$L$ の留数を重みとする対数 Dolbeault 複体の指数定理は見つからなかった。近い文献：$b$-対称多様体の幾何学的量子化（Braverman--Loizides--Song \cite{BLS}）、境界亜群の指数公式（Qiao--So \cite{QS}）、Lie 構造をもつ多様体の擬微分作用素（Ammann--Lauter--Nistor \cite{ALN}）、ファイバー付き境界（Mazzeo--Melrose \cite{MM}）、複素対数 Hodge（Hausel--Hunsicker--Mazzeo \cite{HHM}）。いずれも設定が異なりそのままは適用できない（題名・著者・要旨を確認したのみ\tQ）。

\paragraph{枠組み。}$A=T_X(-\log D)$ を Lie 亜代数（Gualtieri--Li \cite{GL}：対数シンプレクティックで積分可能、滑らかな $D$ で確認\tC\ \tQ）とみて、$U$ 上の複体 $\mathcal D_L=(\wedge^\bullet T\otimes L,\ \dbar+[\sigma,\cdot])$ を、$A$ から決まる重み付き Sobolev 空間（重みは留数 $r$）上の閉作用素と考える。$A$ の意味での主シンボルは可逆（補題 \ref{lem:iso} と局所計算\tV）なので、非 Fredholm 性は境界での指数的根から生じる。局所モデル（$D$ の滑らかな点）では $\nabla z^k=(k-r)z^k\,dz/z$ なので、指数的根は $k=r$ で、存在するのは $r\in\Z$ のとき。

\begin{conj}[指数予想]\tK\label{conj:index}
$X,\sigma,D,L$ を上のとおり（対数シンプレクティック、SNC）とする。
\begin{itemize}
\item[(I1)] $\mathcal D_L$ の、重み $r$ の Sobolev 空間上の閉拡大は、すべての $r_i\notin\Z$ のとき、またそのときに限り Fredholm である。
\item[(I2)] Fredholm のとき、指数（交代和）は $\chi(U)+(-1)^n\#D_{(n)}=(-1)^n\chi(X)$ に等しい（$\int c_n(T(-\log D))$ に角の寄与 $\#D_{(n)}$ を加えたもの）。
\item[(I3)] $r_i\in\Z$ のとき、$j_!$ 型と $Rj_*$ 型の 2 つの閉拡大の核・余核の差は $H^*(\partial U;L)$ の部分商で与えられ（APS 型の壁越え三角形）、$\chi(\partial U)=0$ なので指数は変わらない（系 \ref{cor:wall} と整合）。
\item[(I4)] 次数つきの精密化：$P_{\HH}(t)$ は定理 \ref{thm:K} の形（$\chi(X)t^n$ と、指数的根の核が決める跳びの項）で与えられる。
\end{itemize}
\end{conj}
\paragraph{既に証明されている部分との対応。}(I2) の値は定理 \ref{thm:E}、\ref{thm:K1} で証明済（$\chi(X)$、\tP）。(I3) の三角形は位相的なレベルでは $j_!L_U\to Rj_*L_U\to i_*i^*Rj_*L_U$ として本論文で使われ（$n=2$ で完全に確認\tP\ \tV）、この三角形の指数側の対応物が未証明。(I4) は $n=2$ では定理 \ref{thm:B} \tP、一般次元は定理 \ref{thm:K}（(GEN$_Q$) の下）\tP。
\paragraph{証明すべきこと。}(1) 複素 Lie 亜代数（$J$ なし）上の Fredholm 理論：parametrix の構成と (I1) の「$\Leftrightarrow$」。Ammann--Lauter--Nistor の枠組みを重み付きに拡張する形が自然だが未実施。(2) (I2) の指数公式：b-接束上の Atiyah--Singer 公式 \cite{BLS} の類似を $\wedge^\bullet T$ 係数・重み付きで示す。角（高い余次元の層）の寄与が問題。(3) (I3) の APS 型の壁越え：$H^*(\partial U;L)$ との同定を $n=2$ の滑らかな $D$（Gysin 列）以外で行う。(4) (I4) の次数つき版：スペクトル流の次数つき精密化（文献は見つかっていない）。(5) (GEN$_Q$) の一般の十分条件。

\subsection{三種の神器：到達点}
\begin{center}\small
\begin{tabular}{>{\raggedright\arraybackslash}p{2.4cm}>{\raggedright\arraybackslash}p{6.5cm}>{\raggedright\arraybackslash}p{6.0cm}}
\toprule 神器 & 別の形で成り立ったもの & 成り立たない／未解決\\\midrule
Riemann--Roch & 全 Euler 標数（定理 \ref{thm:E}, \ref{thm:K1}）、Poincar\'e 多項式の公式（定理 \ref{thm:B}, \ref{thm:K}）、正則型の次数別多項式（系 \ref{cor:L}） & 素朴な HRR は不成立。指数定理による統一は予想（予想 \ref{conj:index}）\\\midrule
Serre 消滅 & 横断次数 $a>0$ での消滅（定理 \ref{thm:serre}） & 全次数での消滅は不成立（注意 \ref{rem:total}）\\\midrule
Kodaira 埋め込み & 葉空間版（定理 \ref{thm:K1leaf}）、トーリック Poisson--Kodaira（定理 \ref{thm:PKtoric}） & GH 断面による埋め込みは不可能（命題 \ref{prop:one}）。一般の Poisson--Kodaira は予想（予想 \ref{conj:PK}）\\\bottomrule
\end{tabular}
\end{center}

%==============================================================
\section{未解決問題}\label{sec:open}
%==============================================================
\begin{enumerate}
\item \textbf{一般の Poisson--Kodaira}（予想 \ref{conj:PK}）：$\mathbb P^2$ の $m\ge3$、$\mathbb P^1\times\mathbb P^1$ の Segre 以外の場合、特異な $E$。障害類 $\mathcal J_0\bmod\mathrm{im}\,\mathcal L$ の消滅。
\item \textbf{(GEN$_Q$) の一般の十分条件}、命題 \ref{prop:grQ} の一般 $n$ での証明、ホロモルフィックな対数 Darboux 定理（命題 \ref{prop:genNec}(b)）。
\item \textbf{混合留数の消滅}：一部の $r_i$ が整数のとき $H^k(U;L_U)=0\ (k<n)$（定理 \ref{thm:V} の拡張）。
\item \textbf{局所公式}：準斉次でない孤立特異点、非孤立特異点、$s=1$、非被約な境界（予想 \ref{conj:NR}）、精密化 Riemann--Roch の大域化（予想 \ref{conj:refined}）。三重点の異常 $m=-2,-4$ の説明。
\item \textbf{共鳴 $r_i=1$} の一般の公式（観察 \ref{obs:res}）の証明。
\item \textbf{指数予想}（予想 \ref{conj:index}）：重み付き b-計算の Fredholm 性、指数 $=\chi(U)$、境界のコホモロジーとの関係。
\item \textbf{正則型の拡張}：非平坦なバンドル、hard Lefschetz をみたさないファイバー、統一予想（予想 \ref{conj:U}）の正確な定式化。
\item \textbf{一般の一般化複素構造}に対する Riemann--Roch と Kodaira 型の定理（本研究では未着手）。一般化 Kähler 側の直線束係数の Bochner--Kodaira--Nakano 恒等式。
\item \textbf{文献の未確認項目の解消}：Mokri の定理、LGSX Lemma 4.22 の再導出、Gualtieri \cite{Gua11} \S4 の精読、残る書誌情報（付録の一覧）。
\item \textbf{非トーリック・非積の次元 4 の検証}。候補：$\mathbb P^4$ 上の Feigin--Odesskii 型の楕円 Poisson 構造（$D$ は SNC でない）。
\end{enumerate}

%==============================================================
\appendix
\section{総合版からの変更点}\label{app:changes}
%==============================================================
本論文は総合版（Part I--V の統合、80 頁、2026-11-22 v1）を、テーマ別に 1 本の論文に書き直したものである。主張の格上げはしていない。変更を以下に列挙する。

\subsection{状態表現の更新}
\begin{enumerate}
\item \textbf{局所公式}：総合版の第 II, III 部は「観察・予想」（\tP\ \tV\ \tK）としていたが、第 V 部で準斉次な germ の $s\ne1$ について証明された。本論文は定理 \ref{thm:local} として一箇所にまとめ、(e,$\chi_h$) 基底によるより短い証明に書き換えた（旧証明の符号規約と収束に関する\tQ\ を除去した。残る\tQ\ は、局所モデルと大域 Poisson 加群の茎の同定、Saito の判定法の番号）。規則 $R(s)$ は第 II 部 v2 の訂正版。旧版の閾値の規則 $[s<\mathrm{lct}]\cdot(\cdots)$ は $\mathrm{lct}<1/2$ で誤りであることを注意し、大域の精密化予想（予想 \ref{conj:refined}）は $\mathrm{lct}\ge1/2$ に限定した。
\item \textbf{$n>2$ は未解決との記述}：簡約定理（定理 \ref{thm:R}）を、$\mathrm{gr}\,Q$ の濾過に依らない形（連鎖写像 $\varepsilon$ による (GEN$_Q$)）で証明し直した。一般 $n$ での (GEN$'$) $\Rightarrow$ (GEN$_Q$)（命題 \ref{prop:grQ}）は\tQ\ のまま。定理 V（補遺にあった）は第 7 節の本文に入れ、系 \ref{cor:Vgen} として定理 K（定理 \ref{thm:K}）の消滅部分に接続した。「定理 K の $k<n$ の消滅が未証明」という旧記述は、全 $r_i\notin\Z$ の場合に解消された（混合は未解決）。
\item \textbf{三重点の異常}：$m=-2,-4$ の相殺は観察・予想のままとした（定理 \ref{thm:local} では説明されない）。
\item \textbf{(GEN) の名称}：定理 \ref{thm:T} の仮定 (GEN)（定数行列 $A$ の算術的一般性）と、簡約定理の仮定を区別するため、後者を (GEN$_Q$) と改名した。
\end{enumerate}

\subsection{構成上の変更}
\begin{itemize}
\item 各部ごとの概観・記号表・状態表・「検証の範囲」を、第 1 節の 1 つの状態表と記号表に統合した。通し番号は定義・定理・命題・補題・系・例・注意・観察・予想で共通の 1 本。
\item 旧 Part I の「古典の三種の神器」節は、古典定理と仮定の表（第 2 節）に圧縮した。旧 Part I の可積分性についての注意と、一部の古典定理の仮定表（Kodaira--Nakano など、本論文で使わないもの）は省いた。
\item 旧 Part IV のトーリックの再証明（(GEN$'$) 経由）は、定理 \ref{thm:T} の直接証明があるため、証明を省き、第 7 節で関係を述べるにとどめた。
\item 旧 Part I $\S$ の $\mathbb{CP}^2$ の例は、定理 \ref{thm:B} からの導出（例 \ref{ex:cp2}）として書き直した。
\item 旧 Part III の積の命題（Künneth）は例 \ref{ex:prod} に、非被約な境界は観察 \ref{obs:nr} と予想 \ref{conj:NR} に統合した。
\item 旧 Part II の「Part I への影響の確認」節は、各定理の仮定表に吸収したので独立の節としては残していない。
\end{itemize}

\subsection{旧部の主張と本論文の対応}
\begin{center}\small
\begin{tabular}{>{\raggedright\arraybackslash}p{5.4cm}>{\raggedright\arraybackslash}p{9.2cm}}\toprule
旧（総合版）の主張 & 本論文 \\\midrule
Part I：GH 束 $=$ Poisson 加群、対数的平坦接続 & 命題 \ref{prop:ident}、定理 \ref{thm:main1}\\
Part I：曲面の主定理、閉公式、周期的安定化 & 定理 \ref{thm:main2}, \ref{thm:B}, \ref{thm:A2}\\
Part I：Serre 型定理、双対性 & 定理 \ref{thm:serre}, \ref{thm:dual}\\
Part II：Euler 標数の定理 & 定理 \ref{thm:E}\\
Part II：トーリック（曲面、一般次元）、共鳴 & 定理 \ref{thm:weight}, \ref{thm:C1}, \ref{thm:C2}, \ref{thm:T}、観察 \ref{obs:res}\\
Part II, III, V：局所公式、三重点、非被約 & 定理 \ref{thm:local}、観察 \ref{obs:triple}, \ref{obs:nr}、予想 \ref{conj:refined}, \ref{conj:NR}\\
Part III：Kodaira 型 & 命題 \ref{prop:one}, \ref{prop:leaf}、定理 \ref{thm:K1leaf}, \ref{thm:PKtoric}、予想 \ref{conj:PK}\\
Part IV：正則型、簡約定理 & 定理 \ref{thm:L}、系 \ref{cor:L}、定理 \ref{thm:R}\\
Part V：定理 K1（対数版 Euler 標数） & 定理 \ref{thm:K1}\\
Part V：定理 K、指数予想 & 定理 \ref{thm:K}、予想 \ref{conj:index}\\
補遺：定理 V、Poisson--Kodaira の追加検証 & 定理 \ref{thm:V}、観察 \ref{obs:Segre}\\\bottomrule
\end{tabular}
\end{center}
総合版の「定理 K1（葉空間版 Kodaira）」（第 III 部）と「定理 K1（対数版 Euler 標数）」（第 V 部）は別物であったので、本論文では前者を定理 \ref{thm:K1leaf}、後者を定理 \ref{thm:K1} とした。

\subsection{訂正・追記}
\begin{itemize}
\item 記号 $\rho$ の衝突（トーリックの $\rho\in M_\C$ と、曲面の $[D_i]$ の張る空間の次元）を、後者を $\rho_D$ に改めて解消した。
\item 準斉次 $D_5=z^2w+w^4$ の重みは $(3,2)$、次数 8（総合版の表と本文は正しい。研究メモにあった別の重みは誤り）。局所表の $A_4$ の最後の列を埋めた。
\item Hu の 2 本目の論文の題名は総合版の第 I, III 部の表記（Lie groupoid Riemann--Roch--Hirzebruch theorem and applications, arXiv:2403.13424）に統一した（題名の最終確認は\tQ）。
\end{itemize}

\subsection{検証の範囲とスクリプト}
計算はすべて有理数または有限体（$32003$, $2^{61}-1$）上の厳密な線形代数、あるいは既知の公式の当てはめであり、証明ではない。曲線の選び方への依存、および法 $p$ での還元が標数 0 の結果を代表することは検証していない。
\paragraph{本論文の執筆中に再実行したもの。}
\begin{itemize}
\item \texttt{duality\_ndim.py}, \texttt{duality\_ndim\_exact.py}：$n\le4$、すべて一致。
\item \texttt{roundH\_proof\_check.py}：14 種の準斉次 germ、不一致 0。
\item \texttt{roundK\_check.py}：175 例、188 例、61 例、不一致 0。
\item \texttt{roundJ2\_vanish\_check.py}：$n=4$ の 6 例。
\item \texttt{roundD\_check.py}, \texttt{ncd\_local\_weights.py}, \texttt{ncd\_cochain\_check.py}, \texttt{serre\_rr\_check.py}。
\item 新規に書いたもの \texttt{local\_basis\_check.py}：$(e,\chi_h)$ 基底の式 \eqref{eq:d0}, \eqref{eq:d1} の 210 検査、不一致 0。
\end{itemize}
出力は \texttt{logs/} に保存した。
\paragraph{総合版の記録からの転記で、再実行していないもの。}トーリックの大規模な表（\texttt{roundC\_*.py}, \texttt{toric\_hh.py}, \texttt{toric\_generic\_check2.py}）、三重点（\texttt{local\_unequal*.py}）、非被約（\texttt{roundF\_nonreduced.py}）、Poisson--Kodaira の数値（\texttt{roundE\_*.py}）、共鳴の公式の 108 例、正則型の例（Kodaira--Thurston, Bieberbach）。
\paragraph{文献。}引用した定理の適用範囲を原典で確認していないものが多数残る（各所の\tQ\ 項目：LGSX Thm.~4.19 と Mokri の定理の再導出、Gualtieri 2011 \S4 の精読、Saito の判定法と Cox--Little--Schenck の番号、Pal--Poddar の定義番号、Hu の書誌など）。

\subsection{確信度の低い箇所（まとめ）}
(1) 命題 \ref{prop:ident}（LGSX と Mokri に依存、再導出なし）、(2) 命題 \ref{prop:grQ}（一般 $n$ の $\mathrm{gr}\,Q$ と $\varepsilon$ の連鎖写像性）、(3) 定理 \ref{thm:local} の局所モデルと大域加群の接続、(4) 定理 \ref{thm:dual} の Koszul 符号の整合（手計算のみ）、(5) 定理 \ref{thm:L} の (i)（Gauss--Manin との同定は概略）、(6) 観察 \ref{obs:res} の $N_R$ の定義、(7) 補題 \ref{lem:lift}（見取り図）。

%==============================================================
\begin{thebibliography}{99}\small
\bibitem{AB} M.\ Abouzaid, M.\ Boyarchenko, Local structure of generalized complex manifolds, J.\ Symplectic Geom.\ 4 (2006) 43--62, arXiv:math/0511049.
\bibitem{ALN} B.\ Ammann, R.\ Lauter, V.\ Nistor, Pseudodifferential operators on manifolds with a Lie structure at infinity, arXiv:math/0304044.
\bibitem{AN} Y.\ Akizuki, S.\ Nakano, Note on Kodaira--Spencer's proof of Lefschetz's theorems, Proc.\ Japan Acad.\ 30 (1954) 266--272.
\bibitem{AS} M.\ Atiyah, I.\ Singer, The index of elliptic operators III, Ann.\ Math.\ 87 (1968) 546--604.
\bibitem{Ara} D.\ Arapura, Geometry of cohomology support loci for local systems I, arXiv:alg-geom/9612006; J.\ Algebraic Geom.\ 6 (1997) 563--597 [巻頁は記憶\tQ].
\bibitem{Atiyah} M.\ Atiyah, Vector bundles over an elliptic curve, Proc.\ London Math.\ Soc.\ (3) 7 (1957) 414--452.
\bibitem{BG} A.\ Bailey, M.\ Gualtieri, Integration of generalized complex structures, arXiv:1611.03850.
\bibitem{BLS} M.\ Braverman, Y.\ Loizides, Y.\ Song, Geometric quantization of $b$-symplectic manifolds, arXiv:1910.10016.
\bibitem{Bai} W.\ L.\ Baily, On the imbedding of V-manifolds in projective space, Amer.\ J.\ Math.\ 79 (1957) 403--430 [細部は\tQ].
\bibitem{Bre} G.\ Bredon, \emph{Sheaf Theory}, 2nd ed., Springer GTM 170, 1997 [細部は\tQ].
\bibitem{CF} M.\ Crainic, R.\ L.\ Fernandes, Integrability of Lie brackets, Ann.\ Math.\ 157 (2003) 575--620, arXiv:math/0105033.
\bibitem{CLS} D.\ Cox, J.\ Little, H.\ Schenck, \emph{Toric Varieties}, Graduate Studies in Math.\ 124, AMS 2011 [細部は\tQ].
\bibitem{CMNM} F.\ J.\ Calder\'on-Moreno, L.\ Narv\'aez-Macarro, Logarithmic comparison theorems for free divisors (Dualit\'e et comparaison sur les complexes de de Rham logarithmiques par rapport aux diviseurs libres), Proc.\ London Math.\ Soc.\ 98 (2009), doi:10.1112/plms/pdn043, arXiv:math/0603003 [細部は\tQ].
\bibitem{Cav} G.\ Cavalcanti, New aspects of the $dd^c$-lemma, arXiv:math/0501406; The decomposition of forms and cohomology of generalized complex manifolds, J.\ Geom.\ Phys.\ 57 (2006) 121--132.
\bibitem{Del} P.\ Deligne, \'Equations diff\'erentielles \`a points singuliers r\'eguliers, Lecture Notes in Math.\ 163, Springer 1970.
\bibitem{Del68} P.\ Deligne, Th\'eor\`eme de Lefschetz et crit\`eres de d\'eg\'en\'erescence de suites spectrales, Publ.\ Math.\ IH\'ES 35 (1968) 259--278.
\bibitem{DelFG} P.\ Deligne, Le groupe fondamental du compl\'ement d'une courbe plane n'ayant que des points doubles ordinaires est ab\'elien, S\'em.\ Bourbaki 543 (1979/80).
\bibitem{EV} H.\ Esnault, E.\ Viehweg, Logarithmic de Rham complexes and vanishing theorems, Invent.\ Math.\ 86 (1986) 161--194.
\bibitem{Ful} W.\ Fulton, On the fundamental group of the complement of a node curve, Ann.\ Math.\ 111 (1980) 407--409.
\bibitem{GH} P.\ Griffiths, J.\ Harris, Principles of Algebraic Geometry, Wiley 1978.
\bibitem{GL} M.\ Gualtieri, S.\ Li, Symplectic groupoids of log symplectic manifolds, IMRN 2014, arXiv:1206.3674.
\bibitem{Goto} R.\ Goto, Unobstructed deformations of generalized complex structures induced by $C^\infty$ logarithmic symplectic structures, arXiv:1501.03398 [掲載誌・巻頁は未確認\tQ].
\bibitem{Gua11} M.\ Gualtieri, Generalized complex geometry, Ann.\ Math.\ 174 (2011) 75--123, arXiv:math/0703298.
\bibitem{GuaBr} M.\ Gualtieri, Branes on Poisson varieties, arXiv:0710.2719 (v2, 2010-07；arXiv ページに掲載書誌の記載なし。The Many Facets of Geometry, Oxford Univ.\ Press 2010 は記憶\tQ).
\bibitem{GuaHodge} M.\ Gualtieri, Generalized geometry and the Hodge decomposition, arXiv:math/0409093.
\bibitem{GuaK} M.\ Gualtieri, Generalized K\"ahler geometry, Comm.\ Math.\ Phys.\ 331 (2014) 297--331, arXiv:1007.3485.
\bibitem{HHM} T.\ Hausel, E.\ Hunsicker, R.\ Mazzeo, Hodge cohomology of gravitational instantons, Duke Math.\ J.\ 122 (2004) 485--548, arXiv:math/0207169.
\bibitem{Hart} R.\ Hartshorne, Algebraic Geometry, GTM 52, Springer 1977.
\bibitem{Hit03} N.\ Hitchin, Generalized Calabi--Yau manifolds, Q.\ J.\ Math.\ 54 (2003) 281--308, arXiv:math/0209099.
\bibitem{Hu1} T.\ Hu, Kodaira vanishing theorems for K\"ahler Lie algebroids, arXiv:2403.10876.
\bibitem{Hu2} T.\ Hu, Lie groupoid Riemann--Roch--Hirzebruch theorem and applications, arXiv:2403.13424.
\bibitem{Huy} D.\ Huybrechts, Complex Geometry: An Introduction, Springer 2005.
\bibitem{Kod53} K.\ Kodaira, On a differential-geometric method in the theory of analytic stacks, Proc.\ Nat.\ Acad.\ Sci.\ USA 39 (1953) 1268--1273.
\bibitem{Kod54} K.\ Kodaira, On K\"ahler varieties of restricted type, Ann.\ Math.\ 60 (1954) 28--48.
\bibitem{LGSX1} C.\ Laurent-Gengoux, M.\ Sti\'enon, P.\ Xu, Holomorphic Poisson manifolds and holomorphic Lie algebroids, IMRN 2008, arXiv:0707.4253.
\bibitem{LGSX2} C.\ Laurent-Gengoux, M.\ Sti\'enon, P.\ Xu, Integration of holomorphic Lie algebroids, arXiv:0803.2031.
\bibitem{MM} R.\ Mazzeo, R.\ Melrose, Pseudodifferential operators on manifolds with fibred boundaries, arXiv:math/9812120.
\bibitem{Milnor} J.\ Milnor, \emph{Singular Points of Complex Hypersurfaces}, Annals of Math.\ Studies 61, Princeton Univ.\ Press 1968 [細部は\tQ].
\bibitem{Mok} T.\ Mokri, Matched pairs of Lie algebroids, Glasgow Math.\ J.\ 39 (1997) 167--181, doi:10.1017/S0017089500032055 [年・DOI は出版社ページで確認、巻頁は記憶\tQ].
\bibitem{Mum} D.\ Mumford, Varieties defined by quadratic equations, in: Questions on Algebraic Varieties (C.I.M.E., Varenna 1969), Cremonese 1970 [細部は\tQ].
\bibitem{PP} D.\ Pal, M.\ Poddar, Strong generalized holomorphic principal bundles, arXiv:2404.18113; Ann.\ Sc.\ Norm.\ Super.\ Pisa Cl.\ Sci., doi:10.2422/2036-2145.202501\_026（掲載巻号は\tQ）.
\bibitem{PPT} M.\ Pflaum, H.\ Posthuma, X.\ Tang, The index of geometric operators on Lie groupoids, Indag.\ Math.\ 25 (2014) 1--18, arXiv:1308.0236.
\bibitem{Pym} B.\ Pym, Constructions and classifications of projective Poisson varieties, Lett.\ Math.\ Phys.\ 108 (2018), no.\ 3, 573--632, doi:10.1007/s11005-017-0984-5.
\bibitem{QS} Y.\ Qiao, Y.\ So, Deformation and K-theoretic index formulae on boundary groupoids, arXiv:2304.13345.
\bibitem{SGA4} M.\ Artin, Th\'eor\`eme de finitude pour un morphisme propre; dimension cohomologique des sch\'emas alg\'ebriques affines, in SGA 4, Exp.\ XIV, Lecture Notes in Math.\ 305 [細部は\tQ].
\bibitem{Saito} K.\ Saito, Theory of logarithmic differential forms and logarithmic vector fields, J.\ Fac.\ Sci.\ Univ.\ Tokyo Sect.\ IA Math.\ 27 (1980) 265--291 [判定法の定理番号は記憶\tQ].
\bibitem{Serre} J.-P.\ Serre, Faisceaux alg\'ebriques coh\'erents, Ann.\ Math.\ 61 (1955) 197--278.
\bibitem{Tim} K.\ Timmerscheidt, Mixed Hodge theory for unitary local systems, J.\ reine angew.\ Math.\ 379 (1987) 152--171.
\end{thebibliography}

\end{document}
